% chapters/solutions.tex
% 部分习题解答 (Solutions of Some Exercises)
% 源: 原书 PDF 第 375--437 页 (印刷页 371--434)

\chapter*{部分习题解答}
\addcontentsline{toc}{chapter}{部分习题解答}
\markboth{部分习题解答}{部分习题解答}

在本节中, 为了避免与前面各节的公式混淆, 公式用 (S1), (S2) 等编号.

\solutionhead{1.1}

1. 等式 $\dual{f}{x}=\norm{x}^2$ 蕴含 $\norm{x}\le\norm{f}$.推论  \ref{cor:1.3} 蕴含
$F(x)$ 非空.由 $F(x)$ 的第二种形式显然可见 $F(x)$ 是闭凸集.

2. 在严格凸赋范空间中, 包含在球面内的任何非空凸集都退化为单点.

3. 注意
\begin{equation*}
\dual{f}{y}\le\norm{f}\norm{y}\le\frac12\norm{f}^2+\frac12\norm{y}^2.
\end{equation*}
反之, 设 $f$ 满足
\begin{equation}
\frac12\norm{y}^2-\frac12\norm{x}^2\ge\dual{f}{y-x}\quad\forall y\in E.
\tag{S1}
\label{eq:S1.1}
\end{equation}
先在 \eqref{eq:S1.1} 中取 $y=\lambda x$ ($\lambda\in\R$); 变动 $\lambda$
可见 $\dual{f}{x}=\norm{x}^2$.接着在 \eqref{eq:S1.1} 中选取满足
$\norm{y}=\delta>0$ 的 $y$; 由此推出
\begin{equation*}
\dual{f}{y}\le\frac12\delta^2+\frac12\norm{x}^2.
\end{equation*}
因此我们得到
\begin{equation*}
\delta\norm{f}=\sup_{\substack{y\in E\\\norm{y}=\delta}}\dual{f}{y}
\le\frac12\delta^2+\frac12\norm{x}^2.
\end{equation*}
选取 $\delta=\norm{x}$ 即得出结论.

4. 若 $f\in F(x)$, 则有
\begin{equation*}
\frac12\norm{y}^2-\frac12\norm{x}^2\ge\dual{f}{y-x},
\end{equation*}
若 $g\in F(y)$, 则有
\begin{equation*}
\frac12\norm{x}^2-\frac12\norm{y}^2\ge\dual{g}{x-y}.
\end{equation*}
把这两个不等式相加得出 $\dual{f-g}{x-y}\ge0$.另一方面, 注意
\begin{equation*}
\dual{f-g}{x-y}=\norm{x}^2+\norm{y}^2-\dual{f}{y}-\dual{g}{x}
\ge\norm{x}^2+\norm{y}^2-2\norm{x}\norm{y}.
\end{equation*}

5. 由第 4 问我们已知 $\norm{x}=\norm{y}$.另一方面, 我们有
\begin{equation*}
\dual{F(x)-F(y)}{x-y}=[\norm{x}^2-\dual{F(x)}{y}]+[\norm{y}^2-\dual{F(y)}{x}],
\end{equation*}
括号内的两项都 $\ge0$.由此推出
$\norm{x}^2=\norm{y}^2=\dual{F(x)}{y}=\dual{F(y)}{x}$, 这蕴含
$F(x)\in F(y)$, 从而由第 2 问 $F(x)=F(y)$.

\solutionhead{1.2}

1(a). $\norm{f}_{E^\star}=\max_{1\le i\le n}\abs{f_i}$.

1(b). $f\in F(x)$ 当且仅当对每个 $1\le i\le n$, 有
\begin{equation*}
f_i=\begin{cases}
(\operatorname{sign}x_i)\norm{x}_1 & \text{若}\ x_i\ne0,\\
\text{区间}\ [-\norm{x}_1,+\norm{x}_1]\ \text{中的任意数} & \text{若}\ x_i=0.
\end{cases}
\end{equation*}

2(a). $\norm{f}_{E^\star}=\sum_{i=1}^n\abs{f_i}$.

2(b). 给定 $x\in E$, 考虑集合
\begin{equation*}
I=\{1\le i\le n;\ \abs{x_i}=\norm{x}_\infty\}.
\end{equation*}
则 $f\in F(x)$ 当且仅当成立
\begin{enumerate}[label=(\roman*), itemsep=0.1em]
	\item 对一切 $i\notin I$, $f_i=0$,
	\item 对一切 $i\in I$, $f_ix_i\ge0$, 且 $\sum_{i\in I}\abs{f_i}=\norm{x}_\infty$.
\end{enumerate}

3. $\norm{f}_{E^\star}=\left(\sum_{i=1}^n\abs{f_i}^2\right)^{1/2}$, 且
$f\in F(x)$ 当且仅当对一切 $i=1,2,\dots,n$, $f_i=x_i$.更一般地,
\begin{equation*}
\norm{f}_{E^\star}=\left(\sum_{i=1}^n\abs{f_i}^{p'}\right)^{1/p'},
\end{equation*}
其中 $1/p+1/p'=1$, 且 $f\in F(x)$ 当且仅当对一切 $i=1,2,\dots,n$,
$f_i=\abs{x_i}^{p-2}x_i/\norm{x}_p^{p-2}$.

\solutionhead{1.3}

1. $\norm{f}_{E^\star}=1$ (注意对一切 $\alpha>0$, $f(t^\alpha)=1/(1+\alpha)$).

2. 若存在这样的 $u$, 则会有 $\int_0^1(1-u)dt=0$, 从而 $u\equiv1$; 矛盾.

\solutionhead{1.5}

1. 用 $\mathcal{P}$ 记 $E$ 的所有线性无关子集的族.容易看出 $\mathcal{P}$
(按通常的包含关系偏序) 是归纳的.Zorn 引理蕴含 $\mathcal{P}$ 有极大元,
记作 $(e_i)_{i\in I}$, 它显然是一个代数基.由于对一切 $i\in I$, $e_i\ne0$,
经单位化后可以假设对一切 $i\in I$, $\norm{e_i}=1$.

2. 由于 $E$ 无限维, 可以假设 $\N\subset I$.存在 $E$ 上的 (唯一) 线性泛函,
使得当 $i\in\N$ 时 $f(e_i)=i$, 而当 $i\in I\setminus\N$ 时 $f(e_i)=0$.

3. 假设 $I$ 是可数的, 即 $I=\N$.考虑由 $(e_i)_{1\le i\le n}$ 张成的向量
空间 $F_n$.$F_n$ 是闭的 (见 \S \ref{sec:11.1}), 并且
$\bigcup_{n=1}^\infty F_n=E$.由 Baire 纲定理推出存在某个 $n_0$ 使得
$\operatorname{Int}(F_{n_0})\ne\emptyset$.因此 $E=F_{n_0}$; 矛盾.

\solutionhead{1.7}

1. 设 $x,y\in\overline C$, 从而 $x=\lim x_n$, $y=\lim y_n$, 其中
$x_n,y_n\in C$.于是 $tx+(1-t)y=\lim[tx_n+(1-t)y_n]$, 因此对一切
$t\in[0,1]$, $tx+(1-t)y\in\overline C$.

设 $x,y\in\operatorname{Int}C$, 从而存在某个 $r>0$ 使得 $B(x,r)\subset C$
且 $B(y,r)\subset C$.由此推出对一切 $t\in[0,1]$,
\begin{equation*}
tB(x,r)+(1-t)B(y,r)\subset C.
\end{equation*}
但 $tB(x,r)+(1-t)B(y,r)=B(tx+(1-t)y,r)$ (为什么?).

2. 设 $r>0$ 使得 $B(y,r)\subset C$.对一切 $t\in[0,1]$, 有
\begin{equation*}
tx+(1-t)B(y,r)\subset C,
\end{equation*}
从而 $B(tx+(1-t)y,(1-t)r)\subset C$.由此推出对一切 $t\in[0,1)$,
$tx+(1-t)y\in\operatorname{Int}C$.

3. 固定任何 $y_0\in\operatorname{Int}C$.给定 $x\in\overline C$, 有
$x=\lim_{n\to\infty}\left[\left(1-\frac1n\right)x+\frac1ny_0\right]$.但
$\left(1-\frac1n\right)x+\frac1ny_0\in\operatorname{Int}C$, 因此
$x\in\overline{\operatorname{Int}C}$.这证明了
$\overline C\subset\overline{\operatorname{Int}C}$, 从而
$\overline C=\overline{\operatorname{Int}C}$.

\solutionhead{1.8}

1. 我们已知对一切 $\lambda>0$ 与一切 $x\in E$, $p(\lambda x)=\lambda p(x)$,
且对一切 $x,y\in E$, $p(x+y)\le p(x)+p(y)$.还须验证:

(i) 对一切 $x\in E$, $p(-x)=p(x)$, 这由 $C$ 的对称性推出.

(ii) $p(x)=0\Rightarrow x=0$, 这由 $C$ 有界的事实推出.更精确地说, 设
$L>0$ 使得对一切 $x\in C$, $\norm{x}\le L$.容易看出对一切 $x\in E$,
$p(x)\ge\frac1L\norm{x}$.

2. $C$ 不是有界的.例如考虑序列 $u_n(t)=\sqrt{n}/(1+nt)$, 验证 $u_n\in C$,
而 $\norm{u_n}=\sqrt{n}$.这里 $p(u)=\left(\int_0^1\abs{u(t)}^2dt\right)^{1/2}$
是一个与 $\norm{u}$ 不等价的范数.

\solutionhead{1.9}

1. 令
\begin{equation*}
P=\left\{\lambda=(\lambda_1,\lambda_2,\dots,\lambda_n)\in\R^n;\
\lambda_i\ge0\ \forall i\ \text{且}\ \sum_{i=1}^n\lambda_i=1\right\},
\end{equation*}
于是 $P$ 是 $\R^n$ 的紧子集, 且 $C_n$ 是 $P$ 在连续映射
$\lambda\mapsto\sum_{i=1}^n\lambda_ix_i$ 下的像.

2. 对 $C_n$ 与 $\{0\}$ 应用 Hahn--Banach 第二几何形式.把与分离 $C_n$ 和
$\{0\}$ 的超平面相关联的线性泛函单位化.

4. 对 $C=A-B$ 应用上述构造.

\solutionhead{1.10}

(A) $\Rightarrow$ (B) 是显然的.

(B) $\Rightarrow$ (A). 设 $G$ 为由 $x_i$ ($i\in I$) 张成的向量空间.给定
$x\in G$, 把 $x$ 写成 $x=\sum_{i\in J}\beta_ix_i$, 并令
$g(x)=\sum_{i\in J}\beta_i\alpha_i$.假设 (B) 蕴含这个定义是有意义的, 且
对一切 $x\in G$, $\abs{g(x)}\le M\norm{x}$.接着用推论  \ref{cor:1.2} 把
$g$ 延拓到整个 $E$.

\solutionhead{1.11}

(A) $\Rightarrow$ (B) 再次是显然的.

(B) $\Rightarrow$ (A). 先设 $f_i$ 线性无关 ($1\le i\le n$).令
$\alpha=(\alpha_1,\alpha_2,\dots,\alpha_n)\in\R^n$.考虑由
\begin{equation*}
\phi(x)=(\dual{f_1}{x},\dots,\dual{f_n}{x})
\end{equation*}
定义的映射 $\phi:E\to\R^n$.令 $C=\{x\in E;\ \norm{x}\le M+\varepsilon\}$.
须证明 $\alpha\in\phi(C)$.反设 $\alpha\notin\phi(C)$, 并分离 $\phi(C)$ 与
$\{\alpha\}$ (见习题 1.9).因此存在某个 $\beta=(\beta_1,\dots,\beta_n)\in\R^n$,
$\beta\ne0$, 使得对一切 $x\in C$, $\beta\cdot\phi(x)\le\beta\cdot\alpha$, 即
\begin{equation*}
\dual{\sum\beta_if_i}{x}\le\sum\beta_i\alpha_i\quad\forall x\in C.
\end{equation*}
由此推出 $(M+\varepsilon)\norm{\sum\beta_if_i}\le\sum\beta_i\alpha_i$.利用
假设 (B), 发现 $\sum\beta_if_i=0$.由于 $f_i$ 线性无关, 得出 $\beta=0$;
矛盾.

在一般情形下, 把上述结果应用于 $(f_i)_{1\le i\le n}$ 的极大线性无关子集.

\solutionhead{1.15}

1. 显然 $C\subset C^{\star\star}$ 且 $C^{\star\star}$ 是闭的.反之, 设
$x_0\in C^{\star\star}$ 且 $x_0\notin C$.可以严格分离 $\{x_0\}$ 与 $C$,
从而存在某个 $f_0\in E^\star$ 与某个 $\alpha_0\in\R$ 使得
\begin{equation*}
\dual{f_0}{x}<\alpha_0<\dual{f_0}{x_0}\quad\forall x\in C.
\end{equation*}
由于 $0\in C$, 推出 $\alpha_0>0$; 令 $f=(1/\alpha_0)f_0$, 有
\begin{equation*}
\dual{f}{x}<1<\dual{f}{x_0}\quad\forall x\in C.
\end{equation*}
因此 $f\in C^\star$, 我们导出了矛盾, 因为 $x_0\in C^{\star\star}$.

2. 若 $C$ 是线性子空间, 则
\begin{equation*}
C^\star=\{f\in E^\star;\ \dual{f}{x}=0\ \forall x\in C\}=C^\perp.
\end{equation*}

\solutionhead{1.18}

(a) $\varphi^\star(f)=\begin{cases}-b & \text{若}\ f=a,\\ +\infty & \text{若}\ f\ne a.\end{cases}$

(b) $\varphi^\star(f)=\begin{cases}f\log f-f & \text{若}\ f>0,\\ 0 & \text{若}\ f=0,\\ +\infty & \text{若}\ f<0.\end{cases}$

(c) $\varphi^\star(f)=\abs{f}$.

(d) $\varphi^\star(f)=0$.

(e) $\varphi^\star(f)=\begin{cases}+\infty & \text{若}\ f\ge0,\\ -1-\log\abs{f} & \text{若}\ f<0.\end{cases}$

(f) $\varphi^\star(f)=(1+f^2)^{1/2}$.

(g) $\varphi^\star(f)=\begin{cases}\frac12f^2 & \text{若}\ \abs{f}\le1,\\ +\infty & \text{若}\ \abs{f}>1.\end{cases}$

(h) $\varphi^\star(f)=\frac{1}{p'}\abs{f}^{p'}$, 其中 $\dfrac1p+\dfrac{1}{p'}=1$.

(i) $\varphi^\star(f)=\begin{cases}0 & \text{若}\ 0\le f\le1,\\ +\infty & \text{其他}.\end{cases}$

(j) $\varphi^\star(f)=\begin{cases}\frac{1}{p'}f^{p'} & \text{若}\ f\ge0,\\ 0 & \text{若}\ f<0.\end{cases}$

(k) $\varphi^\star(f)=\begin{cases}+\infty & \text{若}\ f\ge0,\\ -\frac{1}{p'}\abs{f}^{p'} & \text{若}\ f<0.\end{cases}$

(l) $\varphi^\star(f)=\abs{f}+\frac{1}{p'}\abs{f}^{p'}$.

\solutionhead{1.20}

共轭函数定义在 $\ell^{p'}$ 上 ($\dfrac1p+\dfrac{1}{p'}=1$), 由下式给出:

(a) $\varphi^\star(f)=\begin{cases}\dfrac14\sum_{k=1}^\infty\dfrac1k\abs{f_k}^2 & \text{若}\ \sum_{k=1}^\infty\dfrac1k\abs{f_k}^2<+\infty,\\ +\infty & \text{其他}.\end{cases}$

(b) $\varphi^\star(f)=\begin{cases}\sum_{k=2}^{+\infty}a_k\abs{f_k}^{k/(k-1)} & \text{若}\ \sum_{k=2}^{+\infty}a_k\abs{f_k}^{k/(k-1)}<+\infty,\\ +\infty & \text{其他},\end{cases}$
其中 $a_k=(k-1)k^{-k/(k-1)}$.

(c) $\varphi^\star(f)=\begin{cases}0 & \text{若}\ \norm{f}_{\ell^\infty}\le1,\\ +\infty & \text{其他}.\end{cases}$

\solutionhead{1.21}

2. $\varphi^\star=I_A$, 其中 $A=\{[f_1,f_2];\ f_1\le0,\ f_2\le0,\ \text{且}\ 4f_1f_2\ge1\}$.

3. 有
\begin{equation*}
\inf_{x\in E}\{\varphi(x)+\psi(x)\}=0
\end{equation*}
且 $\varphi^\star=I_{D^\perp}$, 其中 $D^\perp=\{[f_1,f_2];\ f_2=0\}$.
由此推出对一切 $f\in E^\star$, $\varphi^\star(-f)+\psi^\star(f)=+\infty$, 从而
\begin{equation*}
\sup_{f\in E^\star}\{-\varphi^\star(-f)-\psi^\star(f)\}=-\infty.
\end{equation*}

4. 定理  \ref{thm:1.12} 的假设不满足: 不存在元素 $x_0\in E$ 使得
$\varphi(x_0)<+\infty$, $\psi(x_0)<+\infty$, 且 $\varphi$ 在 $x_0$ 处连续.

\solutionhead{1.22}

1. 写出
\begin{equation*}
\norm{x-a}\le\norm{x-y}+\norm{y-a}.
\end{equation*}
取 $\inf_{a\in A}$ 给出 $\varphi(x)\le\norm{x-y}+\varphi(y)$.然后交换 $x$
与 $y$.

2. 固定 $x,y\in E$ 与 $t\in[0,1]$.给定 $\varepsilon>0$, 存在 $a\in A$ 与
$b\in A$ 使得
\begin{equation*}
\norm{x-a}\le\varphi(x)+\varepsilon\quad\text{且}\quad
\norm{y-b}\le\varphi(y)+\varepsilon.
\end{equation*}
因此
\begin{equation*}
\norm{tx+(1-t)y-[ta+(1-t)b]}\le t\varphi(x)+(1-t)\varphi(y)+\varepsilon.
\end{equation*}
但 $ta+(1-t)b\in A$, 从而对一切 $\varepsilon>0$,
\begin{equation*}
\varphi(tx+(1-t)y)\le t\varphi(x)+(1-t)\varphi(y)+\varepsilon.
\end{equation*}

3. 由于 $A$ 闭, 有 $A=\{x\in E;\ \varphi(x)\le0\}$, 因此若 $\varphi$ 是
凸的, 则 $A$ 是凸的.

4. 有
\begin{align*}
\varphi^\star(f)&=\sup_{x\in E}\{\dual{f}{x}-\inf_{a\in A}\norm{x-a}\}
=\sup_{x\in E}\sup_{a\in A}\{\dual{f}{x}-\norm{x-a}\}\\
&=\sup_{a\in A}\sup_{x\in E}\{\dual{f}{x}-\norm{x-a}\}
=(I_A)^\star(f)+I_{B_{E^\star}}(f).
\end{align*}

\solutionhead{1.23}

1. 设 $f\in D(\varphi^\star)\cap D(\psi^\star)$.对每个 $x,y\in E$, 有
\begin{align*}
\dual{f}{x-y}-\varphi(x-y)&\le\varphi^\star(f),\\
\dual{f}{y}-\psi(y)&\le\psi^\star(f).
\end{align*}
把这两个不等式相加得出
\begin{equation*}
(\varphi\mathbin{\triangledown}\psi)(x)\ge\dual{f}{x}-\varphi^\star(f)-\psi^\star(f).
\end{equation*}
特别地, $(\varphi\mathbin{\triangledown}\psi)(x)>-\infty$.此外, 我们有
\begin{align*}
(\varphi\mathbin{\triangledown}\psi)^\star(f)
&=\sup_{x\in E}\{\dual{f}{x}-\inf_{y\in E}[\varphi(x-y)+\psi(y)]\}\\
&=\sup_{x\in E}\sup_{y\in E}\{\dual{f}{x}-\varphi(x-y)-\psi(y)\}\\
&=\sup_{y\in E}\sup_{x\in E}\{\dual{f}{x}-\varphi(x-y)-\psi(y)\}\\
&=\varphi^\star(f)+\psi^\star(f).
\end{align*}

2. 须验证对一切 $f,g\in E^\star$ 及一切 $x\in E$,
\begin{equation*}
\dual{f}{x}-\varphi(x)-\psi(x)\le\varphi^\star(f-g)+\psi^\star(g).
\end{equation*}
写成 $\dual{f}{x}=\dual{f-g}{x}+\dual{g}{x}$ 即显然.

3. 给定 $f\in E^\star$, 须证明
\begin{equation}
\sup_{x\in E}\{\dual{f}{x}-\varphi(x)-\psi(x)\}
=\inf_{g\in E^\star}\{\varphi^\star(f-g)+\psi^\star(g)\}.
\tag{S1}
\label{eq:S1.23}
\end{equation}
注意
\begin{equation*}
\sup_{x\in E}\{\dual{f}{x}-\varphi(x)-\psi(x)\}
=-\inf_{x\in E}\{\tilde\varphi(x)+\psi(x)\},
\end{equation*}
其中 $\tilde\varphi(x)=\varphi(x)-\dual{f}{x}$.把定理  \ref{thm:1.12}
应用于函数 $\tilde\varphi$ 与 $\psi$ 给出
\begin{equation*}
\inf_{x\in E}\{\tilde\varphi(x)+\psi(x)\}
=\sup_{g\in E^\star}\{-\tilde\varphi^\star(-g)-\psi^\star(g)\},
\end{equation*}
它恰好对应于 \eqref{eq:S1.23}.

4. 显然有
\begin{align*}
(\varphi^\star\mathbin{\triangledown}\psi^\star)^\star(x)
&=\sup_{f\in E^\star}\{\dual{f}{x}-\inf_{g\in E^\star}[\varphi^\star(f-g)+\psi^\star(g)]\}\\
&=\sup_{f\in E^\star}\sup_{g\in E^\star}\{\dual{f}{x}-\varphi^\star(f-g)-\psi^\star(g)\}\\
&=\sup_{g\in E^\star}\sup_{f\in E^\star}\{\dual{f}{x}-\varphi^\star(f-g)-\psi^\star(g)\}\\
&=\varphi^{\star\star}(x)+\psi^{\star\star}(x).
\end{align*}

\solutionhead{1.24}

1. 我们知道 (命题  \ref{prop:1.10}) 存在某个 $f\in E^\star$ 与常数 $C$ 使得
对一切 $y\in E$, $\varphi(y)\ge\dual{f}{y}-C$.选取 $n\ge\norm{f}$, 则有
$\varphi_n(x)\ge-n\norm{x}-C>-\infty$.

2. 函数 $\varphi_n$ 是两个凸函数的下卷积; 因此 $\varphi_n$ 是凸的 (见习题
1.23 第 7 问).为证明 $\abs{\varphi_n(x_1)-\varphi_n(x_2)}\le n\norm{x_1-x_2}$,
使用与习题 1.22 第 1 问相同的论证.

3. $(\varphi_n)^\star=I_{nB_{E^\star}}+\varphi^\star$ (由习题 1.23 第 1 问).

5. 由第 1 问我们对一切 $y\in E$ 有 $\varphi(y)\ge-\norm{f}\norm{y}-C$, 这给出
\begin{equation*}
n\norm{x-y_n}\le\norm{f}\norm{y_n}+C+\varphi(x)+1/n.
\end{equation*}
由此推出当 $n\to\infty$ 时 $\norm{y_n}$ 保持有界, 因此
$\lim_{n\to\infty}\norm{x-y_n}=0$.另一方面, 我们有
$\varphi_n(x)\ge\varphi(y_n)-1/n$, 由于 $\varphi$ 下半连续, 我们得出
$\liminf_{n\to\infty}\varphi_n(x)\ge\varphi(x)$.

6. 反设存在常数 $C$ 使得沿一个仍记作 $\varphi_n(x)$ 的子序列有
$\varphi_n(x)\le C$.如第 5 问那样选取 $y_n$, 可见 $y_n\to x$.此外,
$\varphi(y_n)\le C+1/n$, 从而
$\varphi(x)\le\liminf_{n\to\infty}\varphi(y_n)\le C$; 矛盾.

\solutionhead{1.25}

4. 对每个固定的 $t>0$, 函数
\begin{equation*}
y\mapsto\frac{1}{2t}\left(\norm{x+ty}^2-\norm{x}^2\right)
\end{equation*}
是凸的.因此作为凸函数的极限, 函数 $y\mapsto[x,y]$ 是凸的.另一方面,
\begin{equation*}
G(x,y)=\sup_{t>0}\left\{-\frac{1}{2t}\left(\norm{x+ty}^2-\norm{x}^2\right)\right\}
\end{equation*}
作为连续函数的上确界是下半连续的.

5. 我们已知 (见习题 1.1 第 3 问)
\begin{equation*}
\frac12\norm{x+ty}^2-\frac12\norm{x}^2\ge\dual{f}{ty},
\end{equation*}
因此对一切 $x,y\in E$ 及一切 $f\in F(x)$,
\begin{equation*}
[x,y]\ge\dual{f}{y}.
\end{equation*}
另一方面, 有
\begin{equation*}
\varphi^\star(f)=\frac12\norm{f}^2-\dual{f}{x}+\frac12\norm{x}^2,
\end{equation*}
以及
\begin{equation*}
\psi^\star(f)=\begin{cases}
0 & \text{若}\ \dual{f}{y}+\alpha\le0,\\
+\infty & \text{若}\ \dual{f}{y}+\alpha>0.
\end{cases}
\end{equation*}
容易验证 $\inf_{z\in E}\{\varphi(z)+\psi(z)\}=0$.由定理  \ref{thm:1.12}
推出存在某个 $f_0\in E^\star$ 使得 $\varphi^\star(f_0)+\psi^\star(-f_0)=0$,
即 $\dual{f_0}{y}\ge\alpha$ 且
$\frac12\norm{f_0}^2-\dual{f_0}{x}+\frac12\norm{x}^2=0$.从而我们有
$\norm{f_0}=\norm{x}$ 且 $\dual{f_0}{x}=\norm{x}^2$, 即 $f_0\in F(x)$.

6. (a) $1<p<\infty$, $[x,y]=\dfrac{\sum\abs{x_i}^{p-2}x_iy_i}{\norm{x}_p^{p-2}}$.

(b) $p=1$, $[x,y]=\norm{x}_1\left(\sum_{x_i\ne0}(\operatorname{sign}x_i)y_i
+\sum_{x_i=0}\abs{y_i}\right)$.

(c) $p=\infty$, $[x,y]=\max_{i\in I}\{x_iy_i\}$, 其中
$I=\{1\le i\le n;\ \abs{x_i}=\norm{x}_\infty\}$.

\solutionhead{1.27}

设 $\widetilde T:E\to F$ 为 $T$ 的连续线性延拓.容易验证
$E=N(\widetilde T)+G$ 且 $N(\widetilde T)\cap G=\{0\}$, 从而
$N(\widetilde T)$ 是 $G$ 的补空间; 矛盾.
\solutionhead{2.1}

不失一般性可以假设 $x_0=0$.

1. 设 $X=\{x\in E;\ \norm{x}\le\rho\}$, 其中 $\rho>0$ 足够小使得 $X\subset D(\varphi)$.
集合 $F_n$ 是闭的且 $\bigcup_{n=1}^\infty F_n=X$.由 Baire 纲定理存在某个 $n_0$
使得 $\operatorname{Int}(F_{n_0})\ne\emptyset$.设 $x_1\in E$ 与 $\rho_1>0$
使得 $B(x_1,\rho_1)\subset F_{n_0}$.给定任何满足 $\norm{x}<\rho_1/2$ 的
$x\in E$, 写出 $x=\frac12(x_1+2x)+\frac12(-x_1)$, 得出
$\varphi(x)\le\frac12n_0+\frac12\varphi(-x_1)$.

2. 存在某个 $\xi\in E$ 与某个常数 $t\in[0,1]$ 使得 $\norm{\xi}=R$ 且
$x_2=tx_1+(1-t)\xi$.由此推出
\begin{equation*}
\varphi(x_2)\le t\varphi(x_1)+(1-t)M,
\end{equation*}
从而 $\varphi(x_2)-\varphi(x_1)\le(1-t)[M-\varphi(x_1)]$.但
$x_2-x_1=(1-t)(\xi-x_1)$, 从而 $\norm{x_2-x_1}\ge(1-t)(R-r)$.因此我们有
\begin{equation*}
\varphi(x_2)-\varphi(x_1)\le\frac{\norm{x_2-x_1}}{R-r}[M-\varphi(x_1)].
\end{equation*}
另一方面, 若 $x_2=0$, 得到 $t\norm{x_1}=(1-t)R$, 因此
$(1-t)=\frac{\norm{x_1}}{\norm{x_1}+R}\le\frac12$.由此推出
$\varphi(0)-\varphi(x_1)\le\frac12[M-\varphi(x_1)]$, 从而
$M-\varphi(x_1)\le2[M-\varphi(0)]$.

\solutionhead{2.2}

由 (i) 我们有 $p(0)\le p(x_n)+p(-x_n)\to0$, 从而 $p(0)\le0$.另一方面由 (i)
$p(0)\le2p(0)$.因此 $p(0)=0$.

接着证明 $p(\alpha_nx_n)\to0$.反设沿某个子序列对某个 $\varepsilon>0$
有 $\abs{p(\alpha_nx_n)}>2\varepsilon$.抽取进一步的子序列, 可以假设对某个
$\alpha\in\R$ 有 $\alpha_n\to\alpha$.为简便起见, 仍用 $(x_n)$ 与 $(\alpha_n)$
记相应的序列.

集合 $F_n$ 是闭的且 $\bigcup_{n\ge1}F_n=\R$.应用 Baire 纲定理, 我们找到某个
$n_0$ 使得 $\operatorname{Int}F_{n_0}\ne\emptyset$.因此存在某个 $\lambda_0\in\R$
与某个 $\delta>0$ 使得对一切 $k\ge n_0$ 及一切满足 $\abs{t}<\delta$ 的 $t$,
有 $\abs{p((\lambda_0+t)x_k)}\le\varepsilon$.另一方面, 注意
\begin{align*}
p(\alpha_kx_k)&\le p((\lambda_0+\alpha_k-\alpha)x_k)+p((\alpha-\lambda_0)x_k),\\
-p(\alpha_kx_k)&\le-p((\lambda_0+\alpha_k-\alpha)x_k)+p((\lambda_0-\alpha)x_k).
\end{align*}
因此对足够大的 $k$, 我们得到 $\abs{p(\alpha_kx_k)}\le2\varepsilon$.矛盾.

最后, 写出
\begin{equation*}
p(\alpha_nx_n)-p(\alpha x)\le p(\alpha_n(x_n-x))+p(\alpha_nx)-p(\alpha x)\to0
\end{equation*}
以及
\begin{equation*}
p(\alpha_nx)\le p(\alpha_n(x-x_n))+p(\alpha_nx_n),
\end{equation*}
从而
\begin{equation*}
p(\alpha_nx_n)-p(\alpha x)\ge-p(\alpha_n(x_n-x))+p(\alpha_nx)-p(\alpha x)\to0.
\end{equation*}

\solutionhead{2.4}

由 (i), 存在线性算子 $T:E\to F^\star$ 使得对一切 $x,y$,
$a(x,y)=\dual{Tx}{y}_{F^\star,F}$.目标是证明 $T$ 是有界算子, 即 $T(B_E)$
在 $F^\star$ 中有界.由推论  \ref{cor:2.5}, 只须固定 $y\in F$ 并验证
$\dual{T(B_E)}{y}$ 有界.这由 (ii) 推出.

\solutionhead{2.6}

1. 有 $\dual{Ax_n-A(x_0+x)}{x_n-x_0-x}\ge0$, 从而
$\dual{Ax_n}{x}\le\varepsilon_n\norm{Ax_n}+C(x)$, 其中
$\varepsilon_n=\norm{x_n-x_0}$, $C(x)=\norm{A(x_0+x)}(1+\norm{x})$ (假设对
一切 $n$, $\varepsilon_n\le1$).由习题 2.5 推出 $(Ax_n)$ 有界; 矛盾.

2. 假设存在 $D(A)$ 中的序列 $(x_n)$ 使得 $x_n\to x_0$ 且
$\norm{Ax_n}\to\infty$.选取 $r>0$ 使得 $B(x_0,r)\subset\operatorname{conv}D(A)$.
对每个满足 $\norm{x}<r$ 的 $x\in E$, 写出
\begin{equation*}
x_0+x=\sum_{i=1}^m t_iy_i,\quad t_i\ge0\ \forall i,\quad
\sum_{i=1}^m t_i=1,\quad y_i\in D(A)\ \forall i
\end{equation*}
(当然 $t_i,y_i$ 与 $m$ 依赖于 $x$).我们有 $\dual{Ax_n-Ay_i}{x_n-y_i}\ge0$,
从而 $t_i\dual{Ax_n}{x_n-y_i}\ge t_i\dual{Ay_i}{x_n-y_i}$.由此推出
\begin{equation*}
\dual{Ax_n}{x_n-x_0-x}\ge\sum_{i=1}^m t_i\dual{Ay_i}{x_n-y_i},
\end{equation*}
这导出 $\dual{Ax_n}{x}\le\varepsilon_n\norm{Ax_n}+C(x)$, 其中
$\varepsilon_n=\norm{x_n-x_0}$,
$C(x)=\sum_{i=1}^m t_i\norm{Ay_i}(1+\norm{x_0-y_i})$.

3. 设 $x_0\in\operatorname{Int}D(A)$.沿用第 1 问中的相同论证, 可证明存在
两个常数 $R>0$ 与 $C$, 使得对一切满足 $\norm{x-x_0}<R$ 的 $x\in D(A)$ 及
一切 $f\in Ax$, 有 $\norm{f}\le C$.

\solutionhead{2.7}

对每个 $x\in\ell^p$ 令 $T_nx=\sum_{i=1}^n\alpha_ix_i$, 于是对每个
$x\in\ell^p$, $T_nx$ 收敛于一个极限.由推论  \ref{cor:2.3} 推出存在常数 $C$
使得对一切 $x\in\ell^p$ 及一切 $n$,
\begin{equation*}
\abs{T_nx}\le C\norm{x}_{\ell^p}.
\end{equation*}
适当地选取 $x$, 可见 $\alpha\in\ell^{p'}$ 且 $\norm{\alpha}_{p'}\le C$.

\solutionhead{2.8}

方法 (ii). 让我们验证 $T$ 的图是闭的.设 $(x_n)$ 为 $E$ 中的序列使得
$x_n\to x$ 且 $Tx_n\to f$.在不等式 $\dual{Tx_n-Ty}{x_n-y}\ge0$ 中取极限
得出
\begin{equation*}
\dual{f-Ty}{x-y}\ge0\quad\forall y\in E.
\end{equation*}
选取 $y=x+tz$ ($t\in\R,z\in E$), 可见 $f=Tx$.

\solutionhead{2.10}

1. 若 $T(M)$ 闭, 则 $M+N(T)=T^{-1}(T(M))$ 也闭.反之, 设 $M+N(T)$ 闭.
由于 $T$ 满, 有 $T((M+N(T))^c)=(T(M))^c$.开映射定理蕴含 $T((M+N(T))^c)$ 开,
从而 $T(M)$ 闭.

2. 若 $M$ 是任意闭子空间且 $N$ 是任意有限维空间, 则 $M+N$ 闭 (见 \S \ref{sec:11.1}).

\solutionhead{2.11}

由开映射定理, 存在常数 $c>0$ 使得 $T(B_E)\supset cB_F$.用 $(e_n)$ 记
$\ell^1$ 的典范基, 即 $e_n=(0,\dots,0,\underset{(n)}{1},0,\dots)$.存在某个
$u_n\in E$ 使得 $\norm{u_n}\le1/c$ 且 $T(u_n)=e_n$.给定
$y=(y_1,y_2,\dots,y_n,\dots)\in\ell^1$, 令 $Sy=\sum_{i=1}^\infty y_iu_i$.
显然级数收敛, 且 $S$ 具有所有要求的性质.

\solutionhead{2.12}

不失一般性可以假设 $T$ 满 (否则用 $R(T)$ 替代 $E$).反设存在 $E$ 中的序列
$(x_n)$ 使得 $\norm{x_n}_E=1$ 且 $\norm{Tx_n}_F+\abs{x_n}<1/n$.由开映射
定理存在常数 $c>0$ 使得 $T(B_E)\supset cB_F$.由于 $\norm{Tx_n}_F<1/n$,
存在某个 $y_n\in E$ 使得 $Tx_n=Ty_n$ 且 $\norm{y_n}_E<1/nc$.把 $x_n$ 写成
$x_n=y_n+z_n$, 其中 $z_n\in N(T)$, $\norm{y_n}_E\to0$ 且 $\norm{z_n}_E\to1$.
另一方面, $\abs{x_n}<1/n$; 因此
$\abs{z_n}<(1/n)+\abs{y_n}\le(1/n)+M\norm{y_n}_E$, 从而 $\abs{z_n}\to0$.
这是不可能的, 因为范数 $\norm{\cdot}_E$ 与 $\abs{\cdot}$ 在有限维空间 $N(T)$
上等价.

\solutionhead{2.13}

首先, 设 $T\in\mathcal{O}$, 从而 $T^{-1}\in\LL(F,E)$ (由推论  \ref{cor:2.7}).
则对满足 $\norm{U}$ 足够小的每个 $U\in\LL(E,F)$, $T+U\in\mathcal{O}$.
事实上, 方程 $Tx+Ux=f$ 可以写成 $x=T^{-1}(f-Ux)$; 只要
$\norm{T^{-1}}\norm{U}<1$, 它 (对每个 $f\in F$) 就有唯一解 (由 Banach
不动点定理; 见定理  \ref{thm:5.7}).

接着, 设 $T\in\mathcal{L}$.由定理  \ref{thm:2.13}, $R(T)$ 闭且在 $F$ 中有
补空间.设 $P:F\to R(T)$ 为连续投影.算子 $PT$ 是由 $E$ 到 $R(T)$ 上的
双射, 从而上面的分析适用.设 $U\in\LL(E,F)$ 满足 $\norm{U}<\delta$; 若
$\delta$ 足够小, 算子 $(PT+PU):E\to R(T)$ 是双射, 从而 $(PT+PU)^{-1}$
作为 $\LL(R(T),E)$ 的元素是良定的.令 $S=(PT+PU)^{-1}P$.显然
$S\in\LL(F,E)$ 且 $S(T+U)=I_E$.

\solutionhead{2.14}

1. 考虑商空间 $\widetilde E=E/N(T)$ 与典范满射 $\pi:E\to\widetilde E$, 使得
对一切 $x\in E$, $\norm{\pi x}_{\widetilde E}=\dist(x,N(T))$.$T$ 在
$\widetilde E$ 上诱导一个单射算子 $\widetilde T$.更精确地说, 把 $T$ 写成
$T=\widetilde T\circ\pi$, 其中 $\widetilde T\in\LL(\widetilde E,F)$, 从而
$R(T)=R(\widetilde T)$.另一方面, 推论  \ref{cor:2.7} 表明 $R(\widetilde T)$
闭当且仅当存在常数 $C$ 使得对一切 $y\in\widetilde E$,
$\norm{y}_{\widetilde E}\le C\norm{\widetilde Ty}$, 或等价地对一切 $x\in E$,
$\norm{\pi x}_{\widetilde E}\le C\norm{\widetilde T\pi x}$.最后一个不等式
即
\begin{equation*}
\dist(x,N(T))\le C\norm{Tx}\quad\forall x\in E.
\end{equation*}

\solutionhead{2.15}

算子 $T:E_1\times E_2\to F$ 是线性的, 有界的且是满射.此外,
$N(T)=N(T_1)\times N(T_2)$ (因为 $R(T_1)\cap R(T_2)=\{0\}$).把习题 2.10
应用于 $M=E_1\times\{0\}$, 可见只要 $M+N(T)$ 闭, $T(M)=R(T_1)$ 就是闭的.
但 $M+N(T)=E_1\times N(T_2)$ 确实是闭的.

\solutionhead{2.16}

设 $\pi$ 为由 $E$ 到 $E/L$ 上的典范满射 (见 \S \ref{sec:11.2}).考虑由
$Tx=\pi x$ ($x\in G$) 定义的算子 $T:G\to E/L$.我们有
\begin{equation*}
\dist(x,N(T))=\dist(x,G\cap L)\le C\dist(x,L)=C\norm{Tx}\quad\forall x\in G.
\end{equation*}
由此推出 (见习题 2.14) $R(T)=\pi(G)$ 是闭的.因此 $\pi^{-1}[\pi(G)]=G+L$
是闭的.

\solutionhead{2.19}

回忆 $N(A^\star)=R(A)^\perp$.

1. 设 $u\in N(A)$ 且 $v\in D(A)$; 对一切 $t\in\R$, 有
\begin{equation*}
\dual{A(u+tv)}{u+tv}\ge-C\norm{A(u+tv)}^2,
\end{equation*}
这蕴含 $\dual{Av}{u}=0$.因此 $N(A)\subset R(A)^\perp$.

2. 配以图范数的 $D(A)$ 是 Banach 空间.配以 $E^\star$ 范数的 $R(A)$ 是
Banach 空间.算子 $A:D(A)\to R(A)$ 满足开映射定理的假设.因此存在常数 $C$,
使得对一切 $f\in R(A)$, 存在 $v\in D(A)$ 满足 $Av=f$ 且
$\norm{v}_{D(A)}\le C\norm{f}$.特别地, $\norm{v}\le C\norm{f}$.给定
$u\in D(A)$, 把上述结果应用于 $f=Au$, 可见存在某个 $v\in D(A)$ 使得
$Au=Av$ 且 $\norm{v}\le C\norm{Au}$.由于 $u-v\in N(A)\subset R(A)^\perp$,
我们有
\begin{equation*}
\dual{Au}{u}=\dual{Av}{u}=\dual{Av}{v}\ge-\norm{Av}\norm{v}\ge-C\norm{Au}^2.
\end{equation*}

\solutionhead{2.21}

1. 区分两种情形: 情形 (i): $f(a)=1$.则 $N(A)=\R a$ 且 $R(A)=N(f)$.
情形 (ii): $f(a)\ne1$.则 $N(A)=\{0\}$ 且 $R(A)=E$.

2. $A$ 不是闭的.否则闭图像定理将蕴含 $A$ 有界, 从而 $f$ 连续.

3. $D(A^\star)=\{u\in E^\star;\ \dual{u}{a}=0\}$, 且对一切
$u\in D(A^\star)$, $A^\star u=u$.

4. $N(A^\star)=\{0\}$ 且 $R(A^\star)=\{u\in E^\star;\ \dual{u}{a}=0\}$.

5. $R(A)^\perp=\{0\}$ 且 $R(A^\star)^\perp=\R a$ (注意 $N(f)$ 在 $E$ 中稠密;
见习题 1.6).由此推出 $N(A^\star)=R(A)^\perp$ 且
$N(A)\subset R(A^\star)^\perp$.注意到在情形 (ii) 中,
$N(A)\ne R(A^\star)^\perp$.

6. 若 $A$ 不是闭的, 可能出现 $N(A)\ne R(A^\star)^\perp$ 的情形.

\solutionhead{2.22}

1. 显然 $D(A)$ 在 $E$ 中稠密.为了验证 $A$ 是闭的, 设 $(u^j)$ 为 $D(A)$
中的序列使得在 $E$ 中 $u^j\to u$ 且 $Au^j\to f$.由此推出对每个 $n$, 当
$j\to\infty$ 时 $u_n^j\to u_n$, 且 $nu_n^j\to f_n$.因此对一切 $n$,
$nu_n=f_n$, 从而 $u\in D(A)$ 且 $Au=f$.

2. $D(A^\star)=\{v=(v_n)\in\ell^\infty;\ (nv_n)\in\ell^\infty\}$,
$A^\star v=(nv_n)$, 且 $\overline{D(A^\star)}=c_0$.

\solutionhead{2.24}

1. 我们有 $D(B^\star)=\{v\in G^\star;\ T^\star v\in D(A^\star)\}$ 且
$B^\star=A^\star T^\star$.

2. 若 $D(A)\ne E$ 且 $T=0$, 则 $B$ 不是闭的.事实上, 设 $(u_n)$ 为 $D(A)$
中的序列使得 $u_n\to u$, 其中 $u\notin D(A)$.则 $Bu_n\to0$, 但 $u\notin D(B)$.

\solutionhead{2.25}

2. 由推论  \ref{cor:2.7}, $T^{-1}\in\LL(F,E)$.由于 $T^{-1}T=I_E$ 且
$TT^{-1}=I_F$, 推出 $T^\star(T^{-1})^\star=I_{E^\star}$ 以及
$(T^{-1})^\star T^\star=I_{F^\star}$.

\solutionhead{2.26}

我们有
\begin{equation*}
\varphi^\star(T^\star f)=\sup_{x\in E}\{\dual{T^\star f}{x}-\varphi(x)\}
=\sup_{y\in R(T)}\{\dual{f}{y}-\psi(y)\}
=-\inf_{y\in F}\{\psi(y)+\zeta(y)\},
\end{equation*}
其中 $\zeta(y)=-\dual{f}{y}+I_{R(T)}(y)$.应用定理  \ref{thm:1.12}, 我们得到
\begin{equation*}
\varphi^\star(T^\star f)=\min_{g\in F^\star}\{\zeta^\star(g)+\psi^\star(-g)\}.
\end{equation*}
但
\begin{equation*}
\zeta^\star(g)=\begin{cases}
0 & \text{若}\ f+g\in R(T)^\perp,\\
+\infty & \text{若}\ f+g\notin R(T)^\perp,
\end{cases}
\end{equation*}
从而
\begin{equation*}
\varphi^\star(T^\star f)=\min_{f+g\in N(T^\star)}\psi^\star(-g)
=\min_{h\in N(T^\star)}\psi^\star(f-h).
\end{equation*}

\solutionhead{2.27}

令 $G=E\times X$ 并考虑算子 $S(x,y)=Tx+y:G\to F$.应用开映射定理, 我们知道
$S$ 是开映射, 从而 $S(E\times(X\setminus\{0\}))=R(T)+(X\setminus\{0\})$ 在
$F$ 中开.因此它的补集 $R(T)$ 是闭的.
\solutionhead{3.1}

应用推论  \ref{cor:2.4}.

\solutionhead{3.2}

注意对一切 $f\in E^\star$, $\dual{f}{\sigma_n}=\dfrac1n\sum_{i=1}^n\dual{f}{x_i}$.
由于 $\dual{f}{x_n}\to\dual{f}{x}$, 推出 $\dual{f}{\sigma_n}\to\dual{f}{x}$.

\solutionhead{3.4}

1. 令 $G_n=\operatorname{conv}\left(\bigcup_{i=n}^\infty\{x_i\}\right)$.由于对拓扑
$\sigma(E,E^\star)$ 有 $x_n\rightharpoonup x$, 由此推出对一切 $n$,
$x\in\overline{G_n}^{\sigma(E,E^\star)}$.另一方面, $G_n$ 是凸集, 其弱拓扑
$\sigma(E,E^\star)$ 闭包与强拓扑闭包重合 (见习题 3.3).因此对一切 $n$,
$x\in\overline{G_n}$ ($G_n$ 的强闭包), 从而存在序列 $(y_n)$ 使得对一切 $n$,
$y_n\in G_n$ 且 $y_n\to x$ 强收敛.

2. 存在 $E$ 中的序列 $(u_k)$ 使得 $u_k\to x$ 且对一切 $k$,
$u_k\in\operatorname{conv}\left(\bigcup_{i=1}^\infty\{x_i\}\right)$.因此存在
严格递增的整数序列 $(n_k)$ 使得对一切 $k$,
$u_k\in\operatorname{conv}\left(\bigcup_{i=1}^{n_k}\{x_i\}\right)$.由
$n_k\le n<n_{k+1}$ 时 $z_n=u_k$ (以及 $1\le n<n_1$ 时 $z_n=x_1$) 定义的序列
$(z_n)$ 具有所有要求的性质.

\solutionhead{3.7}

1. 设 $x\notin A+B$.我们将构造 $0$ 的关于 $\sigma(E,E^\star)$ 的邻域 $W$
使得 $(x+W)\cap(A+B)=\emptyset$.

对每个 $y\in B$, 存在 $0$ 的凸邻域 $V(y)$ 使得 $(x+V(y))\cap(A+y)=\emptyset$
(因为 $A+y$ 闭且 $x\notin A+y$).显然
\begin{equation*}
B\subset\bigcup_{y\in B}\left(y-\frac12V(y)\right),
\end{equation*}
由于 $B$ 是紧集, 存在某个有限集 $I$ 使得
$B\subset\bigcup_{i\in I}\left(y_i-\frac12V(y_i)\right)$, 其中 $y_i\in B$.
令
\begin{equation*}
W=\frac12\bigcap_{i\in I}V(y_i).
\end{equation*}

我们断言 $(x+W)\cap(A+B)=\emptyset$.事实上, 反设存在某个 $w\in W$ 使得
$x+w\in(A+B)$.从而存在某个 $i\in I$ 使得
\begin{equation*}
x+w\in A+y_i-\frac12V(y_i).
\end{equation*}
由于 $V(y_i)$ 是凸的, 推出存在某个 $w'\in V(y_i)$ 使得 $x+w'\in A+y_i$.
因此 $(x+V(y_i))\cap(A+y_i)\ne\emptyset$; 矛盾.

\runin{注记.} 若 $E^\star$ 可分且 $A$ 有界, 则可以使用序列来证明 $A+B$
是闭的, 因为弱拓扑在有界集上是可度量化的 (见定理  \ref{thm:3.29}).这使论证
稍微更容易一些.事实上, 设 $x_n=a_n+b_n$ 为使得 $x_n\rightharpoonup x$
弱收敛的序列, 其中 $a_n\in A,b_n\in B$.由于 $B$ 弱紧 (且可度量化), 存在
子序列使得 $b_{n_k}\rightharpoonup b$, 其中 $b\in B$.从而
$a_{n_k}\rightharpoonup x-b$.但 $A$ 是弱闭的, 因此 $x-b\in A$, 即
$x\in A+B$.

2. 由第 1 问, $(A-B)$ 是弱闭的, 因此是强闭的.于是可以严格分离 $\{0\}$ 与
$(A-B)$.

\solutionhead{3.8}

1. 由于 $V_k$ 是 $0$ 关于 $\sigma(E,E^\star)$ 的邻域, 可以假设 (见命题
 \ref{prop:3.4}) $V_k$ 具有形式
\begin{equation*}
V_k=\{x\in E;\ \abs{\dual{f}{x}}<\varepsilon_k\ \forall f\in F_k\},
\end{equation*}
其中 $\varepsilon_k>0$, $F_k$ 是 $E^\star$ 的有限子集.因此集合
$F=\bigcup_{k=1}^\infty F_k$ 是可数的.我们断言任何 $g\in E^\star$ 可以写成
$F$ 中元素的有限线性组合.事实上, 令
\begin{equation*}
V=\{x\in E;\ \abs{\dual{g}{x}}<1\}.
\end{equation*}
由于 $V$ 是 $0$ 的关于 $\sigma(E,E^\star)$ 的邻域, 存在某个整数 $m$ 使得
$\{x\in E;\ d(x,0)<1/m\}\subset V$, 从而 $V_m\subset V$.设 $x\in E$ 使得
对一切 $f\in F_m$, $\dual{f}{x}=0$.则对一切 $t\in\R$, $tx\in V_m$, 从而
对一切 $t\in\R$, $tx\in V$, 即 $\dual{g}{x}=0$.应用引理  \ref{lem:3.2},
可见 $g$ 是 $F_m$ 中元素的线性组合.

2. 使用与习题 1.5 第 3 问相同的方法.

3. 若 $\dim E^\star<\infty$, 则 $\dim E^{\star\star}<\infty$; 从而
$\dim E<\infty$ (因为存在由 $E$ 到 $E^{\star\star}$ 中的典范单射).

4. 应用下述引理 (它是引理  \ref{lem:3.2} 的容易推论): 假设
$x_1,x_2,\dots,x_k,y\in E$ 满足
\begin{equation*}
[f\in E^\star;\ \dual{f}{x_i}=0\ \forall i]\implies[\dual{f}{y}=0].
\end{equation*}
则存在常数 $\lambda_1,\lambda_2,\dots,\lambda_k$ 使得
$y=\sum_{i=1}^k\lambda_ix_i$.

\solutionhead{3.9}

1. 应用定理  \ref{thm:1.12}, 取
\begin{equation*}
\varphi(x)=\dual{f_0}{x}+I_{B_E}(x)\quad\text{以及}\quad\psi(x)=I_M(x).
\end{equation*}

2. 注意 $B_{E^\star}$ 关于 $\sigma(E^\star,E)$ 是紧的, 而 $M^\perp$ 关于
$\sigma(E^\star,E)$ 是闭的 (为什么?).

\solutionhead{3.11}

只须对序列作论证 (为什么?).设在 $E$ 中 $x_n\to x$ 强收敛, 且
$Ax_n\not\rightharpoonup Ax$ 关于 $\sigma(E^\star,E)$, 即存在某个 $y\in E$
使得 $\dual{Ax_n}{y}\not\to\dual{Ax}{y}$.由习题 2.6 我们已知 $(Ax_n)$
有界.因此存在子序列使得 $\dual{Ax_{n_k}}{y}\to\ell\ne\dual{Ax}{y}$.应用
$A$ 的单调性, 我们有
\begin{equation*}
\dual{Ax_{n_k}-A(x+ty)}{x_{n_k}-x-ty}\ge0.
\end{equation*}
取极限, 得到
\begin{equation*}
-t\ell+t\dual{A(x+ty)}{y}\ge0,
\end{equation*}
这蕴含 $\ell=\dual{Ax}{y}$; 矛盾.

\solutionhead{3.12}

1. 假设 (A) 蕴含对一切 $f\in E^\star$,
$\varphi^\star(f)\ge R\norm{f}+\dual{f}{x_0}-M$.反之, 设 (B) 成立, 并令
$\psi(f)=\varphi^\star(f)-\dual{f}{x_0}$.我们断言存在常数 $k>0$ 与 $C$ 使得
\begin{equation}
\psi(f)\ge k\norm{f}-C\quad\forall f\in E^\star.
\tag{S1}
\label{eq:S3.12}
\end{equation}
平移后总可以假设 $\psi(0)<\infty$ (见命题  \ref{prop:1.10}).固定
$\alpha>\psi(0)$.利用假设 (B), 我们可以找到某个 $r>0$ 使得对一切满足
$\norm{g}\ge r$ 的 $g\in E^\star$, $\psi(g)\ge\alpha$.给定满足
$\norm{f}\ge r$ 的 $f\in E^\star$, 取 $t=r/\norm{f}$, 写出
$\psi(tf)\le t\psi(f)+(1-t)\psi(0)$.由于 $\norm{tf}=r$, 这导出
\begin{equation*}
\alpha-\psi(0)\le\frac{r}{\norm{f}}(\psi(f)-\psi(0)),
\end{equation*}
这建立了断言 \eqref{eq:S3.12}.对 \eqref{eq:S3.12} 取共轭即得到 (A).

2. 函数 $\psi$ 是凸的且关于弱$^\star$拓扑下半连续 (为什么?).假设 (B) 说
对每个 $\lambda\in\R$, 集合 $\{f\in E^\star;\ \psi(f)\le\lambda\}$ 是有界的.
因此它是弱$^\star$紧的 (由定理  \ref{thm:3.16}), 从而 $\inf_{E^\star}\psi$
可以达到.另一方面,
\begin{equation*}
\inf_{E^\star}\psi=-\sup_{f\in E^\star}\{\dual{f}{x_0}-\varphi^\star(f)\}
=-\varphi^{\star\star}(x_0)=-\varphi(x_0).
\end{equation*}
另外, 也可以把定理  \ref{thm:1.12} 应用于函数 $\varphi$ 与 $I_{\{x_0\}}$
(注意 $\varphi$ 在 $x_0$ 处连续; 见习题 2.1).

\solutionhead{3.13}

1. 对每个固定的 $p$, 对一切 $n$ 有 $x_{p+n}\in K_p$.取极限 ($n\to\infty$),
可见 $x\in K_p$, 因为 $K_p$ 是弱闭的 (见定理  \ref{thm:3.7}).另一方面, 设
$V$ 为 $x$ 的关于拓扑 $\sigma(E,E^\star)$ 的凸邻域.存在整数 $N$ 使得对
一切 $n\ge N$, $x_n\in V$.从而对一切 $n\ge N$, $K_n\subset V$, 因此
$\bigcap_{n=1}^\infty K_n\subset V$.由于这对 $x$ 的任何凸邻域 $V$ 都成立,
推出 $\bigcap_{n=1}^\infty K_n\subset\{x\}$ (为什么?).

2. 设 $V$ 为 $x$ 关于拓扑 $\sigma(E,E^\star)$ 的开邻域.令
$K'_n=K_n\cap(V^c)$.由于 $K_n$ 关于 $\sigma(E,E^\star)$ 是紧的 (为什么?),
推出 $K'_n$ 也关于 $\sigma(E,E^\star)$ 是紧的.另一方面,
$\bigcap_{n=1}^\infty K'_n=\emptyset$, 从而存在某个整数 $N$ 使得
$K'_N=\emptyset$, 即 $K_N\subset V$.

3. 可以假设 $x=0$.考虑渐近锥 $C_n=\bigcap_{\lambda>0}\lambda K_n$.由于
$C_n\subset K_n$, 我们推出 $\bigcap_{n=1}^\infty C_n=\{0\}$.令
$S_E=\{x\in E;\ \norm{x}=1\}$.序列 $(C_n\cap S)$ 递减且
$\bigcap_{n=1}^\infty(C_n\cap S)=\emptyset$.从而由紧性, 对某个 $n_0$ 有
$C_{n_0}\cap S=\emptyset$.因此 $C_{n_0}=\{0\}$, 从而 $K_{n_0}$ 有界
(为什么?).因此 $(x_n)$ 有界, 我们归结为第 2 问.

4. 考虑当 $n$ 为奇数时 $x_n=(0,\dots,0,\underset{(n)}{n},0,\dots)$, 当
$n$ 为偶数时 $x_n=0$ 的序列.

\solutionhead{3.18}

2. 反设在 $\ell^1$ 中对拓扑 $\sigma(\ell^1,\ell^\infty)$ 有
$e_{n_k}\rightharpoonup a$ ($k\to\infty$).于是对一切 $\xi\in\ell^\infty$,
有 $\dual{\xi}{e_{n_k}}\to\dual{\xi}{a}$.考虑由
\begin{equation*}
\xi=(0,\dots,0,\underset{(n_1)}{-1},0,\dots,\underset{(n_2)}{1},0,\dots,\underset{(n_3)}{-1},0,\dots)
\end{equation*}
定义的元素 $\xi\in\ell^\infty$.注意 $\dual{\xi}{e_{n_k}}=(-1)^k$ 当
$k\to\infty$ 时不收敛; 矛盾.

3. 设 $E=\ell^\infty$, 从而 $\ell^1\subset E^\star$.令 $f_n=e_n$, 看作
$E^\star$ 中的序列.我们断言 $(f_n)$ 没有关于 $\sigma(E^\star,E)$ 收敛的
子序列.反设对 $\sigma(E^\star,E)$ 有
$f_{n_k}\overset{*}{\rightharpoonup}f$, 即对一切 $\eta\in E$,
$\dual{f_{n_k}}{\eta}\to\dual{f}{\eta}$.如第 2 问中那样选取 $\eta=\xi$,
可见 $\dual{f_{n_k}}{\xi}=(-1)^k$ 不收敛; 矛盾.这里, 配以拓扑
$\sigma(E^\star,E)$ 的集合 $B_{E^\star}$ 是紧的 (由定理  \ref{thm:3.16}),
但它不是可度量化的.应用定理  \ref{thm:3.28}, 我们也可以说 $E=\ell^\infty$
不可分 (另一种证明见第~
\ref{chap:4}~章注记 8 与命题  \ref{prop:11.17}).

\solutionhead{3.19}

1. 注意若在 $\ell^p$ 中 $x_n\to x$ 强收敛, 则对一切 $\varepsilon>0$, 存在
$I$ 使得对一切 $n$, $\sum_{i=I}^\infty\abs{x_i^n}^p\le\varepsilon^p$.

2. 应用习题 3.17.

3. 空间 $B_E$ 对拓扑 $\sigma(E,E^\star)$ 是可度量化的 (由定理  \ref{thm:3.29}).
因此只须在序列上验证 $A$ 的连续性.

\solutionhead{3.20}

1. 考虑由 $(Tx)(t)=\dual{t}{x}$ ($x\in E,t\in B_{E^\star}=K$) 定义的映射
$T:E\to C(K)$.显然
\begin{equation*}
\norm{Tx}=\sup_{t\in K}\abs{(Tx)(t)}=\norm{x}.
\end{equation*}

2. 由于 $K=B_{E^\star}$ 关于 $\sigma(E^\star,E)$ 是可度量化且紧的, $K$ 中
存在稠密的可数子集 $(t_n)$.考虑由
\begin{equation*}
Sx=(\dual{t_1}{x},\dual{t_2}{x},\dots,\dual{t_n}{x},\dots)
\end{equation*}
定义的映射 $S:E\to\ell^\infty$.验证 $\norm{Sx}_{\ell^\infty}=\norm{x}$.

\solutionhead{3.21}

设 $(a_i)$ 为 $E$ 的稠密可数子集.选取第一个子序列使得当 $k\to\infty$ 时
$\dual{f_{n_k}}{a_1}$ 收敛于一个极限.然后从 $(n_k)$ 中抽取子序列使得
$\dual{f_{n'_k}}{a_2}$ 收敛, 等等.

由标准的对角线法, 可以从序列 $(f_n)$ 中抽取子序列 $(g_k)$ 使得对一切 $i$,
$\dual{g_k}{a_i}\to\ell_i$.由于集合 $(a_i)$ 在 $E$ 中稠密, 容易得到对
一切 $a\in E$, $\dual{g_k}{a}\to\ell_a$.由此推出 $g_k$ 对
$\sigma(E^\star,E)$ 收敛于某个 $g$ (见习题 3.16).

\solutionhead{3.22}

(a) $B_E$ 关于 $\sigma(E,E^\star)$ 是可度量化的 (由定理  \ref{thm:3.29}), 且
$0$ 属于 $S=\{x\in E;\ \norm{x}=1\}$ 关于 $\sigma(E,E^\star)$ 的闭包 (见第
3 章注记 2).

(b) 由于 $\dim E=\infty$, $E$ 中存在可分且满足 $\dim E_0=\infty$ 的闭子空间
$E_0$ (为什么?).注意 $E_0$ 自反, 应用情形 (a) (并结合推论  \ref{cor:3.27}).

\solutionhead{3.25}

反设 $C(K)$ 自反.则 $E=\{u\in C(K);\ u(a)=0\}$ 也自反, 且
$\sup_{u\in B_E}f(u)$ 可以达到.

另一方面, 我们断言 $\sup_{u\in B_E}f(u)=1$.事实上, 对一切 $N$, 存在
$u\in E$ 使得在 $K$ 上 $0\le u\le1$ 且对一切 $i=1,\dots,N$, $u(a_i)=1$
(例如应用 Tietze--Urysohn 定理; 参见例如 J.~Munkres \bref{Munkres}{1}).
因此存在某个 $u\in B_E$ 使得 $f(u)=1$.这导出对一切 $n$, $u(a_n)=1$ 且
$u(a)=0$; 矛盾.

\solutionhead{3.26}

1. 给定 $y\in B_F$, 存在某个整数 $n_1$ 使得 $\norm{y-a_{n_1}}<1/2$.由于集合
$\frac12(a_i)_{i>n_1}$ 在 $\frac12B_F$ 中稠密, 存在某个 $n_2>n_1$ 使得
$\norm{y-a_{n_1}-\frac12a_{n_2}}<1/4$.由归纳法构造递增的整数序列
$n_k\uparrow\infty$ 使得
\begin{equation*}
y=a_{n_1}+\frac12a_{n_2}+\frac14a_{n_3}+\dots+\frac{1}{2^{k-1}}a_{n_k}+\dots.
\end{equation*}

2. 反设 $S\in\LL(F,\ell^1)$ 使得 $TS=I_F$.设 $(y_n)$ 为 $F$ 中的任何序列
满足对一切 $n$, $\norm{y_n}=1$ 且 $y_n\rightharpoonup0$ 弱收敛于 $0$.从而
关于 $\sigma(\ell^1,\ell^\infty)$ 有 $Sy_n\rightharpoonup0$, 进而 $Sy_n\to0$
在 $\ell^1$ 中强收敛 (见问题 8).由此推出 $y_n=TSy_n\to0$; 矛盾.

3. 使用定理  \ref{thm:2.12}.

4. $T^\star:F^\star\to\ell^\infty$ 由
$T^\star v=(\dual{v}{a_1},\dual{v}{a_2},\dots,\dual{v}{a_n},\dots)$ 定义.

\solutionhead{3.27}

$B_{E^\star}$ 关于 $\sigma(E^\star,E)$ 是紧的且可度量化的.因此 $B_{E^\star}$
中存在关于 $\sigma(E^\star,E)$ 稠密的可数子集.

1. 显然对一切 $f\in E^\star$, $\norm{f}\le\norm{f}_1\le\sqrt{2}\norm{f}$.

2. 令 $\abs{f}^2=\sum_{n=1}^\infty\frac{1}{2^n}\abs{\dual{f}{a_n}}^2$.注意
范数 $\abs{\cdot}$ 与标量积相关联 (为什么?), 因此它是严格凸的, 即函数
$f\mapsto\abs{f}^2$ 严格凸.更精确地说, 对一切 $t\in[0,1]$ 与一切
$f,g\in E^\star$, 有
\begin{equation}
\abs{tf+(1-t)g}^2+t(1-t)\abs{f-g}^2=t\abs{f}^2+(1-t)\abs{g}^2.
\tag{S1}
\label{eq:S3.27}
\end{equation}
因此, 函数 $f\mapsto\norm{f}^2+\abs{f}^2$ 也是严格凸的.

3. 与第 2 问相同的方法.注意若对一切 $n$, $\dual{b_n}{x}=0$, 则 $x=0$
(为什么?).

4. 给定 $x\in E$, 令
$[x]=\left\{\sum_{n=1}^\infty\frac{1}{2^n}\abs{\dual{b_n}{x}}^2\right\}^{1/2}$,
并用 $[f]$ 记 $[\cdot]$ 在 $E^\star$ 上的对偶范数.注意 $[f]$ 也满足恒等式
\eqref{eq:S3.27}.事实上, 我们有
\begin{align*}
\frac12[tf+(1-t)g]^2&=\sup_{x\in E}\left\{\dual{tf+(1-t)g}{x}-\frac12[x]^2\right\},\\
\frac12[f-g]^2&=\sup_{y\in E}\left\{\dual{f-g}{y}-\frac12[y]^2\right\},
\end{align*}
从而
\begin{align*}
&\frac12[tf+(1-t)g]^2+\frac12t(1-t)[f-g]^2\\
&=\sup_{x,y}\left\{\dual{tf+(1-t)g}{x}+t(1-t)\dual{f-g}{y}-\frac12[x]^2-\frac12t(1-t)[y]^2\right\}.
\end{align*}
作变量代换 $x=t\xi+(1-t)\eta$ 与 $y=\xi-\eta$, 我们得出 \eqref{eq:S3.27}
成立.应用习题 1.23 第 3 问, 可见
\begin{equation*}
\norm{f}_2^2=\inf_{h\in E^\star}\{\norm{f-h}_1^2+[h]^2\}=\min_{h\in E^\star}\{\norm{f-h}_1^2+[h]^2\}.
\end{equation*}
我们断言函数 $f\mapsto\norm{f}_2^2$ 严格凸.事实上, 给定 $f,g\in E^\star$,
固定 $h_1,h_2\in E^\star$ 使得 $\norm{f}_2^2=\norm{f-h_1}_1^2+[h_1]^2$,
$\norm{g}_2^2=\norm{g-h_2}_1^2+[h_2]^2$.对每个 $t\in(0,1)$, 有
\begin{align*}
\norm{tf+(1-t)g}_2^2&\le\norm{tf+(1-t)g-(th_1+(1-t)h_2)}_1^2+[th_1+(1-t)h_2]^2\\
&<t\norm{f}_2^2+(1-t)\norm{g}_2^2,
\end{align*}
除非 $f-h_1=g-h_2$ 且 $h_1=h_2$, 即 $f=g$.

\solutionhead{3.28}

由于 $E$ 自反, $\sup_{x\in B_E}\dual{f}{x}$ 由某个唯一点 $x_0\in B_E$ 达到.
则 $x=x_0\norm{f}$ 满足 $f\in F(x)$.

另外, 也可以考虑由 $E^\star$ 到 $E^{\star\star}$ 中的对偶映射 $F^\star$.
集合 $F^\star(f)$ 非空 (由推论  \ref{cor:1.3}).固定任何 $\xi\in F^\star(f)$.
由于 $E$ 自反, 存在某个 $x\in E$ 使得 $Jx=\xi$ ($J$ 是由 $E$ 到
$E^{\star\star}$ 中的典范单射).我们有 $\norm{\xi}=\norm{f}=\norm{x}$ 以及
$\dual{\xi}{f}=\norm{f}^2=\dual{f}{x}$.因此 $f\in F(x)$.

\runin{唯一性.} 设 $x_1,x_2$ 使得 $f\in F(x_1)$ 且 $f\in F(x_2)$.则
$\norm{x_1}=\norm{x_2}=\norm{f}$, 因此若 $x_1\ne x_2$, 有
$\norm{\dfrac{x_1+x_2}{2}}<\norm{f}$.另一方面,
$\dual{f}{x_1}=\dual{f}{x_2}=\norm{f}^2$, 从而若 $x_1\ne x_2$, 有
$\norm{f}^2=\dual{f}{\dfrac{x_1+x_2}{2}}<\norm{f}^2$, 矛盾.

\solutionhead{3.29}

1. 反设存在 $M_0>0$, $\varepsilon_0>0$ 与两个序列 $(x_n),(y_n)$ 使得
$\norm{x_n}\le M$, $\norm{y_n}\le M$, $\norm{x_n-y_n}>\varepsilon_0$, 以及
\begin{equation}
\norm{\frac{x_n+y_n}{2}}^2>\frac12\norm{x_n}^2+\frac12\norm{y_n}^2-\frac1n.
\tag{S1}
\label{eq:S3.29}
\end{equation}
考虑仍记作 $(x_n)$ 与 $(y_n)$ 的子序列使得 $\norm{x_n}\to a$ 且
$\norm{y_n}\to b$.我们发现 $a+b\ge\varepsilon_0$ 且
$\frac12a^2+\frac12b^2\le\left(\frac{a+b}{2}\right)^2$.因此 $a=b\ne0$.

令 $x'_n=\dfrac{x_n}{\norm{x_n}}$ 与 $y'_n=\dfrac{y_n}{\norm{y_n}}$.对足够大的
$n$, 有 $\norm{x'_n-y'_n}\ge(\varepsilon_0/a)+o(1)$.由一致凸性存在
$\delta_0>0$ 使得 $\norm{\dfrac{x'_n+y'_n}{2}}\le1-\delta_0$.从而
\begin{equation*}
\norm{\frac{x_n+y_n}{2}}\le a(1-\delta_0)+o(1).
\end{equation*}
由 \eqref{eq:S3.29} 我们有 $\norm{\dfrac{x_n+y_n}{2}}^2\ge a^2+o(1)$.因此
$a^2\le a^2(1-\delta_0)^2+o(1)$; 矛盾.

\solutionhead{3.32}

1. 下确界可以达到, 因为 $E$ 自反, 我们可以应用推论  \ref{cor:3.23}.唯一性
来自空间 $E$ 严格凸的事实, 从而函数 $y\mapsto\norm{y-x}^2$ 严格凸.

2. 设 $(y_n)$ 为极小化序列; 令 $d_n=\norm{x-y_n}$ 与 $d=\inf_{y\in C}\norm{x-y}$,
从而 $d_n\to d$.设 $(y_{n_k})$ 为使得 $y_{n_k}\rightharpoonup z$ 弱收敛的序列.
从而 $z\in C$ 且 $\norm{x-z}\le d$ (为什么?).由此推出
$x-y_{n_k}\rightharpoonup x-z$ 弱收敛且 $\norm{x-y_{n_k}}\to d=\norm{x-z}$,
因此 (见命题  \ref{prop:3.32}) $y_{n_k}\to z$ 强收敛.极限的唯一性蕴含整个序列
$(y_n)$ 强收敛于 $P_Cx$.这个论证是标准的, 我们将多次使用它.为方便读者起见,
我们回忆它: 反设 $(y_n)$ 不收敛于 $y=P_Cx$.则存在 $\varepsilon>0$ 与子序列
$(y_{m_j})$ 使得对一切 $j$, $\norm{y_{m_j}-y}\ge\varepsilon$.从 $(y_{m_j})$
中我们抽取 (按上述论证) 进一步的子序列 $(y_{n_k})$ 使得 $y_{n_k}\to P_Cx$.
由于 $(y_{n_k})$ 是 $(y_{m_j})$ 的子序列, 对一切 $k$ 有
$\norm{y_{n_k}-y}\ge\varepsilon$, 从而 $\norm{P_Cx-y}\ge\varepsilon$.矛盾.

3 与 4. 反设存在某个 $\varepsilon_0>0$ 与序列 $(x_n),(y_n)$ 使得
$\norm{x_n}\le M$, $\norm{y_n}\le M$, $\norm{x_n-y_n}\to0$, 且对一切 $n$,
$\norm{P_Cx_n-P_Cy_n}\ge\varepsilon_0$.我们有
\begin{equation*}
\norm{x_n-P_Cx_n}\le\norm{\frac{x_n-P_Cx_n+P_Cy_n}{2}}\le\norm{\frac{x_n+y_n}{2}-\frac{P_Cx_n+P_Cy_n}{2}}+o(1),
\end{equation*}
同理
\begin{equation*}
\norm{y_n-P_Cy_n}\le\norm{\frac{x_n+y_n}{2}-\frac{P_Cx_n+P_Cy_n}{2}}+o(1).
\end{equation*}
由此推出
\begin{equation}
\frac12\norm{x_n-P_Cx_n}^2+\frac12\norm{y_n-P_Cy_n}^2\le\norm{\frac{x_n+y_n}{2}-\frac{P_Cx_n+P_Cy_n}{2}}^2+o(1).
\tag{S1}
\label{eq:S3.32a}
\end{equation}
另一方面, 若令 $a_n=x_n-P_Cx_n$ 与 $b_n=y_n-P_Cy_n$, 则
$\norm{a_n-b_n}\ge\varepsilon_0+o(1)$ 且 $\norm{a_n}\le M', \norm{b_n}\le M'$.
利用习题 3.29, 我们知道存在某个 $\delta_0>0$ 使得
\begin{equation*}
\norm{\frac{a_n+b_n}{2}}^2\le\frac12\norm{a_n}^2+\frac12\norm{b_n}^2-\delta_0,
\end{equation*}
即
\begin{equation}
\norm{\frac{x_n+y_n}{2}-\frac{P_Cx_n+P_Cy_n}{2}}^2\le\frac12\norm{x_n-P_Cx_n}^2+\frac12\norm{y_n-P_Cy_n}^2-\delta_0.
\tag{S2}
\label{eq:S3.32b}
\end{equation}
结合 \eqref{eq:S3.32a} 与 \eqref{eq:S3.32b} 导出矛盾.

5. 与第 1 问相同的论证.

6. 我们有
\begin{equation}
n\norm{y_n-x}^2+\varphi(y_n)\le n\norm{y-x}^2+\varphi(y)\quad\forall y\in D(\varphi).
\tag{S3}
\label{eq:S3.32c}
\end{equation}
由于 $\varphi$ 被一个连续仿射函数由下方控制 (见命题  \ref{prop:1.10}), 我们看到
当 $n\to\infty$ 时 $(y_n)$ 保持有界.设 $(y_{n_k})$ 为使得
$y_{n_k}\rightharpoonup z$ 弱收敛的子序列.注意 $z\in\overline{D(\varphi)}$
(为什么?).由 \eqref{eq:S3.32c} 我们得到对一切 $y\in D(\varphi)$,
$\norm{z-x}\le\norm{y-x}$, 从而对一切 $y\in\overline{D(\varphi)}$ 成立.因此
$z=P_Cx$, 其中 $C=\overline{D(\varphi)}$.再次使用 \eqref{eq:S3.32c} 给出
\begin{equation*}
\limsup_{n\to\infty}\norm{y_n-x}\le\norm{y-x}\quad\forall y\in D(\varphi),
\end{equation*}
特别地对 $y=z$ 成立.我们得出 $y_{n_k}\to z$ 强收敛, 最后极限的唯一性表明
(如上) 整个序列 $(y_n)$ 强收敛于 $P_Cx$.
\solutionhead{4.3}

2. 注意 $h_n=\frac12(\abs{f_n-g_n}+f_n+g_n)$.

3. 注意 $f_ng_n-fg=(f_n-f)g_n+f(g_n-g)$, 且由控制收敛定理, 在 $L^p(\Omega)$
中 $f(g_n-g)\to0$.

\solutionhead{4.5}

1. 回忆 $L^1(\Omega)\cap L^\infty(\Omega)\subset L^p(\Omega)$, 更精确地
$\norm{f}_p^p\le\norm{f}_\infty^{p-1}\norm{f}_1$.由于 $\Omega$ 是 $\sigma$-有限
的, 可以写成 $\Omega=\bigcup_n\Omega_n$, 其中对一切 $n$, $\abs{\Omega_n}<\infty$.
给定 $f\in L^p(\Omega)$, 验证 $f_n=\chi_{\Omega_n}T_nf\in L^1(\Omega)\cap L^\infty(\Omega)$,
且当 $n\to\infty$ 时在 $L^p(\Omega)$ 中 $f_n\to f$.

2. 设 $(f_n)$ 为 $L^p(\Omega)\cap L^q(\Omega)$ 中的序列, 使得在 $L^p(\Omega)$
中 $f_n\to f$ 且 $\norm{f_n}_q\le1$.我们假设 (通过抽取子序列) $f_n\to f$
几乎处处 (见定理  \ref{thm:4.9}).由 Fatou 引理推出 $f\in L^q(\Omega)$ 且
$\norm{f}_q\le1$.

3. 由第 2 问我们已知 $f\in L^q(\Omega)$, 从而对介于 $p$ 与 $q$ 之间的每个
$r$, $f\in L^r(\Omega)$.另一方面, 可以写出 $\dfrac1r=\dfrac\alpha p+\dfrac{1-\alpha}{q}$,
其中 $0<\alpha\le1$, 我们得到
\begin{equation*}
\norm{f_n-f}_r\le\norm{f_n-f}_p^\alpha\norm{f_n-f}_q^{1-\alpha}
\le\norm{f_n-f}_p^\alpha(2C)^{1-\alpha}.
\end{equation*}

\solutionhead{4.6}

1. 我们有 $\norm{f}_p\le\norm{f}_\infty\abs{\Omega}^{1/p}$, 从而
$\limsup_{p\to\infty}\norm{f}_p\le\norm{f}_\infty$.另一方面, 固定
$0<k<\norm{f}_\infty$, 并令 $A=\{x\in\Omega;\ \abs{f(x)}>k\}$.显然
$\abs{A}\ne0$ 且 $\norm{f}_p\ge k\abs{A}^{1/p}$.由此推出
$\liminf_{p\to\infty}\norm{f}_p\ge k$, 因此
$\liminf_{p\to\infty}\norm{f}_p\ge\norm{f}_\infty$.

2. 固定 $k>C$, 设 $A$ 如上定义.则 $k^p\abs{A}\le\norm{f}_p^p\le C^p$, 从而对
一切 $p\ge1$, $\abs{A}\le(C/k)^p$.令 $p\to\infty$, 可见 $\abs{A}=0$.

3. $f(x)=\log\abs{x}$.

\solutionhead{4.7}

考虑由 $Tu=au$ 定义的算子 $T:L^p(\Omega)\to L^q(\Omega)$.我们断言 $T$ 的
图是闭的.事实上, 设 $(u_n)$ 为 $L^p(\Omega)$ 中的序列使得在 $L^p(\Omega)$
中 $u_n\to u$ 且在 $L^q(\Omega)$ 中 $au_n\to f$.抽取子序列可以假设
$u_n\to u$ 几乎处处且 $au_n\to f$ 几乎处处.从而 $f=au$ 几乎处处, 于是
$f=Tu$.由闭图像定理 (定理  \ref{thm:2.9}) 推出 $T$ 有界, 因此存在常数 $C$
使得
\begin{equation}
\norm{au}_q\le C\norm{u}_p\quad\forall u\in L^p(\Omega).
\tag{S1}
\label{eq:S4.7}
\end{equation}

情形 1: $p<\infty$.由 \eqref{eq:S4.7} 推出
\begin{equation*}
\int_\Omega\abs{a}^q\abs{v}\le C^q\norm{v}_{p/q}\quad\forall v\in L^{p/q}(\Omega).
\end{equation*}
因此映射 $v\mapsto\int_\Omega\abs{a}^qv$ 是 $L^{p/q}(\Omega)$ 上的连续线性
泛函, 从而 $\abs{a}^q\in L^{(p/q)'}(\Omega)$.

情形 2: $p=\infty$.在 \eqref{eq:S4.7} 中选取 $u\equiv1$.

\solutionhead{4.8}

1. 配以范数 $\norm{\cdot}_1$ 的 $X$ 是 Banach 空间.对每个 $n$, $X_n$ 是
$X$ 的闭子集 (见习题 4.5).另一方面, $X=\bigcup_nX_n$.事实上, 对每个
$f\in X$, 存在某个 $q>1$ 使得 $f\in L^q(\Omega)$.因此只要
$1+(1/n)\le q$, 就有 $f\in L^{1+1/n}(\Omega)$, 此外
\begin{equation*}
\norm{f}_{1+1/n}\le\norm{f}_1^{\alpha_n}\norm{f}_q^{1-\alpha_n},
\end{equation*}
其中 $\dfrac{1}{1+(1/n)}=\dfrac{\alpha_n}{1}+\dfrac{1-\alpha_n}{q}$.由
Baire 纲定理推出存在某个整数 $n_0$ 使得 $\operatorname{Int}X_{n_0}\ne\emptyset$.
因此 $X\subset L^{1+1/n_0}(\Omega)$.

2. 恒等映射 $I:X\to L^p(\Omega)$ 是图为闭的线性算子.因此它是有界算子.

\solutionhead{4.9}

对每个 $t\in\R$, 在 $\Omega$ 上几乎处处有 $f(x)t\le j(f(x))+j^\star(t)$,
积分得到
\begin{equation*}
\left(\frac{1}{\abs{\Omega}}\int_\Omega f\right)t
\le\frac{1}{\abs{\Omega}}\int_\Omega j(f)+j^\star(t).
\end{equation*}
因此
\begin{equation*}
j\left(\frac{1}{\abs{\Omega}}\int_\Omega f\right)
=\sup_{t\in\R}\left\{\left(\frac{1}{\abs{\Omega}}\int_\Omega f\right)t-j^\star(t)\right\}
\le\frac{1}{\abs{\Omega}}\int_\Omega j(f).
\end{equation*}

\solutionhead{4.10}

1. 设 $u_1,u_2\in D(J)$ 且 $t\in[0,1]$.函数 $x\mapsto j(tu_1(x)+(1-t)u_2(x))$
是可测的 (因为 $j$ 连续).另一方面, $j(tu_1+(1-t)u_2)\le tj(u_1)+(1-t)j(u_2)$.
回忆存在常数 $a$ 与 $b$ 使得对一切 $s\in\R$, $j(s)\ge as+b$ (见命题  \ref{prop:1.10}).
由此推出 $j(tu_1+(1-t)u_2)\in L^1(\Omega)$, 且
$J(tu_1+(1-t)u_2)\le tJ(u_1)+(1-t)J(u_2)$.

2. 先设 $j\ge0$.我们断言对每个 $\lambda\in\R$, 集合
$\{u\in L^p(\Omega);\ J(u)\le\lambda\}$ 是闭的.事实上, 设 $(u_n)$ 为
$L^p(\Omega)$ 中的序列使得在 $L^p(\Omega)$ 中 $u_n\to u$ 且 $\int_\Omega j(u_n)\le\lambda$.
抽取子序列可以假设 $u_n\to u$ 几乎处处.由 Fatou 引理推出 $j(u)\in L^1(\Omega)$
且 $\int_\Omega j(u)\le\lambda$.因此 $J$ 下半连续.在一般情形下, 设
$\tilde j(s)=j(s)-(as+b)\ge0$.我们已知 $\widetilde J$ 下半连续, 从而
$J(u)=\widetilde J(u)+a\int_\Omega u+b\abs{\Omega}$ 也下半连续.

3. 我们先断言: 对使 $j^\star(f)\in L^1(\Omega)$ 的每个 $f\in L^{p'}(\Omega)$, 有
\begin{equation*}
J^\star(f)\le\int_\Omega j^\star(f).
\end{equation*}
事实上, 在 $\Omega$ 上几乎处处有 $fu-j(u)\le j^\star(f)\ \forall u\in L^p(\Omega)$,
从而
\begin{equation*}
\sup_{u\in D(J)}\left\{\int_\Omega fu-J(u)\right\}\le\int_\Omega j^\star(f).
\end{equation*}

反向不等式的证明更精细, 需要某种``正则化''步骤.先设 $1<p<\infty$, 并令
\begin{equation*}
j_n(t)=j(t)+\frac1n\abs{t}^p,\quad t\in\R.
\end{equation*}
我们断言对一切 $f\in L^{p'}(\Omega)$,
\begin{equation}
J_n^\star(f)=\int_\Omega j_n^\star(f).
\tag{S1}
\label{eq:S4.10}
\end{equation}
事实上, 设 $f\in L^{p'}(\Omega)$.对几乎处处固定的 $x\in\Omega$,
$\sup_{u\in\R}\left\{f(x)u-j(u)-\frac1n\abs{u}^p\right\}$ 由某个唯一的元素
$u=u(x)$ 达到.显然我们有
$j(u(x))+\frac1n\abs{u(x)}^p-f(x)u(x)\le j(0)$.由此推出 $u\in L^p(\Omega)$
且 $j(u)\in L^1(\Omega)$ (为什么?).我们得出
\begin{equation*}
J_n^\star(f)=\sup_{v\in D(J)}\left\{\int_\Omega fv-J_n(v)\right\}
\ge\int_\Omega\left(fu-j(u)-\frac1n\abs{u}^p\right)=\int_\Omega j_n^\star(f).
\end{equation*}
由于我们已经建立了反向不等式, 可见 \eqref{eq:S4.10} 成立.

接着令 $n\uparrow\infty$.显然 $J\le J_n$, 从而 $J_n^\star\le J^\star$, 即
$\int_\Omega j_n^\star(f)\le J^\star(f)$.我们断言对每个 $s\in\R$, 当
$n\uparrow\infty$ 时 $j_n^\star(s)\uparrow j^\star(s)$.事实上, 我们知道
$j_n^\star=j^\star\mathbin{\triangledown}\left(\frac1n\abs{\cdot}^p\right)^\star$
(见习题 1.23), 然后可以像习题 1.24 那样论证.由单调收敛定理得出: 若
$f\in D(J^\star)$, 则 $j^\star(f)\in L^1(\Omega)$ 且
$\int_\Omega j^\star(f)\le J^\star(f)$.

最后, 若 $p=1$, 可以例如使用 $j_n(t)=j(t)+\frac1nt^2$ 来修改上述方法.

4. 先假设在 $\Omega$ 上几乎处处有 $f(x)\in\partial j(u(x))$, 我们有
\begin{equation*}
j(v)-j(u(x))\ge f(x)(v-u(x))\quad\forall v\in\R,\ \text{在}\ \Omega\ \text{上几乎处处}.
\end{equation*}
选取 $v=0$, 可见 $j(u)\in L^1(\Omega)$, 从而对一切 $v\in D(J)$,
$J(v)-J(u)\ge\int_\Omega f(v-u)$.

反之, 假设 $f\in\partial J(u)$.则我们有 $J(u)+J^\star(f)=\int_\Omega fu$.
因此 $j(u)\in L^1(\Omega)$, $j^\star(f)\in L^1(\Omega)$, 且
$\int_\Omega\{j(u)+j^\star(f)-fu\}=0$.由于几乎处处有 $j(u)+j^\star(f)-fu\ge0$,
我们发现几乎处处有 $j(u)+j^\star(f)-fu=0$, 即几乎处处 $f(x)\in\partial j(u(x))$.

\solutionhead{4.11}

令 $f=u^\alpha$, $g=v^\alpha$, 以及 $p=1/\alpha$.须证明
\begin{equation}
\left(\int_\Omega f\right)^p+\left(\int_\Omega g\right)^p
\le\left[\int_\Omega(f^p+g^p)^{1/p}\right]^p.
\tag{S1}
\label{eq:S4.11}
\end{equation}
令 $a=\int_\Omega f$ 与 $b=\int_\Omega g$, 于是我们有
\begin{align*}
a^p+b^p&=\int_\Omega(a^{p-1}f+b^{p-1}g)\le\int_\Omega(a^p+b^p)^{1/p'}(f^p+g^p)^{1/p}\\
&=(a^p+b^p)^{1/p'}\int_\Omega(f^p+g^p)^{1/p}.
\end{align*}
由此推出 $(a^p+b^p)^{1/p}\le\int_\Omega(f^p+g^p)^{1/p}$, 即 \eqref{eq:S4.11}.

\solutionhead{4.12}

1. 只须证明
\begin{equation*}
\inf_{t\in[-1,+1]}\left\{\frac{(\abs{t}^p+1)^{1-s}(\abs{t}^p+1-2\abs{\frac{t+1}{2}}^p)^s}{\abs{t-1}^p}\right\}>0,
\end{equation*}
或等价地证明
\begin{equation*}
\inf_{t\in[-1,+1]}\left\{\frac{\abs{t}^p+1-2\abs{\frac{t+1}{2}}^p}{\abs{t-1}^2}\right\}>0.
\end{equation*}
但函数 $\phi(t)=\abs{t}^p+1-2\abs{\frac{t+1}{2}}^p$ 满足对一切 $t\in[-1,+1)$,
$\phi(t)>0$, 且 $\phi(1)=\phi'(1)=0$, $\phi''(1)>0$.

\solutionhead{4.14}

1 与 2. 令 $g_n=\chi_{S_n(\alpha)}$, 从而 $g_n\to0$ 几乎处处且 $\abs{g_n}\le1$.
由控制收敛定理 (因为 $\abs{\Omega}<\infty$) 推出 $\int_\Omega g_n\to0$, 即
$\abs{S_n(\alpha)}\to0$.

3. 给定任何整数 $m\ge1$, 我们可以把第 2 问应用于 $\alpha=1/m$, 找到整数
$N_m$ 使得 $\abs{S_{N_m}(1/m)}<\delta/2^m$.令 $\Omega_m=S_{N_m}(1/m)$,
我们得到对一切 $k\ge N_m$ 及一切 $x\in\Omega\setminus\Omega_m$,
$\abs{f_k(x)-f(x)}\le1/m$.最后, 令 $A=\bigcup_{m=1}^\infty\Omega_m$, 从而
$\abs{A}<\delta$.我们断言 $f_n\to f$ 在 $\Omega\setminus A$ 上一致收敛.
事实上, 给定 $\varepsilon>0$, 固定满足 $m_0>1/\varepsilon$ 的整数 $m_0$.显然
对一切 $k\ge N_{m_0}$ 及一切 $x\in\Omega\setminus\Omega_{m_0}$,
$\abs{f_k(x)-f(x)}<\varepsilon$, 从而对一切 $k\ge N_{m_0}$ 及一切
$x\in\Omega\setminus A$, $\abs{f_k(x)-f(x)}<\varepsilon$.

4. 给定 $\varepsilon>0$, 先利用 (i) 固定某个 $\delta>0$, 然后利用第 3 问
固定某个 $A$.我们得到对一切 $n$, $\int_A\abs{f_n}^p\le\varepsilon$, 且
$f_n\to f$ 在 $\Omega\setminus A$ 上一致收敛.由 Fatou 引理推出
$\int_A\abs{f}^p\le\varepsilon$, 从而
\begin{equation*}
\int_\Omega\abs{f_n-f}^p=\int_A\abs{f_n-f}^p+\int_{\Omega\setminus A}\abs{f_n-f}^p
\le2^p\varepsilon+\abs{\Omega}\norm{f_n-f}_{L^\infty(\Omega\setminus A)}^p.
\end{equation*}

\solutionhead{4.15}

1(iv). 注意对一切 $\phi\in C_c(\Omega)$, $\int_\Omega f_n\phi\to0$.反设
$f_{n_k}\rightharpoonup f$ 弱收敛于 $\sigma(L^1,L^\infty)$.由此推出对一切
$\phi\in C_c(\Omega)$, $\int_\Omega f\phi=0$, 从而 (由推论  \ref{cor:4.24})
$f=0$ 几乎处处.此外 $\int_\Omega f_{n_k}\to\int_\Omega f=0$; 但
$\int_\Omega f_{n_k}=\int_0^{n_k}e^{-t}dt\to1$; 矛盾.

2(iv). 注意对一切 $\phi\in C_c(\Omega)$, $\int_\Omega g_n\phi\to0$, 并利用
$C_c(\Omega)$ 在 $L^{p'}(\Omega)$ 中稠密的事实 (因为 $p'<\infty$).

\solutionhead{4.16}

1. 我们先验证: 若序列 $(f_n)$ 满足
\begin{equation}
f_n\rightharpoonup\tilde f\quad\text{弱收敛于}\ \sigma(L^p,L^{p'})
\tag{S1}
\label{eq:S4.16a}
\end{equation}
且
\begin{equation}
f_n\to f\quad\text{几乎处处},
\tag{S2}
\label{eq:S4.16b}
\end{equation}
则 $f=\tilde f$ 几乎处处.

事实上, 由习题 3.4 我们知道存在 $L^p(\Omega)$ 中的序列 $(g_n)$ 使得
$g_n\in\operatorname{conv}\{f_n,f_{n+1},\dots\}$ 且在 $L^p(\Omega)$ 中
$g_n\to\tilde f$ 强收敛.由 \eqref{eq:S4.16b} 推出 $g_n\to f$ 几乎处处.
另一方面 (由定理  \ref{thm:4.9}), 存在子序列 $(g_{n_k})$ 使得
$g_{n_k}\to\tilde f$ 几乎处处.因此 $f=\tilde f$ 几乎处处.

现在在假设 (i) 与 (ii) 下验证 $f_n\rightharpoonup f$ 弱收敛于 $\sigma(L^p,L^{p'})$.
存在弱收敛于某个极限 $\tilde f$ 的子序列 $(f_{n_k})$.由前面的讨论我们知道
$f=\tilde f$ 几乎处处.``极限的唯一性''蕴含整个序列 $(f_n)$ 弱收敛于 $f$.

3. 第一种方法. 写出
\begin{equation}
\norm{f_n-f}_q\le\norm{f_n-T_kf_n}_q+\norm{T_kf_n-T_kf}_q+\norm{T_kf-f}_q.
\tag{S3}
\label{eq:S4.16c}
\end{equation}
注意对每个 $k>0$, 有
\begin{equation*}
\int_\Omega\abs{f_n-T_kf_n}^q\le\int_{[\abs{f_n}\ge k]}\abs{f_n}^q.
\end{equation*}
另一方面, $\int_\Omega\abs{f_n}^p\le C^p$, 从而
\begin{equation*}
k^{p-q}\int_{[\abs{f_n}\ge k]}\abs{f_n}^q\le C^p.
\end{equation*}
由此推出对一切 $n$,
\begin{equation}
\norm{f_n-T_kf_n}_q\le\left(\frac{C^p}{k^{p-q}}\right)^{1/q}.
\tag{S4}
\label{eq:S4.16d}
\end{equation}
取极限 ($n\to\infty$), 借助 Fatou 引理我们得到
$\norm{f-T_kf}_q\le\left(\dfrac{C^p}{k^{p-q}}\right)^{1/q}$.给定
$\varepsilon>0$, 固定足够大的 $k$ 使得 $(C^p/k^{p-q})^{1/q}<\varepsilon$.
显然 (由控制收敛定理) 当 $n\to\infty$ 时 $\norm{T_kf_n-T_kf}_q\to0$, 从而
存在某个整数 $N$ 使得对一切 $n\ge N$, $\norm{T_kf_n-T_kf}_q<\varepsilon$.
结合 \eqref{eq:S4.16c} 与 \eqref{eq:S4.16d}, 可见对一切 $n\ge N$,
$\norm{f_n-f}_q<3\varepsilon$.

第二种方法. 由 Egorov 定理, 给定 $\delta>0$, 存在某个 $A\subset\Omega$ 使得
$\abs{A}<\delta$ 且 $f_n\to f$ 在 $\Omega\setminus A$ 上一致收敛.写出
\begin{align*}
\int_\Omega\abs{f_n-f}^q&=\int_{\Omega\setminus A}+\int_A
\le\norm{f_n-f}_{L^\infty(\Omega\setminus A)}^q\abs{\Omega}+\norm{f_n-f}_p^q\abs{A}^{1-(q/p)}\\
&\le\norm{f_n-f}_{L^\infty(\Omega\setminus A)}^q\abs{\Omega}+(2C)^q\delta^{1-(q/p)},
\end{align*}
这导出对一切 $\delta>0$, $\limsup_{n\to\infty}\int_\Omega\abs{f_n-f}^q\le(2C)^q\delta^{1-(q/p)}$.

\solutionhead{4.17}

1. 由齐次性, 只须验证
\begin{equation*}
\sup_{t\in[-1,+1]}\left\{\frac{\abs{\abs{t+1}^p-\abs{t}^p-1}}{\abs{t}^{p-1}+\abs{t}}\right\}<\infty.
\end{equation*}

\solutionhead{4.18}

1. 首先容易验证对一切 $a,b\in(0,1)$, $\int_a^bu_n(t)dt\to(b-a)\overline f$.
只要 $1<p\le\infty$, 这就蕴含 $u_n\rightharpoonup\overline f$ 弱收敛于
$\sigma(L^p,L^{p'})$ (因为 $p'<\infty$, 阶梯函数在 $L^{p'}$ 中稠密).
当 $p=1$ 即 $f\in L^1_{\text{loc}}(\R)$ 时, 存在一个 $T$-周期函数
$g\in L^\infty(\R)$ 使得 $\frac1T\int_0^T\abs{f-g}<\varepsilon$ (其中
$\varepsilon>0$ 任意固定).令 $v_n(x)=g(nx)$, $x\in(0,1)$, 并设 $\phi\in L^\infty(0,1)$.
我们有
\begin{equation*}
\abs{\int u_n\phi-\overline f\int\phi}\le3\varepsilon\norm{\phi}_\infty+\abs{\int v_n\phi-\overline g\int\phi},
\end{equation*}
从而对一切 $\varepsilon>0$, $\limsup_{n\to\infty}\abs{\int u_n\phi-\overline f\int\phi}\le3\varepsilon\norm{\phi}_\infty$.
由此推出 $u_n\rightharpoonup\overline f$ 弱收敛于 $\sigma(L^1,L^\infty)$.

2. $\lim_{n\to\infty}\norm{u_n-\overline f}_p=\left(\frac1T\int_0^T\abs{f-\overline f}^p\right)^{1/p}$.

3. (i) $u_n\overset{*}{\rightharpoonup}0$ 关于 $\sigma(L^\infty,L^1)$.

(ii) $u_n\overset{*}{\rightharpoonup}\frac12(\alpha+\beta)$ 关于 $\sigma(L^\infty,L^1)$.

\solutionhead{4.20}

1. 设 $(u_n)$ 为 $L^p(\Omega)$ 中的序列使得在 $L^p(\Omega)$ 中 $u_n\to u$
强收敛.存在子序列使得 $u_{n_k}(x)\to u(x)$ 几乎处处, 且对一切 $k$,
$\abs{u_{n_k}}\le v$, 其中 $v\in L^p(\Omega)$ (见定理  \ref{thm:4.9}).由
控制收敛定理推出在 $L^q(\Omega)$ 中 $Au_{n_k}\to Au$ 强收敛.``极限的唯一性''
蕴含整个序列 $(Au_n)$ 在 $L^q(\Omega)$ 中强收敛于 $Au$.

2. 考虑习题 4.18 第 3(ii) 问中定义的序列 $(u_n)$.注意 $u_n\rightharpoonup\frac12(\alpha+\beta)$,
而 $Au_n\rightharpoonup\frac12(a(\alpha)+a(\beta))$.由此推出对一切 $\alpha,\beta\in\R$,
$a\left(\frac{\alpha+\beta}{2}\right)=\frac12(a(\alpha)+a(\beta))$, 因此 $a$
必为仿射函数.

\solutionhead{4.21}

1. 验证对每个有界区间 $I$, $\int_I u_n(t)dt\to0$.然后利用 (具有紧支集的)
阶梯函数在 $L^{p'}(\R)$ 中的稠密性.

2. 我们再次断言对每个有界区间 $I$, $\int_I u_n(t)dt\to0$.事实上, 给定
$\varepsilon>0$, 固定 $\delta>0$ 使得 $\delta(\norm{u_0}_\infty+\abs{I})<\varepsilon$.
令 $E=[\abs{u_0}>\delta]$, 写出
\begin{equation*}
\int_I u_n(t)dt=\int_{I+n}u_0(t)dt=\int_{(I+n)\cap E}u_0+\int_{(I+n)\cap E^c}u_0.
\end{equation*}
选取充分大的 $N$ 使得对一切 $n>N$, $\abs{(I+n)\cap E}<\delta$.我们得到
对一切 $n\ge N$, $\abs{\int_I u_n(t)dt}\le\delta\norm{u_0}_\infty+\delta\abs{I}<\varepsilon$.
然后利用 (具有紧支集的) 阶梯函数在 $L^1(\R)$ 中的稠密性.

3. 反设 $u_{n_k}\rightharpoonup u$ 弱收敛于 $\sigma(L^1,L^\infty)$.考虑由
$f=\sum_i(-1)^i\chi_{(-n_i,-n_{i+1})}$ 定义的函数 $f\in L^\infty(\R)$.注意
$\int u_{n_k}f=(-1)^k$ 不收敛.

\solutionhead{4.22}

1. 为了证明 (B) $\Rightarrow$ (A), 利用由形如 $\chi_E$ (其中 $E$ 可测且
$\abs{E}<\infty$) 的函数张成的向量空间在 $L^{p'}(\Omega)$ 中稠密的事实
(只要 $p'<\infty$).

2. 利用由形如 $\chi_E$ (其中 $E\subset\Omega$ 且 $E$ 可测) 的函数张成的向量
空间在 $L^\infty(\Omega)$ 中稠密的事实.

4. 给定 $\varepsilon>0$, 固定某个满足 $\abs{\omega}<\infty$ 的可测子集
$\omega\subset\Omega$ 使得
\begin{equation}
\int_{\omega^c}f<\varepsilon.
\tag{S1}
\label{eq:S4.22}
\end{equation}
我们有 $\int_{\omega^c}f_n=\int_{\omega^c}f+\left(\int_\omega f-\int_\omega f_n\right)+\left(\int_\Omega f_n-\int_\Omega f\right)$,
从而由 (b) 与 (c), $\int_{\omega^c}f_n=\int_{\omega^c}f+o(1)$.另一方面,
$\int_F f_n=\int_{F\cap\omega}f_n+\int_{F\cap\omega^c}f_n=\int_{F\cap\omega}f+\int_{F\cap\omega^c}f_n+o(1)$,
从而 $\int_F f_n-\int_F f=\int_{F\cap\omega^c}(f_n-f)+o(1)$.结合 \eqref{eq:S4.22},
我们得到 $\abs{\int_F f_n-\int_F f}\le2\varepsilon+o(1)$.由此推出
$\int_F f_n\to\int_F f$.最后, 我们利用由函数 $\chi_F$ ($F\subset\Omega$ 可测且
$\abs{F}\le\infty$) 张成的向量空间在 $L^\infty(\Omega)$ 中稠密的事实.

\solutionhead{4.23}

1. 设 $(u_n)$ 为 $C$ 中的序列使得在 $L^p(\Omega)$ 中 $u_n\to u$ 强收敛.存在
子序列 $(u_{n_k})$ 使得 $u_{n_k}\to u$ 几乎处处.从而 $u\ge f$ 几乎处处.

2. 假设 $u\in L^\infty(\Omega)$ 满足对满足 $f\phi\in L^1(\Omega)$ 且 $\phi\ge0$
的一切 $\phi\in L^1(\Omega)$, $\int_\Omega u\phi\ge\int_\Omega f\phi$.我们
断言 $u\ge f$ 几乎处处.事实上, 写出 $\Omega=\bigcup_n\Omega_n$, 其中
$\abs{\Omega_n}<\infty$, 并令 $\Omega'_n=\Omega_n\cap[\abs{f}<n]$, 从而
$\bigcup_n\Omega'_n=\Omega$.令 $A=[u<f]$.选取 $\phi=\chi_{A\cap\Omega'_n}$,
我们发现 $\int_{A\cap\Omega'_n}\abs{f-u}\le0$, 从而对一切 $n$,
$\abs{A\cap\Omega'_n}=0$.由此推出 $\abs{A}=0$.

3. 注意若固定满足 $f\phi\in L^1(\Omega)$ 的 $\phi\in L^1(\Omega)$, 则集合
$\{u\in L^\infty(\Omega);\ \int_\Omega u\phi\ge\int_\Omega f\phi\}$ 关于拓扑
$\sigma(L^\infty,L^1)$ 是闭的.

\solutionhead{4.24}

1. 对每个 $\phi\in L^1(\R^N)$, 我们有
\begin{equation*}
\int v_n\phi=\int u\zeta_n(\check\rho_n\star\phi)
=\int u\zeta_n(\check\rho_n\star\phi-\phi)+\int u\zeta_n\phi,
\end{equation*}
从而
\begin{equation*}
\abs{\int v_n\phi-\int v\phi}\le\norm{u}_\infty\norm{\check\rho_n\star\phi-\phi}_1
+\norm{u}_\infty\norm{(\zeta_n-\zeta)\phi}_1.
\end{equation*}
右边的第一项由定理  \ref{thm:4.22} 趋于零, 而右边的第二项由控制收敛定理趋于零.

2. 设 $B=B(x_0,R)$, 设 $\chi$ 表示 $B(x_0,R+1)$ 的特征函数.令
$\tilde v_n=\rho_n\star(\zeta_n\chi u)$.注意在 $B(x_0,R)$ 上 $\tilde v_n=v_n$,
因为 $\operatorname{supp}(\tilde v_n-v_n)\subset B(0,1/n)+B(x_0,R+1)^c$.另一方面,
我们有
\begin{align*}
\int_B\abs{v_n-v}&=\int_B\abs{\tilde v_n-\chi v}\le\int_{\R^N}\abs{\tilde v_n-\chi v}\\
&\le\int_{\R^N}\abs{\rho_n\star(\zeta_n-\zeta)\chi u}+\int_{\R^N}\abs{(\rho_n\star\chi v)-\chi v}\\
&\le\int_{\R^N}\abs{(\zeta_n-\zeta)\chi u}+\int_{\R^N}\abs{(\rho_n\star\chi v)-\chi v}\to0.
\end{align*}

\solutionhead{4.25}

1. 用 $\bar u$ 记 $u$ 在 $\Omega$ 外延拓为 $0$ 所得的函数.令
$\Omega_n=\{x\in\Omega;\ \dist(x,\partial\Omega)>2/n\ \text{且}\ \abs{x}<n\}$.
设 $\zeta_n$ (相应地 $\zeta$) 表示 $\Omega_n$ (相应地 $\Omega$) 的特征函数,
从而在 $\R^N$ 上 $\zeta_n\to\zeta$.令 $v_n=\rho_n\star(\zeta_n\bar u)$.我们
知道 $v_n\in C_c^\infty(\Omega)$, 且对每个球 $B$, $\int_B\abs{v_n-u}\to0$ (由习题
4.24).因此对每个球 $B$, 存在 (依赖于 $B$ 的) 子序列在 $B$ 上几乎处处收敛于
$u$.由对角线法我们可以构造在 $\R^N$ 上几乎处处收敛于 $u$ 的子序列 $(v_{n_k})$.

\solutionhead{4.26}

1. 假设 $A<\infty$.证明 $f\in L^1(\Omega)$ 且 $\norm{f}_1\le A$.对一切
$\phi\in C_c(\Omega)$, 我们有 $\abs{\int_\Omega f\phi}\le A\norm{\phi}_\infty$.
设 $K\subset\Omega$ 为任何紧子集, 设 $\psi\in C_c(\Omega)$ 为满足在 $\Omega$
上 $0\le\psi\le1$ 且在 $K$ 上 $\psi=1$ 的函数.设 $u$ 为 $L^\infty(\Omega)$
中的任何函数.利用习题 4.25 我们可以构造 $C_c(\Omega)$ 中的序列 $(u_n)$ 使得
$\norm{u_n}_\infty\le\norm{u}_\infty$ 且在 $\Omega$ 上 $u_n\to u$ 几乎处处.
我们有 $\abs{\int_\Omega f\psi u_n}\le A\norm{u}_\infty$.由控制收敛定理令
$n\to\infty$, 得到对一切 $u\in L^\infty(\Omega)$, $\abs{\int_\Omega f\psi u}\le A\norm{u}_\infty$.
选取 $u=\operatorname{sign}(f)$, 发现对 $\Omega$ 的每个紧子集 $K$,
$\int_K\abs{f}\le A$.由此推出 $f\in L^1(\Omega)$ 且 $\norm{f}_1\le A$.

2. 假设 $B<\infty$.对满足 $\phi\ge0$ 的一切 $\phi\in C_c(\Omega)$, 有
$\int_\Omega f\phi\le B\norm{\phi}_\infty$.使用与第 1 问相同的方法, 得到
对满足 $u\ge0$ 的一切 $u\in L^\infty(\Omega)$, $\int_\Omega f\psi u\le B\norm{u}_\infty$.
选取 $u=\chi_{[f>0]}$ 发现 $\int_K f^+\le B$.

\solutionhead{4.27}

先考察一个抽象框架.设 $E$ 为向量空间, 设 $f,g$ 为 $E$ 上的两个线性泛函,
满足 $f\not\equiv0$.假设
\begin{equation*}
[\phi\in E\ \text{且}\ f(\phi)>0]\implies[g(\phi)\ge0].
\end{equation*}
我们断言存在常数 $\lambda\ge0$ 使得 $g=\lambda f$.事实上, 固定任何满足
$f(\phi_0)=1$ 的 $\phi_0\in E$.对每个 $\phi\in E$ 与每个 $\varepsilon>0$,
有 $f(\phi-f(\phi)\phi_0+\varepsilon\phi_0)=\varepsilon>0$, 从而
$g(\phi-f(\phi)\phi_0+\varepsilon\phi_0)\ge0$.由此推出对一切 $\phi\in E$,
$g(\phi)\ge\lambda f(\phi)$, 从而 $g=\lambda f$, 其中 $\lambda=g(\phi_0)\ge0$.

\runin{应用.} $E=C_c^\infty(\Omega)$, $f(\phi)=\int_\Omega u\phi$, 以及
$g(\phi)=\int_\Omega v\phi$.

\solutionhead{4.30}

1 与 2. 注意 $\dfrac{1}{p'}+\dfrac{1}{q'}+\dfrac1r=1$ 以及 $(1-\alpha)r=p$,
$(1-\beta)r=q$.对几乎所有的 $x\in\R^N$, 写出
\begin{equation*}
\abs{f(x-y)g(y)}=\phi_1(y)\phi_2(y)\phi_3(y),
\end{equation*}
其中 $\phi_1(y)=\abs{f(x-y)}^\alpha$, $\phi_2(y)=\abs{g(y)}^\beta$, 以及
$\phi_3(y)=\abs{f(x-y)}^{1-\alpha}\abs{g(y)}^{1-\beta}$.显然
$\phi_1\in L^{q'}(\R^N)$ 且 $\phi_2\in L^{p'}(\R^N)$.另一方面,
$\abs{\phi_3(y)}^r=\abs{f(x-y)}^p\abs{g(y)}^q$.由定理  \ref{thm:4.15} 推出:
对几乎所有的 $x\in\R^N$, 函数 $y\mapsto\abs{\phi_3(y)}^r$ 可积.由 Hölder
不等式 (见习题 4.4) 推出: 对几乎所有的 $x\in\R^N$, 函数 $y\mapsto\abs{f(x-y)g(y)}$
可积, 且
\begin{equation*}
\int_{\R^N}\abs{f(x-y)}\abs{g(y)}dy\le\norm{f}_p^\alpha\norm{g}_q^\beta
\left(\int_{\R^N}\abs{f(x-y)}^p\abs{g(y)}^qdy\right)^{1/r}.
\end{equation*}
因此
\begin{equation*}
\abs{(f\star g)(x)}^r\le\norm{f}_p^{\alpha r}\norm{g}_q^{\beta r}
\int_{\R^N}\abs{f(x-y)}^p\abs{g(y)}^qdy,
\end{equation*}
从而
\begin{equation*}
\int_{\R^N}\abs{(f\star g)(x)}^rdx\le\norm{f}_p^{\alpha r}\norm{g}_q^{\beta r}
\norm{f}_p^p\norm{g}_q^q=\norm{f}_p^r\norm{g}_q^r.
\end{equation*}

3. 若 $1<p<\infty$ 且 $1<q<\infty$, 存在 $C_c(\R^N)$ 中的序列 $(f_n)$ 与
$(g_n)$ 使得在 $L^p(\R^N)$ 中 $f_n\to f$, 在 $L^q(\R^N)$ 中 $g_n\to g$.
则 $f_n\star g_n\in C_c(\R^N)$, 此外 $\norm{(f_n\star g_n)-(f\star g)}_\infty\to0$.
由此推出当 $\abs{x}\to\infty$ 时 $(f\star g)(x)\to0$.

\solutionhead{4.34}

给定任何 $\varepsilon>0$, 存在 $\mathcal{F}$ 在 $L^p(\R^N)$ 中被半径为
$\varepsilon$ 的球的有限覆盖, 设 $\mathcal{F}\subset\bigcup_{i=1}^kB(f_i,\varepsilon)$.

2. 对每个 $i$, 存在某个 $\delta_i>0$ 使得对一切满足 $\abs{h}<\delta_i$ 的
$h\in\R^N$, $\norm{\tau_hf_i-f_i}_{L^p(\R^N)}<\varepsilon$ (见引理  \ref{lem:4.3}).
令 $\delta=\min_{1\le i\le k}\delta_i$.容易验证对一切 $f\in\mathcal{F}$ 及一切
满足 $\abs{h}<\delta$ 的 $h\in\R^N$, $\norm{\tau_hf-f}_p<3\varepsilon$.

3. 对每个 $i$, 存在某个有界开集 $\Omega_i\subset\R^N$ 使得
$\norm{f_i}_{L^p(\R^N\setminus\Omega_i)}<\varepsilon$.令
$\Omega=\bigcup_{i=1}^k\Omega_i$, 验证对一切 $f\in\mathcal{F}$,
$\norm{f}_{L^p(\R^N\setminus\Omega)}<2\varepsilon$.

\solutionhead{4.37}

1. 写出
\begin{align*}
\int_I u_n(x)\phi(x)dx&=\int_{-n}^{+n}f(t)\left[\phi\left(\frac{t}{n}\right)-\phi(0)\right]dt
+\phi(0)\int_{-n}^{+n}f(t)dt\\
&=A_n+B_n;
\end{align*}
由 Lebesgue 定理 $A_n\to0$, 且由于 $\int_{-\infty}^{+\infty}f(t)dt=0$,
$B_n\to0$.

2. 注意对所有 $\delta>0$,
\begin{equation*}
\int_0^\delta\abs{u_n(x)}dx=\int_0^{n\delta}\abs{f(t)}dt\to\int_0^\infty\abs{f(t)}dt>0.
\end{equation*}

3. 反设成立.我们将对一切 $\phi\in C([-1,+1])$ 有 $\int_I u\phi=0$, 从而
(由推论  \ref{cor:4.24}) $u\equiv0$.另一方面, 若选取 $\phi=\chi_{(0,1)}$,
我们得到 $\int_I u_n\phi=\int_0^nf(t)dt\to\int_0^{+\infty}f(t)dt>0$.矛盾.

\solutionhead{4.38}

2. 验证对一切 $\phi\in C^1([0,1])$, 当 $n\to\infty$ 时
$\int_I u_n\phi=\int_I\phi+O(1/n)$.然后利用 $\norm{u_n}_1$ 有界以及
$C^1([0,1])$ 在 $C([0,1])$ 中稠密的事实.

3. 序列 $(u_n)$ 不能是等度可积的, 因为 $\abs{\operatorname{supp}u_n}\to0$ 且
$1=\int_I u_n=\int_{\operatorname{supp}u_n}\abs{u_n}$.

4. 若 $u_{n_k}\rightharpoonup u$ 弱收敛于 $\sigma(L^1,L^\infty)$, 则由第
2 问与推论  \ref{cor:4.24}, 我们会有 $u\equiv1$.选取进一步的子序列 $(u_{n'_k})$
使得 $\sum_k\abs{\operatorname{supp}u_{n'_k}}<1$.设 $\phi=\chi_A$, 其中
$A=I\setminus\left(\bigcup_k\operatorname{supp}u_{n'_k}\right)$, 从而
$\abs{A}>0$.我们对一切 $k$ 有 $\int_I u_{n'_k}\phi=0$, 从而
$0=\int_I\phi=\abs{A}$.矛盾.

5. 考虑满足 $\sum_k\abs{\operatorname{supp}u_{n_k}}<\infty$ 的子序列 $(u_{n_k})$.
令 $B_k=\bigcup_{j\ge k}(\operatorname{supp}u_{n_j})$ 且 $B=\bigcap_kB_k$.
显然当 $k\to\infty$ 时 $\abs{B_k}\to0$, 从而 $\abs{B}=0$.若 $x\notin B$,
则存在某个 $k_0$ 使得对一切 $k\ge k_0$, $u_{n_k}(x)=0$.
\solutionhead{5.1}

1. 对 $a=u+v$ 与 $b=v$ 应用平行四边形法则给出 $(u,2v)=2(u,v)$.

2. 计算 (i) $-$ (ii) $+$ (iii).

3. 注意由 $(\ ,\ )$ 的定义, 映射 $\lambda\in\R\mapsto(\lambda u,v)$ 是连续的.

\solutionhead{5.2}

设 $A$ 为满足 $0<\abs{A}<\abs{\Omega}$ 的可测集, 选取满足 $A\cap B=\emptyset$ 且
$0<\abs{B}<\abs{\Omega}$ 的可测集 $B$.令 $u=\chi_A$ 与 $v=\chi_B$.

先设 $1\le p<\infty$.我们有
\begin{equation*}
\norm{u+v}_p^p=\norm{u-v}_p^p=\abs{A}+\abs{B},
\end{equation*}
从而 $\norm{u+v}_p^2+\norm{u-v}_p^2=2(\abs{A}+\abs{B})^{2/p}$.另一方面,
我们有 $2(\norm{u}_p^2+\norm{v}_p^2)=2(\abs{A}^{2/p}+\abs{B}^{2/p})$.最后,
注意对一切 $\alpha,\beta>0$, 当 $p<2$ 时
$(\alpha+\beta)^{2/p}>\alpha^{2/p}+\beta^{2/p}$, 当 $p>2$ 时
$(\alpha+\beta)^{2/p}<\alpha^{2/p}+\beta^{2/p}$.用相同的函数 $u$ 与 $v$
考察 $p=\infty$ 的情形.

\solutionhead{5.3}

验证
\begin{equation}
2(t_nu_n-t_mu_m,u_n-u_m)=(t_n+t_m)\abs{u_n-u_m}^2+(t_n-t_m)(\abs{u_n}^2-\abs{u_m}^2),
\tag{S1}
\label{eq:S5.3}
\end{equation}
这蕴含对一切 $m,n$, $(t_n-t_m)(\abs{u_n}^2-\abs{u_m}^2)\le0$.

1. 设 $n>m$, 从而 $t_n\ge t_m$, 进而 $\abs{u_n}\le\abs{u_m}$ (注意若
$t_n=t_m$, 由 \eqref{eq:S5.3} 有 $u_n=u_m$).另一方面, 对 $n>m$, 我们有
\begin{equation*}
(t_n+t_m)\abs{u_n-u_m}^2\le(t_n-t_m)(\abs{u_m}^2-\abs{u_n}^2)\le t_n(\abs{u_m}^2-\abs{u_n}^2),
\end{equation*}
从而 $\abs{u_n-u_m}^2\le\abs{u_m}^2-\abs{u_n}^2$.由此推出当 $n\uparrow\infty$
时 $\abs{u_n}\downarrow\ell$, 且 $(u_n)$ 是 Cauchy 序列.

2. 设 $n>m$, 从而 $t_m\ge t_n$ 且 $\abs{u_m}\le\abs{u_n}$.对 $n>m$, 有
\begin{equation*}
(t_n+t_m)\abs{u_n-u_m}^2\le(t_m-t_n)(\abs{u_n}^2-\abs{u_m}^2)\le t_m(\abs{u_n}^2-\abs{u_m}^2),
\end{equation*}
从而 $\abs{u_n-u_m}^2\le\abs{u_n}^2-\abs{u_m}^2$.我们有下述抉择:

(i) 或者当 $n\uparrow\infty$ 时 $\abs{u_n}\uparrow\infty$,

(ii) 或者当 $n\uparrow\infty$ 时 $\abs{u_n}\uparrow\ell<\infty$, 此时 $(u_n)$ 是 Cauchy 序列.

另一方面, 令 $v_n=t_nu_n$ 与 $s_n=1/t_n$, 得到
\begin{equation*}
(s_nv_n-s_mv_m,v_n-v_m)\le0,
\end{equation*}
从而由第 1 问 $(v_n)$ 收敛于一个极限.由此推出若 $t_n\to t>0$, 则 $(u_n)$
也收敛于一个极限.最后若 $t_n\to0$, 两种情形 (i) 与 (ii) 均可出现.例如取
$H=\R$, 情形 (i) 取 $u_n=C/t_n$, 情形 (ii) 取 $u_n=C$.

\solutionhead{5.4}

注意
\begin{equation*}
\abs{v-u}^2=\abs{v-f}^2-\abs{u-f}^2+2(f-u,v-u).
\end{equation*}

\solutionhead{5.5}

1. 设 $K=\overline{\bigcup_nK_n}$.我们断言 $u_n\to u=P_Kf$.首先注意序列
$d_n=\abs{f-u_n}=\dist(f,K_n)$ 单调递增且有上界.从而当 $n\uparrow\infty$
时 $d_n\uparrow\ell<\infty$.接着应用平行四边形法则 (取 $a=f-u_n,b=f-u_m$),
得到
\begin{equation*}
\abs{f-\frac{u_n+u_m}{2}}^2+\abs{\frac{u_n-u_m}{2}}^2
=\frac12(\abs{f-u_n}^2+\abs{f-u_m}^2).
\end{equation*}
由此推出若 $m\ge n$, 则 $\abs{u_n-u_m}^2\le2(d_m^2-d_n^2)$.因此 $(u_n)$
收敛于极限 $u$, 且显然 $u\in K$.另一方面, 对一切 $v\in K_n$ 有
$\abs{f-u_n}\le\abs{f-v}$, 特别地对一切 $v\in K$ 有 $\abs{f-u_n}\le\abs{f-v}$.
取极限得到对一切 $v\in K$, $\abs{f-u}\le\abs{f-v}$.

2. 显然 $K=\overline{\bigcup_nK_n}$ 是凸的 (为什么?).我们断言
$u_n\to u=P_Kf$.首先注意序列 $d_n=\abs{f-u_n}=\dist(f,K_n)$ 单调递减,
从而 $d_n\to\ell$.接着 (用上面的相同方法), 若 $m\ge n$, 有
$\abs{u_n-u_m}^2\le2(d_n^2-d_m^2)$.因此 $(u_n)$ 收敛于极限 $u$, 且显然
$u\in K$.最后注意只要 $m\ge n$, 对一切 $v\in K_n$ 就有 $\abs{f-u_m}\le\abs{f-v}$.
取极限 ($m\to\infty$) 导出对一切 $v\in\bigcup_nK_n$ 有 $\abs{f-u}\le\abs{f-v}$,
由稠密性对一切 $v\in K$ 成立.

序列 $(\alpha_n)$ 单调递减, 从而收敛于极限 $\alpha$.我们断言
$\alpha=\inf_K\varphi$.首先, 显然 $\inf_K\varphi\le\alpha_n$, 从而
$\inf_K\varphi\le\alpha$.另一方面, 设 $u$ 为 $K$ 中的任意元素, 设
$u_n=P_{K_n}u$.在不等式 $\alpha_n\le\varphi(u_n)$ 中取极限, 得到
$\alpha\le\varphi(u)$ (因为 $u_n\to u$).由此推出 $\alpha\le\inf_K\varphi$.

\solutionhead{5.6}

1. 例如考虑 $\norm{u}\ge1$ 且 $\norm{v}\le1$ 的情形.我们有
\begin{equation*}
\norm{Tu-Tv}=\norm{\frac{u}{\norm{u}}-v}=\frac{\norm{(u-v)+(v-v\norm{u})}}{\norm{u}}
\le\norm{u-v}+\norm{u}-1\le2\norm{u-v},
\end{equation*}
因为 $\norm{u}\le\norm{u-v}+\norm{v}\le\norm{u-v}+1$.

2. 设 $u=(1,0)$ 与 $v=(1,\alpha)$.则有
$\norm{Tu-Tv}=2\abs{\alpha}/(1+\abs{\alpha})$, 而 $\norm{u-v}=\abs{\alpha}$.
选取 $\alpha\ne0$ 且任意小即得出结论.

3. $T$ 与 $P_{B_E}$ 重合.只须验证若 $\norm{u}\ge1$, 则对一切 $v\in B_E$,
\begin{equation*}
\left(u-\frac{u}{\norm{u}},v-\frac{u}{\norm{u}}\right)\le0.
\end{equation*}

\solutionhead{5.10}

(i) $\Rightarrow$ (ii). 写出对一切 $t\in(0,1)$ 与一切 $v\in K$,
$F(u)\le F((1-t)u+tv)$, 这蕴含
\begin{equation*}
\frac1t[F(u+t(v-u))-F(u)]\ge0.
\end{equation*}
令 $t\to0$ 即得 (ii).

(ii) $\Rightarrow$ (i). 我们断言对一切 $u,v\in H$,
$F(v)-F(u)\ge(F'(u),v-u)$.事实上, 函数 $t\in\R\mapsto\phi(t)=F(u+t(v-u))$
属于 $C^1$ 类且是凸的.因此 $\phi(1)-\phi(0)\ge\phi'(0)$.

\solutionhead{5.12}

$T$ 满当且仅当 $E$ 完备.

1. 通过令对一切 $u,v\in E$,
\begin{equation*}
((T(u),T(v)))=(u,v),
\end{equation*}
把 $E$ 的标量积转移到 $R(T)$ 上.注意对一切 $f,g\in R(T)$,
$\abs{((f,g))}\le\norm{f}_{E^\star}\norm{g}_{E^\star}$.标量积 $(( , ))$
可以通过连续性与稠密性延拓到 $\overline{R(T)}$ 上, 它现在配有 Hilbert
空间的结构.

2. 固定任何 $f\in E^\star$.映射 $g\in R(T)\mapsto\dual{f}{T^{-1}(g)}$ 是
$R(T)$ 上的连续线性泛函.它可以 (通过连续性) 延拓到 $\overline{R(T)}$.在
$\overline{R(T)}$ 中应用 Riesz--Fréchet 表示定理, 得到某个元素
$h\in\overline{R(T)}$ 使得对一切 $g\in R(T)$, $((h,g))=\dual{f}{T^{-1}(g)}$.
从而对一切 $v\in E$, $((h,T(v)))=\dual{f}{v}$.另一方面, 对一切
$h\in\overline{R(T)}$ 及一切 $v\in E$, 有 $((h,Tv))=\dual{h}{v}$ (当
$h\in R(T)$ 时这显然).由此推出 $f=h$, 从而 $f\in R(T)$, 即 $R(T)=E^\star$.

3. 我们构造了等距映射 $T:E\to E^\star$, 且 $R(T)$ 在 $E^\star$ 中稠密.由于
$E^\star$ 完备, 我们得出 (在同构意义下) $E^\star$ 是 $E$ 的完备化.

\solutionhead{5.13}

1. 我们断言平行四边形法则成立.事实上, 设 $f\in F(u)$ 与 $g\in F(v)$.则
$f\pm g\in F(u\pm v)$, 从而我们有
\begin{equation*}
\dual{f+g}{u+v}=\norm{u+v}^2\quad\text{与}\quad\dual{f-g}{u-v}=\norm{u-v}^2.
\end{equation*}
把这些关系式相加得出 $2(\norm{u}^2+\norm{v}^2)=\norm{u+v}^2+\norm{u-v}^2$.

2. 设 $T:E\to E^\star$ 为习题 5.12 中引入的映射.我们断言 $F(u)=\{T(u)\}$.
显然 $T(u)\in F(u)$.另一方面, 我们知道 $E^\star$ 关于对偶范数
$\norm{\cdot}_{E^\star}$ 是 Hilbert 空间.特别地, $E^\star$ 严格凸, 从而
(见习题 1.1) $F(u)$ 退化为单点.

\solutionhead{5.14}

凸性不等式 $a(tu+(1-t)v,tu+(1-t)v)\le ta(u,u)+(1-t)a(v,v)$ 等价于
$t(1-t)a(u-v,u-v)\ge0$.

考虑由 $a(u,v)=(Au,v)\ \forall u,v\in H$ 定义的算子 $A\in\LL(H)$.则
$F'(u)=Au+A^\star u$, 因为我们有
\begin{equation*}
F(u+h)-F(u)=(Au+A^\star u,h)+a(h,h).
\end{equation*}

\solutionhead{5.15}

首先通过连续性把 $S$ 延拓为算子 $\overline S:G\to F$.接着令
$T=\overline S\circ P_G$, 其中 $P_G$ 表示由 $H$ 到 $G$ 上的投影算子.

\solutionhead{5.18}

(ii) $\Rightarrow$ (i). 假设 (ii) 蕴含 $T$ 是单射且 $R(T)$ 闭.从而 $R(T)$
有补空间 (因为 $H$ 是 Hilbert 空间).由定理  \ref{thm:2.13} 我们推出 $T$ 有左逆.

(i) $\Rightarrow$ (ii). 假设 (i) 蕴含 $T$ 是单射且 $R(T)$ 闭.然后使用定理~
\ref{thm:2.21}.

\solutionhead{5.19}

注意
\begin{equation*}
\limsup_{n\to\infty}\abs{u_n-u}^2=\limsup_{n\to\infty}(\abs{u_n}^2-2(u_n,u)+\abs{u}^2)\le0.
\end{equation*}

\solutionhead{5.20}

1. 若 $u\in N(S)$, 则对一切 $v\in H$, $(Sv,v-u)\ge0$; 用 $tv$ 替换 $v$,
可见对一切 $v\in H$, $(Sv,u)=0$.反之, 若 $u\in R(S)^\perp$, 则对一切
$v\in H$, $(Sv-Su,v)\ge0$; 用 $tv$ 替换 $v$, 可见对一切 $v\in H$, $(Su,v)=0$
(另见问题 16).

2. 应用推论  \ref{cor:5.8} (Lax--Milgram).

3. 方法 (a). 令 $u_t=(I+tS)^{-1}f$.

若 $f\in N(S)$, 则对一切 $t>0$, $u_t=f$.

若 $f\in R(S)$, 写出 $f=Sv$, 从而 $u_t+S(tu_t-v)=0$.由此推出
$(u_t,tu_t-v)\le0$, 从而 $\abs{u_t}\le(1/t)\abs{v}$.因此当 $t\to\infty$
时 $u_t\to0$.由稠密性, 仍可证明对每个 $f\in\overline{R(S)}$, 当 $t\to\infty$
时 $u_t\to0$.

在一般情形 $f\in H$ 下, 把 $f$ 写成 $f=f_1+f_2$, 其中 $f_1=P_{N(S)}f$ 且
$f_2=P_{\overline{R(S)}}f$.

方法 (b). 我们有 $u_t+tSu_t=f$, 从而 $\abs{u_t}\le\abs{f}$.抽取子序列
$t_n\to\infty$, 可以假设 $u_{t_n}\rightharpoonup u$ 弱收敛且 $Su=0$ (为什么?),
即 $u\in N(S)$.由第 1 问我们知道对一切 $v\in N(S)$, $(Su_t,v)=0$, 从而对一切
$v\in N(S)$, $(f-u_t,v)=0$.取极限, 发现对一切 $v\in N(S)$, $(f-u,v)=0$.
因此 $u=P_{N(S)}f$, ``极限的唯一性''蕴含当 $t\to\infty$ 时 $u_t\rightharpoonup u$
弱收敛.另一方面, 我们有 $(Su_t,u_t)\ge0$, 即 $(f-u_t,u_t)\ge0$, 从而
$\limsup_{t\to\infty}\abs{u_t}^2\le(f,u)=\abs{u}^2$.由此推出当 $t\to\infty$
时 $u_t\to u$ 强收敛.

\solutionhead{5.21}

1. 令 $S=I-T$ 并应用习题 5.20 第 1 问.

2. 写出 $f=u-Tu$, 注意 $\sigma_n(f)=\dfrac1n(u-T^nu)$.

3. 首先验证对一切 $f\in R(I-T)$, $\lim_{n\to\infty}\sigma_n(f)=0$.接着把一般的
$f\in H$ 分解为 $f=f_1+f_2$, 其中 $f_1\in N(I-T)$ 且
$f_2\in N(I-T)^\perp=\overline{R(I-T)}$.则我们有
$\sigma_n(f)=\sigma_n(f_1)+\sigma_n(f_2)=f_1+\sigma_n(f_2)$.

4. 依次对 $u,Su,S^2u,\dots,S^iu,\dots$ 应用不等式 (1), 并把所得的不等式相加.
注意对一切 $i=0,1,\dots,n$,
\begin{equation*}
\abs{S^nu-S^{n+1}u}\le\abs{S^iu-S^{i+1}u}.
\end{equation*}

5. 写出 $f=u-Tu=2(u-Su)$, 我们得到 $\abs{\mu_n(f)}\le2\abs{u}/\sqrt{n+1}$.

6. 使用与第 3 问相同的方法.

\solutionhead{5.25}

2. 设 $m>n$.对 $f=u_m$ 与 $v=P_Ku_n$ 应用习题 5.4, 得到
\begin{align*}
\abs{P_Ku_n-P_Ku_m}^2&\le\abs{P_Ku_n-u_m}^2-\abs{P_Ku_m-u_m}^2\\
&\le\abs{P_Ku_n-u_n}^2-\abs{P_Ku_m-u_m}^2.
\end{align*}
因此 $(P_Ku_n)$ 是 Cauchy 序列.

3. 可以假设 $u_{n_k}\rightharpoonup u$ 弱收敛.回忆对一切 $v\in K$,
$(u_n-P_Ku_n,v-P_Ku_n)\le0$.沿子序列 $n_k$ 取极限导出对一切 $v\in K$,
$(u-\ell,v-\ell)\le0$.由于 $u\in K$, 可以取 $v=u$ 并得出 $u=\ell$.再次,
``极限的唯一性''蕴含 $u_n\rightharpoonup\ell$ 弱收敛.

4. 对每个 $v\in K$, $\lim_{n\to\infty}\abs{u_n-v}^2$ 存在, 从而对每个
$v,w\in K$, $\lim_{n\to\infty}(u_n,v-w)$ 也存在.由此推出对每个 $z\in H$,
$\phi(z)=\lim_{n\to\infty}(u_n,z)$ 存在.由 Riesz--Fréchet 表示定理我们可以把
$\phi(z)$ 写成 $\phi(z)=(u,z)$, 其中 $u\in H$.最后注意对一切 $v\in K$,
$(u-\ell,v-\ell)\le0$, 从而 $\ell=P_Ku$.

5. 经平移与伸缩后总可以假设 $K=B_H$.从而 $\abs{u_n}\downarrow\alpha$.
若 $\alpha<1$, 则对充分大的 $n$, $u_n=P_Ku_n$ (我们已知 $P_Ku_n$ 强收敛).
若 $\alpha\ge1$, 则 $P_Ku_n=u_n/\abs{u_n}$ 强收敛, 从而 $u_n$ 也强收敛.

6. 回忆对一切 $v\in K$, $(u_n-P_Ku_n,v-P_Ku_n)\le0$, 从而对一切 $v\in K$,
$(u_n-\ell,v-\ell)\le\varepsilon_n$, 其中 $\varepsilon_n\to0$ ($\varepsilon_n$
依赖于 $v$).把这些不等式相加得出对一切 $v\in K$,
$(\sigma_n-\ell,v-\ell)\le\varepsilon'_n$, 其中 $\varepsilon'_n\to0$.假设
$\sigma_{n_k}\rightharpoonup\sigma$ 弱收敛, 则 $\sigma\in K$ 满足对一切
$v\in K$, $(\sigma-\ell,v-\ell)\le0$.因此 $\sigma=\ell$, ``极限的唯一性''
蕴含 $\sigma_n\rightharpoonup\ell$ 弱收敛.

\solutionhead{5.26}

3. 注意 $\sqrt{n}u_n$ 有界, 且对每个固定的 $j$, 当 $n\to\infty$ 时
$(\sqrt{n}u_n,e_j)\to0$.

\solutionhead{5.27}

设 $F$ 为由 $E_n$ 张成的向量空间的闭包.我们知道 (见定理  \ref{thm:5.9} 的证明)
对一切 $u\in H$, $\sum_{n=1}^\infty\abs{P_{E_n}u}^2=\abs{P_Fu}^2$, 从而对一切
$u\in D$, $\abs{P_Fu}=\abs{u}$.由此推出对一切 $u\in D$, $\abs{P_Fu}^2=\abs{u}^2$,
因此对一切 $u\in D$, $P_{F^\perp}u=0$.从而对一切 $u\in H$, $P_{F^\perp}u=0$,
即 $F^\perp=\{0\}$, 于是 $F=H$.

\solutionhead{5.28}

1. 由命题  \ref{prop:3.25}, $V$ 可分.考虑 $V$ 的稠密可数子集 $(v_n)$, 并像
定理  \ref{thm:5.11} 的证明那样得出结论.

\solutionhead{5.29}

2. 若 $2<p<\infty$, 使用不等式 $\norm{u}_p\le\norm{u}_\infty^{1-2/p}\norm{u}_2^{2/p}$.
注意每个无限维 Hilbert 空间 (可分或不可分) 都包含一个无限标准正交序列.

6. 在 $\Omega$ 上积分, 我们发现 $k\le M^2\abs{\Omega}$, 这给出了 $E$ 的维数的
一个上界.

\solutionhead{5.30}

1. 对每个固定的 $t\in[0,1]$, 考虑函数 $u_t(s)=p(s)\chi_{[0,t]}(s)$, 并写出
$\sum_{n=1}^\infty\abs{(u_t,e_n)}^2\le\norm{u_t}_2^2$.

2. (2) 中的等号蕴含对几乎所有的 $t\in[0,1]$, (1) 中等号成立.因此对几乎所有的
$t\in[0,1]$, $u_t=\sum_{n=1}^\infty(u_t,e_n)e_n$, 从而
$u_t\in E=\overline{\operatorname{span}\{e_n\}}$.还须验证由函数 $(u_t)$ 张成的
空间在 $L^2$ 中稠密.设 $f\in L^2$ 使得对几乎所有的 $t$, $\int_0^1fu_t=0$.由此
推出对一切 $t\in[0,1]$, $\int_0^tfp=0$, 从而 $fp=0$ 几乎处处.

\solutionhead{5.31}

容易验证当 $i\ne j$ 时 $(\phi_i,\phi_j)=0$.令 $n=2^{p+1}-1$.设 $E$ 为由
$\{\phi_0,\phi_1,\dots,\phi_n\}$ 张成的空间, $F$ 为由区间
$(\frac{i}{2^{p+1}},\frac{i+1}{2^{p+1}})$ ($0\le i\le2^{p+1}-1$) 的特征函数张成的
空间.显然 $E\subset F$, $\dim E=n+1=2^{p+1}$, 且 $\dim F=2^{p+1}$.因此 $E=F$.

\solutionhead{5.32}

2. 函数 $u=r_1r_2$ 与所有函数 $(r_i)_{i\ge0}$ 正交, 且 $u\ne0$.因此
$(r_i)_{i\ge0}$ 不是基.

3. 容易验证 $(w_n)_{n\ge0}$ 是标准正交系, 且 $w_0=r_0$, 对一切 $\ell\ge0$,
$w_{2^\ell}=r_{\ell+1}$.为了证明 $(w_n)_{n\ge0}$ 是基, 可以使用与习题 5.31
相同的论证.
\solutionhead{6.2}

3. 考虑定义在 $[0,1]$ 上的函数序列
\begin{equation*}
u_n(t)=\begin{cases}
0 & \text{若}\ 0\le t\le\frac12,\\
n\left(t-\frac12\right) & \text{若}\ \frac12<t\le\frac12+\frac1n,\\
1 & \text{若}\ \frac12+\frac1n<t\le1.
\end{cases}
\end{equation*}
注意 $T(u_n)\to f$, 但 $f\notin T(B_E)$, 因为 $f\notin C^1([0,1])$.

\solutionhead{6.3}

反证法.若结论不成立, 则存在某个 $\delta>0$ 使得对一切 $u\in E$,
$\norm{Tu}_F\ge\delta\norm{u}_E$.因此 $R(T)$ 闭.考虑由 $T_0=T$ 定义的
算子 $T_0:E\to R(T)$.显然 $T_0$ 是双射.由推论 2.6,
% 原书此处印作 "By Corollary 2.6" (印刷页 413); 第~
% \ref{chap:2}~章并无编号 2.6 的推论
% (2.6 是定理~
\ref{thm:2.6} 开映射定理), 有界逆定理是推论~
\ref{cor:2.7}. 依技能规范保留原样
% 并记录, 不擅自更正, 故此处不用 \ref 以免产生未定义引用.
$T_0^{-1}\in\LL(R(T),E)$.另一方面, $T_0\in\KK(E,R(T))$.因此 $B_E$ 紧
且 $\dim E<\infty$.

\solutionhead{6.5}

设 $T:V\to\ell^2$ 为由
\begin{equation*}
Tu=(\sqrt{\lambda_1}u_1,\sqrt{\lambda_2}u_2,\dots,\sqrt{\lambda_n}u_n,\dots)
\end{equation*}
定义的算子.显然对一切 $u\in V$, $\abs{Tu}_{\ell^2}=\norm{u}_V$, 且 $T$ 是由
$V$ 到 $\ell^2$ 上的满射.由于 $\ell^2$ 完备, 推出 $V$ 也完备.

考虑由 $J_nu=(u_1,u_2,\dots,u_n,0,0,\dots)$ 定义的算子 $J_n:V\to\ell^2$.容易
验证 $\norm{J_n-I}_{\LL(V,\ell^2)}\to0$, 从而由 $V$ 到 $\ell^2$ 中的典范
单射是紧的.

\solutionhead{6.7}

1. 假设 $T$ 是由 $E$ 弱到 $F$ 强的连续映射.则对每个 $\varepsilon>0$, 存在
$0$ 在 $E$ 弱中的邻域 $V$ 使得 $x\in V\Rightarrow\norm{Tx}<\varepsilon$.
可以假设 $V$ 具有形式
\begin{equation*}
V=\{x\in E;\ \abs{\dual{f_i}{x}}<\delta\ \forall i=1,2,\dots,n\},
\end{equation*}
其中 $f_1,\dots,f_n\in E^\star$, $\delta>0$.

设 $M=\{x\in E;\ \dual{f_i}{x}=0\ \forall i=1,\dots,n\}$, 从而对一切 $x\in M$,
$Tx=0$.另一方面, $M$ 具有有限余维 (见 \S \ref{sec:2.4} 例 2).因此
$E=M+N$, 其中 $\dim N<\infty$.由此推出 $R(T)=T(N)$ 是有限维的.

2. 注意若在 $E$ 中 $u_n\rightharpoonup u$ 弱收敛, 则在 $F$ 中
$Tu_n\rightharpoonup Tu$ 弱收敛.另一方面, $(Tu_n)$ 在 $F$ 中有紧闭包
(关于强拓扑).因此 $Tu_n\to Tu$ (参见例如习题 3.5).

6. 注意 $T^\star\in\LL(E^\star,(c_0)^\star)$.但 $(c_0)^\star=\ell^1$ (见
\S \ref{sec:11.3}).由于 $E^\star$ 自反, 由第 5 问推出 $T^\star$ 是紧的.
因此 (由定理  \ref{thm:6.4}) $T$ 紧.

\solutionhead{6.8}

1. 存在常数 $c$ 使得 $B_{R(T)}\subset cT(B_E)$, 因此 $R(T)$ 的单位球是紧的.

2. 设 $E_0$ 为 $N(T)$ 的补空间.则 $T_0=T|_{E_0}$ 是由 $E_0$ 到 $R(T)$ 上的
双射.因此 $\dim E_0=\dim R(T)<\infty$.

\solutionhead{6.9}

1. (A) $\Rightarrow$ (B):
设 $E_0$ 为 $N(T)$ 的补空间, 设 $P:E\to N(T)$ 为相应的投影算子.则
$T_0=T|_{E_0}$ 是由 $E_0$ 到 $R(T)$ 上的双射.由开映射定理存在常数 $C$
使得对一切 $u\in E_0$, $\norm{u}_E\le C\norm{Tu}_F$.由此推出对一切 $u\in E$,
\begin{equation*}
\norm{u}_E\le\norm{u-Pu}_E+\norm{Pu}_E\le C\norm{Tu}_F+\norm{Pu}_E.
\end{equation*}

(C) $\Rightarrow$ (A):
(i) 为了验证 $N(T)$ 中的单位球是紧的, 设 $(u_n)$ 为 $N(T)$ 中的序列使得
$\norm{u_n}_E\le1$.由于 $(Q(u_n))$ 在 $G$ 中有紧闭包, 可以抽取在 $G$ 中收敛的
子序列 $(Q(u_{n_k}))$.应用 (C), 可见 $(u_{n_k})$ 是 Cauchy 序列.

(ii) 引入 $N(T)$ 的补空间, 还可以额外假设 $T$ 是单射.设 $(u_n)$ 为 $E$ 中的
序列使得 $Tu_n\to f$.先验证 $(u_n)$ 有界.若不然, 令 $v_n=u_n/\norm{u_n}$.
应用 (C), 可见子序列 $(v_{n_k})$ 是 Cauchy 序列.设 $v_{n_k}\to v$, 其中
$v\in N(T)$ 且 $\norm{v}=1$; 不可能.因此 $(u_n)$ 有界, 我们可以抽取在 $G$ 中
收敛的子序列 $(Q(u_{n_k}))$.再次应用 (C), 发现 $(u_{n_k})$ 是 Cauchy 序列.

为了恢复习题 2.12 中的结果, 写出
\begin{equation*}
\norm{u}_E\le C(\norm{Tu}_F+\norm{Pu}_E)\le C(\norm{Tu}_F+\abs{Pu}),
\end{equation*}
因为 $N(T)$ 上的所有范数等价.此外,
\begin{equation*}
\abs{Pu}\le\abs{u-Pu}+\abs{u}\le C\norm{u-Pu}_E+\abs{u}\le C\norm{Tu}_F+\abs{u}.
\end{equation*}

2. 注意
\begin{equation*}
\norm{u}_E\le C(\norm{Tu}_F+\norm{Pu}_E)\le C(\norm{(T+S)u}_F+\norm{Pu}_E+\norm{Su}_F),
\end{equation*}
并考虑由 $Qu=[Pu,Su]$ 定义的紧算子 $Q:E\to E\times F$.

\solutionhead{6.10}

1. 注意对一切 $u\in E$,
\begin{align*}
\abs{Q(1)}\norm{u}&\le\norm{Q(1)u-Q(T)u}+\norm{Q(T)u}\\
&=\norm{\widetilde Q(T)(u-Tu)}+\norm{Q(T)u}\le C(\norm{u-Tu}+\norm{Q(T)u}).
\end{align*}

2. 蕴含关系 $N(I-T)=\{0\}\Rightarrow R(I-T)=E$ 的证明: 反设
$R(I-T)=E_1\ne E$.令 $E_n=(I-T)^nE$.则 $(E_n)$ 是闭子空间的递减序列.选取
$u_n\in E_n$ 使得 $\norm{u_n}=1$ 且 $\dist(u_n,E_{n+1})\ge1/2$.写出
\begin{equation*}
Q(T)u_n-Q(T)u_m=Q(T)u_n-Q(1)u_n+Q(1)u_n-Q(1)u_m+Q(1)u_m-Q(T)u_m.
\end{equation*}
因此对 $m>n$, 有 $\norm{Q(T)u_n-Q(T)u_m}\ge\abs{Q(1)}/2$, 这是不可能的.

关于逆命题, 沿用定理  \ref{thm:6.6} 的证明中叙述的论证.

3. 沿用定理  \ref{thm:6.6} 的证明中的相同记号, 写出 $S=T+\widetilde\Lambda\circ P$.
这里 $S\notin\KK(E)$, 但 $\widetilde\Lambda\circ P\in\KK(E)$.因此
$Q(S)\in\KK(E)$ (为什么?).然后像定理  \ref{thm:6.6} 的证明那样继续.

\solutionhead{6.11}

1. 存在整数 $n_0\ge1$ 使得 $\operatorname{Int}F_{n_0}\ne\emptyset$, 从而
$B(u_0,\rho)\subset F_{n_0}$.对每个 $u\in F$ 及满足 $\abs{\lambda}<\rho/\norm{u}$
的 $\lambda$, 有 $u_0+\lambda u\in F_{n_0}$.因此
\begin{equation*}
\abs{\lambda}\abs{u(x)-u(y)}\le\abs{u_0(x)-u_0(y)}+n_0d(x,y)^{1/n_0}
\le2n_0d(x,y)^{1/n_0}.
\end{equation*}
由此推出对一切 $x,y\in K$,
\begin{equation*}
\abs{u(x)-u(y)}\le\frac{2n_0}{\rho}\norm{u}d(x,y)^{1/n_0}.
\end{equation*}

2. Ascoli--Arzelà 定理蕴含 $B_F$ 是紧的.

\solutionhead{6.13}

反设存在某个 $\varepsilon_0>0$ 与序列 $(u_n)$ 使得 $\norm{u_n}_E=1$ 且
$\norm{Tu_n}_F\ge\varepsilon_0+n\abs{u_n}$.则 $\abs{u_n}\to0$, 可以假设
$u_{n_k}\rightharpoonup u$ 弱收敛.但函数 $u\mapsto\abs{u}$ 是凸且连续的.
因此它对弱拓扑下半连续, 从而 $u=0$.由此推出 $Tu_n\to0$.不可能.

\solutionhead{6.15}

1. 若 $u=f+\lambda(T-\lambda I)^{-1}f$, 则 $\lambda u=T(u-f)$, 从而
$\abs{\lambda}\norm{u}\le\norm{T}(\norm{u}+\norm{f})$.

2. 由命题  \ref{prop:6.7} 的证明我们知道: 若 $\mu\in\R$ 使得
$\abs{\mu-\lambda}\norm{(T-\lambda I)^{-1}}<1$, 则 $\mu\in\rho(T)$.因此
$\dist(\lambda,\sigma(T))\ge1/\norm{(T-\lambda I)^{-1}}$.

4. $(U-I)^{-1}=\frac12(T-I)$.

6. 注意关系式 $Uu-\lambda u=f$ 等价于
\begin{equation*}
Tu-\frac{\lambda+1}{\lambda-1}u=\frac{1}{\lambda-1}(f-Tf).
\end{equation*}

\solutionhead{6.16}

2. $(T-\lambda I)^{-1}=\dfrac{1}{1-\lambda^n}\sum_{i=0}^{n-1}\lambda^{n-i-1}T^i$.

3. $(T-\lambda I)^{-1}=-\sum_{i=0}^{n-1}\lambda^{-i-1}T^i$.

4. $(I-T)^{-1}=(I-T^n)^{-1}\sum_{i=0}^{n-1}T^i$.

\solutionhead{6.18}

1. $\norm{S_r}=\norm{S_\ell}=1$.注意 $S_\ell\circ S_r=I$, 从而
$S_r\notin\KK(E),S_\ell\notin\KK(E)$.

3. 对每个 $\lambda\in[-1,+1]$, 算子 $(S_r-\lambda I)$ 不是满射: 例如若
$f=(-1,0,0,\dots)$, 则方程 $S_rx-\lambda x=f$ 没有解 $x\in\ell^2$.

4. $N(S_\ell-\lambda I)=\R(1,\lambda,\lambda^2,\dots)$.

6. $S_r^\star=S_\ell$ 且 $S_\ell^\star=S_r$.

7. 写出 $S_rx-\lambda x=f$, 我们有
\begin{equation*}
\abs{x}=\abs{S_rx}=\abs{\lambda x+f}\le\abs{\lambda}\abs{x}+\abs{f}.
\end{equation*}
从而 $\abs{S_rx-\lambda x}\ge(1-\abs{\lambda})\abs{x}$, 因此 $R(S_r-\lambda I)$ 闭.
应用定理  \ref{thm:2.19} 给出
\begin{equation*}
R(S_r-\lambda I)=N(S_\ell-\lambda I)^\perp
=\left\{x\in\ell^2;\ \sum_{i=1}^\infty\lambda^{i-1}x_i=0\right\}
\end{equation*}
以及 $R(S_\ell-\lambda I)=N(S_r-\lambda I)^\perp=E$.

8. 我们有 $R(S_r\pm I)=N(S_\ell\pm I)^\perp=E$ 以及
$R(S_\ell\pm I)=N(S_r\pm I)^\perp=E$.由第 3 问已知 $R(S_r\pm I)\ne E$.
另一方面, $R(S_\ell\pm I)\ne E$; 否则, 由于 $S_\ell\pm I$ 是单射, 会有
$\pm1\in\rho(S_\ell)$, 矛盾.

9. 若对一切 $n$, $\alpha_n\ne0$, 则 $\operatorname{EV}(S_r\circ M)=\emptyset$;
若对某个 $n$, $\alpha_n=0$, 则 $\operatorname{EV}(S_r\circ M)=\{0\}$.

10. 总可以假设 $\alpha\ne0$; 否则 $S_r\circ M$ 紧, 结论显然.

我们证明对满足 $\abs{\lambda}>\abs{\alpha}$ 的每个 $\lambda$, $(T-\lambda I)$
是双射.注意 $M=\alpha I+K$, 其中 $K$ 是紧算子.令 $T=S_r\circ M$, 得到
$T=\alpha S_r+K_1$ 且 $(T-\lambda I)=(\alpha S_r-\lambda I)+K_1=J\circ(I+K_2)$,
其中 $J=(\alpha S_r-\lambda I)$ 是双射, $K_1,K_2$ 紧.应用定理  \ref{thm:6.6}(c),
只须验证 $N(T-\lambda I)=\{0\}$.这已在第 9 问中建立.

我们证明对 $\abs{\lambda}\le\abs{\alpha}$, $(T-\lambda I)$ 不是双射.反设
$(T-\lambda I)$ 是双射.写出
$(S_r-\frac\lambda\alpha I)=\frac1\alpha(T-\lambda I)-\frac1\alpha K_1=J'\circ(I+K_3)$,
其中 $J'$ 是双射且 $K_3$ 紧.再次应用定理  \ref{thm:6.6}(c), 可见
$(S_r-\frac\lambda\alpha I)$ 单射当且仅当它满射.但由第 2 与第 3 问已知
对 $\abs{\lambda}\le\abs{\alpha}$, $(S_r-\frac\lambda\alpha I)$ 是单射且不是满射.
矛盾.

11. $\sigma(S_r\circ M)=[-\sqrt{\abs{ab}},+\sqrt{\abs{ab}}]$.事实上, 若
$\abs{\lambda}\le\sqrt{\abs{ab}}$, 算子 $(S_r\circ M-\lambda I)$ 不是满射,
因为例如 $f=(-1,0,0,\dots)\notin R(S_r\circ M-\lambda I)$.另一方面, 若
$\abs{\lambda}>\sqrt{\abs{ab}}$, 算子 $(S_r\circ M-\lambda I)$ 是双射, 因为
$(S_r\circ M)^2=abS_r^2$.从而 $\norm{(S_r\circ M)^2}\le\abs{ab}$, 可以应用
习题 6.16 第 4 问.

\solutionhead{6.20}

1. 注意 $\abs{Tu(x)-Tu(y)}\le\abs{x-y}^{1/p'}\norm{u}_p$.若 $1<p<\infty$,
可以应用 Ascoli--Arzelà 定理得出 $T(B_E)$ 在 $C([0,1])$ 中有紧闭包, 从而在
$L^p(0,1)$ 中有紧闭包.若 $p=1$, 应用定理  \ref{thm:4.26}.

2. $\operatorname{EV}(T)=\emptyset$.首先注意 $0\notin\operatorname{EV}(T)$.
事实上, 方程 $Tu=0$ 蕴含对一切 $a,b\in[0,1]$, $\int_0^1u\chi_{[a,b]}=0$.
若 $1<p<\infty$, 利用阶梯函数在 $L^{p'}$ 中的稠密性得出 $u\equiv0$.当 $p=1$
时, 通过用阶梯函数一致逼近 $\phi$, 证明对一切 $\phi\in C([0,1])$ 有
$\int_0^1u\phi=0$.由推论  \ref{cor:4.24} 得出 $u\equiv0$.

3. 对 $\lambda\ne0$ 与 $f\in C([0,1])$, 令 $u=(T-\lambda I)^{-1}f$.则
$v(x)=\int_0^xu(t)dt$ 满足 $v\in C^1([0,1])$ 且 $v-\lambda v'=f$, $v(0)=0$.
因此
\begin{equation*}
u(x)=-\frac1\lambda f(x)-\frac{1}{\lambda^2}\int_0^xe^{(x-t)/\lambda}f(t)dt.
\end{equation*}
同样的公式对 $f\in L^p$ 仍然有效 (用稠密性论证).

4. $(T^\star v)(x)=\int_x^1v(t)dt$.

\solutionhead{6.22}

2. 反设存在某个 $\mu\in Q(\sigma(T))$ 使得 $\mu\notin\sigma(Q(T))$.则
$\mu=Q(\lambda)$, 其中 $\lambda\in\sigma(T)$, 且 $Q(T)-Q(\lambda)I=S$ 是双射.
我们可以写出对一切 $t\in\R$, $Q(t)-Q(\lambda)=(t-\lambda)\widetilde Q(t)$,
从而 $(T-\lambda I)\widetilde Q(T)=\widetilde Q(T)(T-\lambda I)=S$.因此
$T-\lambda I$ 是双射, $\lambda\in\rho(T)$; 矛盾.

3. 取 $E=\R^2$, $T=\begin{pmatrix}0&1\\-1&0\end{pmatrix}$ 且 $Q(t)=t^2$.
则 $\operatorname{EV}(T)=\sigma(T)=\emptyset$ 且
$\operatorname{EV}(T^2)=\sigma(T^2)=\{-1\}$.

4. $T^2+I$ 由 Lax--Milgram 定理是双射.每个没有实根的 2 次多项式可以写成
$Q(t)=t^2+at+b=(t+a/2)^2+b-a^2/4$, 其中 $b-a^2/4>0$, 我们可以再次应用
Lax--Milgram 定理.若多项式 $Q(t)$ 没有实根, 则其根为共轭复数.我们可以写出
$Q(t)=Q_1(t)Q_2(t)\dots Q_\ell(t)$, 其中每个 $Q_i(t)$ 是没有实根的 2 次
多项式.由于 $Q_i(T)$ 是双射, 对 $Q(T)$ 也成立同样结论.

5. (i) 反设 $\mu\in\operatorname{EV}(Q(T))$ 且 $\mu\notin Q(\operatorname{EV}(T))$.
则存在 $u\ne0$ 使得 $Q(T)u=\mu u$.写出
\begin{equation*}
Q(t)-\mu=(t-t_1)(t-t_2)\dots(t-t_q)\widetilde Q(t),
\end{equation*}
其中 $t_i$ 是多项式 $Q(t)-\mu$ 的实根, $\widetilde Q$ 没有实根.由于
$\mu\notin Q(\operatorname{EV}(T))$, 对一切 $i$ 有 $t_i\notin\operatorname{EV}(T)$.
我们有 $(T-t_1I)\dots(T-t_kI)\widetilde Q(T)u=0$.由于这个乘积中的每个因子
都是单射, 我们得出 $u=0$.矛盾.

(ii) 像 (i) 那样论证.

\solutionhead{6.23}

3. 在 $E=\R^2$ 中取 $T(u_1,u_2)=(u_2,0)$.则 $T^2=0$, 从而 $r(T)=0$, 而 $\norm{T}=1$.

5. 在 $E=\R^3$ 中取 $T(u_1,u_2,u_3)=(u_2,-u_1,0)$.则 $\sigma(T)=\{0\}$.
利用 $T^3=-T$ 的事实, 容易看出 $r(T)=1$.

\runin{评注.} 若在 $\CC$ 上的 Banach 空间中工作, 情况完全不同; 见 \S \ref{sec:11.4}.
在那里总有 $r(T)=\max\{\abs{\lambda};\ \lambda\in\sigma(T)\}$.在本例中取
$E=\CC^3$, 我们有 $\sigma(T)=\{0,+i,-i\}$, 从而 $r(T)=1$.

6. 假设该公式对 $T^n$ 成立, 我们有
\begin{align*}
(T^{n+1}u)(t)&=\frac{1}{(n-1)!}\int_0^tds\int_0^s(s-\tau)^{n-1}u(\tau)d\tau\\
&=\frac{1}{(n-1)!}\int_0^tu(\tau)\left[\int_\tau^t(s-\tau)^{n-1}ds\right]d\tau
=\frac{1}{n!}\int_0^t(t-\tau)^nu(\tau)d\tau.
\end{align*}

7. 考虑定义在 $\R$ 上的函数 $f$ 与 $g$:
\begin{equation*}
f(t)=\begin{cases}\frac{1}{(n-1)!}t^{n-1} & \text{若}\ 0\le t\le1,\\ 0 & \text{其他},\end{cases}
\qquad
g(t)=\begin{cases}u(t) & \text{若}\ 0\le t\le1,\\ 0 & \text{其他},\end{cases}
\end{equation*}
于是对 $0\le t\le1$, 有
$(f\star g)(t)=\int_0^1(t-\tau)u(\tau)d\tau=(T^nu)(t)$.我们推出
\begin{equation*}
\norm{f\star g}_{L^p(0,1)}\le\norm{f\star g}_{L^p(\R)}\le\norm{f}_{L^1(\R)}\norm{g}_{L^p(\R)}
=\frac{1}{n!}\norm{u}_{L^p(0,1)}.
\end{equation*}

8. 应用 Stirling 公式.

\solutionhead{6.24}

2. (v) $\Rightarrow$ (vi). 对每个 $\varepsilon>0$, $T_\varepsilon=T+\varepsilon I$
是双射且 $\sigma(T_\varepsilon)\subset[\varepsilon,1+\varepsilon]$.从而
$\sigma(T_\varepsilon^{-1})\subset\left[\frac{1}{1+\varepsilon},\frac1\varepsilon\right]$.
对 $T_\varepsilon^{-1}$ 应用命题  \ref{prop:6.9} 给出对一切 $v\in H$,
$(T_\varepsilon^{-1}v,v)\ge\frac{1}{1+\varepsilon}\abs{v}^2$, 即对一切 $u\in H$,
$(T_\varepsilon u,u)\ge\frac{1}{1+\varepsilon}\abs{T_\varepsilon u}^2$.

3. 令 $U=2T-I$.显然 (vii) 等价于对一切 $u\in H$, $\abs{u}\le\abs{Uu}$ (记作 (vii$'$)).
应用定理  \ref{thm:2.20}, 可见 (vii) $\Rightarrow$ $(-1,+1)\subset\rho(U)=2\rho(T)-1$.
从而 (vii) $\Rightarrow$ (viii).反之, (viii) $\Rightarrow$ $(-1,+1)\subset\rho(U)$.
因此 $\sigma(U)\subset(-\infty,-1]\cup[1,+\infty)$ 且 $\sigma(U^{-1})\subset[-1,+1]$.
由命题  \ref{prop:6.9} 我们知道 $\norm{U^{-1}}\le1$, 即 (vii$'$) 成立.

\solutionhead{6.25}

由构造, 在 $X$ 上有 $M\circ(I+K)=I$, 在 $R(I+K)$ 上有 $(I+K)\circ M=I$.
给定任何 $x\in E$, 写出 $x=x_1+x_2$, 其中 $x_1\in X,x_2\in N(I+K)$.则
$M\circ(I+K)(x)=x_1=x-Px$, 其中 $P$ 是到 $N(I+K)$ 上的投影算子.
对任何 $x\in E$, 有
\begin{equation*}
(I+K)\circ\widetilde M(x)=(I+K)\circ M\circ Q(x)=Qx=x-\widetilde Px,
\end{equation*}
其中 $\widetilde P$ 是到 $R(I+K)$ 在 $E$ 中的补空间上的有限秩投影算子.
\solutionhead{8.8}

4. 我们有
\begin{equation*}
u'_n=\zeta_nu'+\zeta'_nu.
\end{equation*}
显然由控制收敛定理, 在 $L^p$ 中 $\zeta_nu'\to u'$.还须证明在 $L^p$ 中
$\zeta'_nu\to0$.注意
\begin{equation*}
\norm{\zeta'_nu}_p^p\le C\int_{1/n}^{2/n}n^p\abs{u(x)}^pdx,
\end{equation*}
其中 $C=\norm{\zeta'}_{L^\infty}^p$.

当 $p=1$ 时, 由于 $u\in C([0,1])$ 且 $u(0)=0$, 有
\begin{equation*}
n\int_{1/n}^{2/n}\abs{u(x)}dx\le\max_{x\in[1/n,2/n]}\abs{u(x)}\to0\quad(n\to\infty).
\end{equation*}
当 $p>1$ 时, 我们有
\begin{equation*}
n^p\int_{1/n}^{2/n}\abs{u(x)}^pdx=n^p\int_{1/n}^{2/n}x^p\frac{\abs{u(x)}^p}{x^p}dx
\le2^p\int_{1/n}^{2/n}\frac{\abs{u(x)}^p}{x^p}dx\to0
\end{equation*}
(由第 1 问).

\solutionhead{8.9}

1. 由习题 8.8 第 1 问已知 $\dfrac{u'(x)}{x}\in L^p$.另一方面,
\begin{equation*}
u(x)=\int_0^xu'(t)dt=xu'(x)-\int_0^xu''(t)tdt,
\end{equation*}
从而
\begin{equation*}
\frac{u(x)}{x^2}=\frac{u'(x)}{x}-\frac{1}{x^2}\int_0^xu''(t)tdt.
\end{equation*}
但
\begin{equation*}
\frac{1}{x^2}\abs{\int_0^xu''(t)tdt}\le\frac1x\int_0^x\abs{u''(t)}dt\in L^p
\end{equation*}
(同上).

2. 我们有 $v\in C^1((0,1))$ 且
\begin{equation*}
v'(x)=-\frac{u(x)}{x^2}+\frac{u'(x)}{x}\in L^p
\end{equation*}
(由第 1 问).此外, 当 $x\to0$ 时
\begin{equation*}
v(x)=\frac{u(x)}{x}=\frac1x\int_0^xu'(t)dt\to0
\end{equation*}
(因为 $u\in C^1([0,1])$ 且 $u'(0)=0$).

3. 只须证明当 $n\to0$ 时 $\norm{\zeta'_nu'}_p+\norm{\zeta''_nu}_p\to0$.但
\begin{equation*}
\norm{\zeta'_nu'}_p^p\le Cn^p\int_{1/n}^{2/n}\abs{u'(x)}^pdx
\le2^pC\int_{1/n}^{2/n}\frac{\abs{u'(x)}^p}{x^p}dx,
\end{equation*}
以及
\begin{equation*}
\norm{\zeta''_nu}_p^p\le Cn^{2p}\int_{1/n}^{2/n}\abs{u(x)}^pdx
\le4^pC\int_{1/n}^{2/n}\frac{\abs{u(x)}^p}{x^{2p}}dx,
\end{equation*}
由 $\dfrac{u'(x)}{x}\in L^p$ 与 $\dfrac{u(x)}{x^2}\in L^p$ 即得出结论.

4. 设 $u\in X_m$.由归纳假设有 $u'\in X_{m-1}$ 且 $\dfrac{u'(x)}{x^{m-1}}\in L^p(I)$.
接着注意到
\begin{equation*}
\frac{u(x)}{x^m}=\frac{1}{x^m}\int_0^x\frac{u'(t)}{t^{m-1}}t^{m-1}dt.
\end{equation*}
再次应用 Hardy 不等式 (见问题 34 第 C 部分) 得到
\begin{equation*}
\frac{\abs{u(x)}}{x^m}\le\frac1x\int_0^x\frac{\abs{u'(t)}}{t^{m-1}}dt\in L^p(I).
\end{equation*}
为了证明 $\dfrac{u(x)}{x^{m-1}}\in X_1$, 注意
\begin{equation*}
D\left(\frac{u(x)}{x^{m-1}}\right)=\frac{Du(x)}{x^{m-1}}-(m-1)\frac{u(x)}{x^m}\in L^p(I),
\end{equation*}
且当 $x\to0$ 时
\begin{equation*}
\frac{\abs{u(x)}}{x^{m-1}}\le\frac{1}{x^{m-1}}\int_0^x\frac{\abs{u'(t)}}{t^{m-1}}t^{m-1}dt
\le\int_0^x\frac{\abs{u'(t)}}{t^{m-1}}dt\to0
\end{equation*}
(因为 $\dfrac{u'(t)}{t^{m-1}}\in L^p(I)$).

5. 只须对满足 $0\le\ell\le k-1$ 的每个整数 $\ell$ 验证 $D^\ell v\in X_1$.
但 $D^\ell v$ 是形如 $\dfrac{D^{j+\alpha}u(x)}{x^{m-j-k+\ell-\alpha}}$ (其中
$\alpha$ 是满足 $0\le\alpha\le\ell$ 的整数) 的函数的线性组合.然后使用第 4 问.

6. 只须证明当 $\alpha+\beta=m$ 且 $1\le\alpha\le m$ 时, 在 $L^p(I)$ 中
$(D^\alpha\zeta_n)(D^\beta u)\to0$.但 $\abs{D^\alpha\zeta_n(x)}\le Cn^\alpha$,
从而
\begin{align*}
\int_0^1\abs{D^\alpha\zeta_n(x)}^p\abs{D^\beta u(x)}^pdx
&\le Cn^{\alpha p}\int_{1/n}^{2/n}\abs{\frac{D^\beta u(x)}{x^\alpha}}^px^{\alpha p}dx\\
&\le C\int_{1/n}^{2/n}\abs{\frac{D^\beta u(x)}{x^\alpha}}^pdx\to0,
\end{align*}
因为由第 4 问 $\dfrac{D^\beta u(x)}{x^\alpha}\in L^p(I)$.

8. 为证明 $v\in C([0,1])$, 注意 $v(x)=\dfrac1x\int_0^xu'(t)dt$ 且 $u'\in C([0,1])$
满足 $u'(0)=0$.接着证明 $v\in W^{1,1}(I)$.分部积分可见
\begin{equation*}
v'(x)=\frac{1}{x^2}\int_0^xu''(t)tdt,
\end{equation*}
直接计算给出
\begin{equation*}
\norm{v'}_1\le\int_0^1\abs{u''(t)}(1-t)dt\le\norm{u''}_1.
\end{equation*}

9. 令 $u(x)=\int_0^x(1+\abs{\log t})^{-1}dt$.显然 $u\in W^{2,1}(I)$ 满足
$u(0)=u'(0)=0$, 此外 $\dfrac{u'(x)}{x}\notin L^1(I)$.关系式
$\dfrac{u(x)}{x^2}=\dfrac{u'(x)}{x}-v'(x)$ 结合第 8 问表明
$\dfrac{u(x)}{x^2}\notin L^1(I)$.

\solutionhead{8.10}

4. 显然当 $n\to\infty$ 时, $v'_n(x)=G'(nu(x))u'(x)\to f(x)$ 几乎处处, 其中
\begin{equation*}
f(x)=\begin{cases}
0 & \text{若}\ u(x)\ne0,\\
u'(x) & \text{若}\ u(x)=0.
\end{cases}
\end{equation*}

6. 对一切 $\phi\in C_c^1(I)$, 有 $\int_0^1v_n\phi'=-\int_0^1v'_n\phi$.令
$n\to\infty$ 取极限得出对一切 $\phi\in C_c^1(I)$, $\int_0^1f\phi=0$, 因此在
$I$ 上 $f=0$ 几乎处处, 即在 $[u=0]$ 上 $u'(x)=0$ 几乎处处.

\solutionhead{8.12}

1. 使用习题 8.2 以及 $\liminf_{n\to\infty}\norm{u'_n}_{L^p}\ge\norm{u'}_{L^p}$
的事实.

2. 考虑习题 8.2 中的序列 $(u_n)$.我们有 $\norm{u_n}_{L^1}\le1/2$ 且
$\norm{u'_n}_{L^1}=1$.因此 $\frac23u_n\in B_1$.另一方面, 在 $L^1$ 中
$\frac23u_n\to\frac23u$, 其中
\begin{equation*}
u(x)=\begin{cases}
0 & \text{若}\ x\in(0,1/2),\\
1 & \text{若}\ x\in(1/2,1).
\end{cases}
\end{equation*}
但 $u\notin W^{1,1}$.因此 $B_1$ 在 $L^1$ 中不是闭的.

\solutionhead{8.16}

2. $R(A)=L^p(0,1)$ 且 $N(A)=\{0\}$.

3. $v\in D(A^\star)$ 当且仅当 $v\in L^{p'}$ 且存在常数 $C$ 使得对一切
$u\in D(A)$, $\abs{\int_0^1vu'}\le C\norm{u}_p$.特别地,
$v\in D(A^\star)\Rightarrow v\in W^{1,p'}$, 此时
\begin{equation}
\abs{u(1)v(1)-\int_0^1uv'}\le C\norm{u}_p\quad\forall u\in D(A).
\tag{S1}
\label{eq:S8.16}
\end{equation}
由 \eqref{eq:S8.16} 推出对一切 $u\in D(A)$,
\begin{equation*}
\abs{u(1)}\abs{v(1)}\le(C+\norm{v'}_{p'})\norm{u}_p.
\end{equation*}
由于存在 $D(A)$ 中的序列 $(u_n)$ 使得 $u_n(1)=1$ 且 $\norm{u_n}_p\to0$,
由此推出 $v(1)=0$.于是我们证明了 $v\in D(A^\star)\Rightarrow v\in W^{1,p'}$
且 $v(1)=0$.容易推出
\begin{equation*}
D(A^\star)=\{v\in W^{1,p'};\ v(1)=0\},
\end{equation*}
且 $A^\star v=-v'$.

4. 我们有 $N(\widetilde A)=\{0\}$,
$R(\widetilde A)=\left\{f\in L^p;\ \int_0^1f(t)dt=0\right\}$, 以及
$(\widetilde A)^\star v=-v'$, 定义域为 $D((\widetilde A)^\star)=W^{1,p'}$.

\solutionhead{8.17}

在确定 $D(A^\star)$ 时, 记住下述事实是有用的: 设 $I=(0,1)$, $1<p\le\infty$.
假设 $u\in L^p(I)$ 满足
\begin{equation}
\abs{\int_Iu\phi'}\le C\norm{\phi}_{p'}\quad\forall\phi\in C_c^1(I)\ \text{使得}\ \int_I\phi=0,
\tag{S1}
\label{eq:S8.17}
\end{equation}
则 $u\in W^{1,p}(I)$.

事实上, 固定满足 $\int_I\psi_0=1$ 的函数 $\psi_0\in C_c^1(I)$.设 $\zeta$
为 $C_c^1(I)$ 中的任何函数.把 $\phi=\zeta-(\int_I\zeta)\psi_0$ 代入
\eqref{eq:S8.17}, 我们得到
\begin{equation*}
\abs{\int_Iu\zeta'}\le C\norm{\zeta}_{p'}+C'\abs{\int_I\zeta},
\end{equation*}
其中 $C'$ 只依赖于 $u$ 与 $\psi_0$.因此由命题  \ref{prop:8.3}, $u\in W^{1,p}(I)$.

当 $Au=u''-xu'$ 时, 我们有 $A^\star v=v''+xv'+v$.注意下述恒等式
\begin{equation*}
A^\star(e^{-x^2/2}u)=e^{-x^2/2}Au\quad\forall u\in H^2(I),
\end{equation*}
这允许在各种边界条件下计算 $N(A^\star)$.

\solutionhead{8.19}

给定 $f\in L^2(0,1)$, 令 $F(x)=\int_0^xf(t)dt$.则
\begin{equation*}
\varphi^\star(f)=\begin{cases}
\dfrac12\int_0^1F^2(x)dx & \text{若}\ \int_0^1f(t)dt=0,\\
+\infty & \text{其他}.
\end{cases}
\end{equation*}
事实上, 若 $\int_0^1f(t)dt=0$, 则对一切 $v\in H^1(0,1)$,
$\int_0^1fv=\int_0^1F'v=-\int_0^1Fv'$, 且
\begin{align*}
\varphi^\star(f)&=\sup_{v\in H^1}\left\{\int_0^1fv-\frac12\int_0^1v'^2\right\}
=\sup_{v\in H^1}\left\{-\int_0^1Fv'-\frac12\int_0^1v'^2\right\}\\
&=\sup_{w\in L^2}\left\{-\int_0^1Fw-\frac12\int_0^1w^2\right\}
=\frac12\int_0^1F^2.
\end{align*}

\solutionhead{8.21}

2. 设 $U$ 为满足 $-(pU')'+qU=f$ 在 $(0,1)$ 上且 $U(1)=0$ 的任何函数.则
\begin{equation*}
\int_0^1fv_0=p(0)(U'(0)-k_0U(0)).
\end{equation*}
因此若 $\int_0^1fv_0=0$, 任何这样的函数 $U$ 满足 $U'(0)=k_0U(0)$.由于
$U(0)$ 可以任意选取, 我们看到解集是一维的.

\solutionhead{8.22}

1. 函数 $\rho(x)=x$ 属于 $H^1(0,1)$, 但 $\sqrt{\rho(x)}=\sqrt{x}\notin H^1(0,1)$.

2. 对在 $(0,1)$ 上满足 $\rho\ge0$ 的每个 $\rho\in H^1(0,1)$, 令
$\gamma_\varepsilon=\sqrt{\rho+\varepsilon}$.由于函数 $t\mapsto\sqrt{t+\varepsilon}$
在 $[0,+\infty)$ 上属于 $C^1$ 类, 我们推出 $\gamma_\varepsilon\in H^1(0,1)$,
此外
\begin{equation*}
\gamma'_\varepsilon=\frac{\rho'}{2\sqrt{\rho+\varepsilon}},
\end{equation*}
从而在集合 $[\rho>0]$ 上 $\abs{\gamma'_\varepsilon}\le\mu$.另一方面, 我们
知道在集合 $[\rho=0]$ 上 $\rho'=0$ 几乎处处 (见习题 8.10), 从而在 $[\rho=0]$
上 $\abs{\gamma'_\varepsilon}\le\mu$ 几乎处处.因此在 $(0,1)$ 上几乎处处
$\abs{\gamma'_\varepsilon}\le\mu$.

由此推出若 $\mu\in L^2$, 则当 $\varepsilon\to0$ 时
$\norm{\gamma'_\varepsilon}_{L^2}\le C$.由于当 $\varepsilon\to0$ 时在
$C([0,1])$ 中 $\gamma_\varepsilon\to\sqrt{\rho}$, 且在 $L^2(0,1)$ 中
$\gamma'_\varepsilon\to\mu$, 我们得出 (见习题 8.2) $\sqrt{\rho}\in H^1(0,1)$
且 $(\sqrt{\rho})'=\mu$.

反之, 若 $\sqrt{\rho}\in H^1(0,1)$, 令 $\gamma=\sqrt{\rho}$, 从而 $\rho=\gamma^2$
且 $\rho'=2\gamma\gamma'$.因此在 $[\rho>0]$ 上 $\mu=\gamma'$ 几乎处处,
此外在 $[\rho=0]$ 上 $\mu=\gamma'$ 几乎处处 (因为在 $[\gamma=0]=[\rho=0]$ 上
$\gamma'=0$ 几乎处处).

\solutionhead{8.24}

1. 可以选取 $C_\varepsilon=1+1/\varepsilon$.

2. 弱表述为: $u\in H^1(I)$,
\begin{equation*}
a(u,v)=\int_0^1(u'v'+kuv)-u(1)v(1)=\int_0^1fv\quad\forall v\in H^1(I).
\end{equation*}
显然 $a(u,v)$ 是 $H^1(0,1)$ 上的连续双线性形式.由第 1 问, 例如当 $k>2$ 时
它是强制的.相应的极小化问题为
\begin{equation*}
\min_{v\in H^1}\left\{\frac12\int_0^1(v'^2+kv^2)-\frac12v(1)^2-\int_0^1fv\right\}.
\end{equation*}

3. 设 $g\in L^2(I)$, 设 $v\in H^2(I)$ 为 (1) 的相应解 (其中 $f$ 换成 $g$).
我们有
\begin{align*}
(Tf,g)_{L^2}&=\int_0^1ug=\int_0^1u(-v''+kv)\\
&=-u(1)v'(1)+u(0)v'(0)+u'(1)v(1)-u'(0)v(0)+\int_0^1(-u''+ku)v\\
&=-u(1)v(1)+u(1)v(1)+\int_0^1fv=(f,Tg)_{L^2}.
\end{align*}
因此 $T$ 自伴.它紧是因为它是由 $L^2(I)$ 到 $H^1(I)$ 中的有界算子, 且
$H^1(I)\subset L^2(I)$ 紧嵌入 (见定理  \ref{thm:8.8}).

4. 由 \S \ref{sec:6.4} 的结果我们知道存在 $L^2(I)$ 中的序列 $(u_n)$ 满足
$Tu_n=\mu_nu_n$, 其中 $\norm{u_n}_{L^2}=1$, 对一切 $n$ 有 $\mu_n>0$, 且
$\mu_n\to0$.因此我们有 $-u''_n+ku_n=\dfrac{1}{\mu_n}u_n$, 从而在 $I$ 上
$-u''_n=\left(\dfrac{1}{\mu_n}-k\right)u_n$.

5. 排除了 $\lambda=0$ 的值 (为什么?).若 $\lambda>0$, 有
$u(x)=A\cos\sqrt{\lambda}x+B\sin\sqrt{\lambda}x$, 其中常数 $A$ 与 $B$ 经
调整以满足边界条件, 即 $B=0$ 且 $A(\cos\sqrt{\lambda}+\sqrt{\lambda}\sin\sqrt{\lambda})=0$,
从而 $A\ne0$ 当且仅当 $\sqrt{\lambda}$ 是方程 $\tan t=-1/t$ 的解 (通过考察
图形可见它有无限正解序列 $t_n\to\infty$).若 $\lambda<0$, 有
$u(x)=Ae^{\sqrt{\abs{\lambda}}x}+Be^{-\sqrt{\abs{\lambda}}x}$.把它与边界
条件结合给出 $A=B$ 且
$A\sqrt{\abs{\lambda}}e^{\sqrt{\abs{\lambda}}}-B\sqrt{\abs{\lambda}}e^{-\sqrt{\abs{\lambda}}}=Ae^{\sqrt{\abs{\lambda}}}+Be^{-\sqrt{\abs{\lambda}}}$.
为了有某个 $u\not\equiv0$, $\lambda$ 必须满足
$\sqrt{\abs{\lambda}}(e^{\sqrt{\abs{\lambda}}}-e^{-\sqrt{\abs{\lambda}}})=e^{\sqrt{\abs{\lambda}}}+e^{-\sqrt{\abs{\lambda}}}$,
即 $t=\sqrt{\abs{\lambda}}$ 是方程 $e^{2t}=\dfrac{t+1}{t-1}$ 的解.考察图形表明
存在唯一解 $t_0>1$, 此时 $\lambda=-t_0^2$.

\solutionhead{8.25}

2. 反设存在 $H^1(I)$ 中的序列 $(u_n)$ 使得 $a(u_n,u_n)\to0$ 且
$\norm{u_n}_{H^1(I)}=1$.抽取子序列 $(u_{n_k})$ 可以假设在 $L^2$ 中
$u'_{n_k}\rightharpoonup u'$ 弱收敛且 $u_n\to u$ 强收敛.由下半连续性 (见命题
 \ref{prop:3.5}) 有 $\liminf\int_I(u'_{n_k})^2\ge\int_I(u')^2$, 从而
$a(u,u)=0$, 因此 $u=0$.但 $\int_I(u'_{n_k})^2=1-\int_I u_{n_k}^2$, 从而
\begin{equation*}
a(u_{n_k},u_{n_k})=\int_I(u'_{n_k})^2+\left(\int_0^1u_{n_k}\right)^2
=1-\int_I u_{n_k}^2+\left(\int_0^1u_{n_k}\right)^2\to1.
\end{equation*}
矛盾.

4. 对一切 $v\in H^1(I)$, 有 $\int_Iu'v'=\int_Igv$, 其中
$g=f-(\int_0^1u)\chi_{(0,1)}$.因此 $u\in H^2(I)$ 且满足
\begin{equation*}
\begin{cases}
-u''+(\int_0^1u)\chi_{(0,1)}=f & \text{在}\ I\ \text{上},\\
u'(0)=u'(2)=0.
\end{cases}
\end{equation*}

5. $u\in C^2(I)$ 当且仅当 $\int_0^1u=0$.这发生当且仅当 $\int_If=0$.

8. $T$ 的特征值为正, 且若 $1/\lambda$ 是特征值, 必有非平凡函数 $u\not\equiv0$
满足
\begin{equation*}
\begin{cases}
-u''+\int_0^1u=\lambda u & \text{在}\ (0,1)\ \text{上},\\
-u''=\lambda u & \text{在}\ (1,2)\ \text{上},\\
u'(0)=u'(2)=0,\\
u(1-)=u(1+),\quad u'(1-)=u'(1+).
\end{cases}
\end{equation*}
因此在 $(0,1)$ 上 $u(x)=\dfrac k\lambda+A\cos(\sqrt{\lambda}x)$, 在 $(1,2)$
上 $u(x)=A'\cos(\sqrt{\lambda}(x-2))$, 其中常数 $k,A,A'$ 由关系式
$u(1-)=u(1+)$, $u'(1-)=u'(1+)$ 以及 $k=\int_0^1u$ 确定.我们得出或者
$\sin\sqrt{\lambda}=0$, 即 $\lambda=n^2\pi^2$ ($n=1,2,\dots$); 或者 $\lambda$
是方程 $\tan\sqrt{\lambda}=2\sqrt{\lambda}(1-\lambda)$ 的解.

\solutionhead{8.26}

3. 令 $a(v,v)=\int_I pv'^2+qv^2$.我们有 $(S_Nf-S_Df,f)=\int_If(u_N-u_D)$.
由极小化性质已知 $\frac12a(u_N,u_N)-\int_Ifu_N\le\frac12a(u_D,u_D)-\int_Ifu_D$.
另一方面, $a(u_N,u_N)=\int_Ifu_N$ 且 $a(u_D,u_D)=\int_Ifu_D$.因此
$\int_If(u_N-u_D)\ge0$.

6. 令 $a_i(v,v)=a(v,v)+k_iv^2(0)$ ($i=1,2$), 以及 $u_{k_1}=u_1,u_{k_2}=u_2$.
由于 $u_i$ 是 $\frac12a_i(v,v)-\int_Ifv$ 在 $V=\{v\in H^1(I);\ v(1)=0\}$
上的极小元, 我们有
\begin{equation*}
\frac12a(u_2,u_2)+\frac12k_2u_2^2(0)-\int_Ifu_2
\le\frac12a(u_1,u_1)+\frac12k_2u_1^2(0)-\int_Ifu_1.
\end{equation*}
另一方面, $a(u_1,u_1)+k_1u_1^2(0)=\int_Ifu_1$, 以及
$a(u_2,u_2)+k_2u_2^2(0)=\int_Ifu_2$.因此
\begin{equation*}
-\frac12\int_Ifu_2\le\frac12\int_Ifu_1+\frac12(k_2-k_1)u_1^2(0)-\int_Ifu_1,
\end{equation*}
从而
\begin{equation*}
(S_{k_2}f-S_{k_1}f,f)=\int_If(u_2-u_1)\ge(k_1-k_2)u_1^2(0)\ge0.
\end{equation*}

\solutionhead{8.27}

4. 问题 $-\phi''+\phi=1$ 在 $I$ 上, $\phi(-1)=\phi(1)=0$ 的解由
$\phi(x)=1+A(e^x+e^{-x})$ 给出, 其中 $A=-e/(e^2+1)$.由 $\phi$ 的唯一性必有
$u=\lambda u(0)\phi$.因此
\begin{equation*}
\lambda_0=\frac{1}{\phi(0)}=\frac{e^2+1}{(e-1)^2}.
\end{equation*}

5. 方程 (1) 变为 $u=S(f+\lambda u(0))=Sf+\lambda u(0)S1=Sf+\lambda u(0)\phi$.
从而 $u(0)(1-\lambda\phi(0))=(Sf)(0)$, 即
\begin{equation*}
u(0)=\frac{\lambda_0(Sf)(0)}{\lambda_0-\lambda},
\end{equation*}
且 $u=Sf+\dfrac{\lambda\lambda_0(Sf)(0)\phi}{\lambda_0-\lambda}$ 是所求的解.

6. 当 $\lambda=\lambda_0$ 时, 解 $u$ 的存在性蕴含 $(Sf)(0)=0$.反之, 设
$(Sf)(0)=0$.则 (1) 的解必须具有形式 $u=Sf+A\phi$ (对某个常数 $A$).直接计算
表明任何这样的 $u$ 满足 $-u''+u=f+A$.但 $u(0)=(Sf)(0)+\frac{A}{\lambda_0}=\frac{A}{\lambda_0}$.
因此我们有 $-u''+u=f+\lambda_0u(0)$, 即对任何 $A$, (1) 都成立.因此当
$\lambda=\lambda_0$ 时 (1) 的所有解的集合为 $Sf+\R\phi$.

\solutionhead{8.29}

2. 解 $u\in H^1(0,1)$ 的存在性与唯一性来自 Lax--Milgram 定理.特别地, $u$
满足对一切 $v\in H_0^1(0,1)$, $\int_0^1u'v'=\int_0^1(f-u)v$, 因此
$u'\in H^1(0,1)$, 即 $u\in H^2(0,1)$; 此外在 $(0,1)$ 上 $-u''+u=f$.利用
$u\in H^2(0,1)$ 的信息, 我们现在可以写出
\begin{align*}
a(u,v)&=\int_0^1(-u''+u)v+u'(1)v(1)-u'(0)v(0)+(u(1)-u(0))(v(1)-v(0))\\
&=\int_0^1fv\quad\forall v\in H^1(0,1).
\end{align*}
从而对一切 $v\in H^1(0,1)$,
\begin{equation*}
(u'(1)+u(1)-u(0))v(1)-(u'(0)+u(1)-u(0))v(0)=0.
\end{equation*}
由于 $v(0)$ 与 $v(1)$ 是任意的, 我们得出
\begin{equation*}
u'(1)+u(1)-u(0)=0\quad\text{且}\quad u'(0)+u(1)-u(0)=0.
\end{equation*}

5. 使用定理  \ref{thm:8.19} 的证明中的相同函数 $G$, 取 $v=G(-u)$, 由于
$f\ge0$ 与 $G\ge0$, 有 $a(u,G(-u))=\int_0^1fG(-u)\ge0$.另一方面,
\begin{align*}
a(u,G(-u))&=-\int_0^1u'^2G'(-u)-\int_0^1(-u)G(-u)+(u(1)-u(0))(G(-u(1))-G(-u(0)))\\
&\le-\int_0^1(-u)G(-u)
\end{align*}
(因为 $G$ 单调递增).由此推出 $\int_0^1(-u)G(-u)\le0$, 从而 $-u\le0$.

7. 设 $1/\lambda$ 为特征值, $u$ 为相应的特征函数.则 $-u''+u=\lambda u$ 在
$(0,1)$ 上, $u'(0)=u(0)-u(1)$, $u'(1)=u(0)-u(1)$.由于
$a(u,u)=\lambda\int_0^1u^2\ge\int_0^1u^2$, 可见 $\lambda\ge1$.此外 $\lambda=1$
是对应于 $u=\text{常数}$ 的特征值.现设 $\lambda>1$ 并令 $\alpha=\sqrt{\lambda-1}$.
必有 $u(x)=A\cos\alpha x+B\sin\alpha x$.为了满足边界条件, 需要满足方程组
\begin{align*}
B\alpha&=A-A\cos\alpha-B\sin\alpha,\\
-A\alpha\sin\alpha+B\alpha\cos\alpha&=A-A\cos\alpha-B\sin\alpha.
\end{align*}
这个方程组有非平凡解当且仅当 $2(1-\cos\alpha)+\alpha\sin\alpha=0$, 即
$\sin(\alpha/2)=0$ 或 $(\alpha/2)+\tan(\alpha/2)=0$.

\solutionhead{8.34}

1. 设 $u$ 为古典解.则对一切 $v\in H^1(0,1)$,
\begin{equation*}
-u'(1)v(1)+u'(0)v(0)+\int_0^1(u'v'+uv)=\int_0^1fv.
\end{equation*}
令 $V=\{v\in H^1(0,1);\ v(0)=v(1)\}$.若 $v\in V$, 我们得到
\begin{equation*}
a(u,v)=\int_0^1(u'v'+uv)=\int_0^1fv+kv(0).
\end{equation*}
弱表述为: $u\in V$ 且对一切 $v\in V$, $a(u,v)=\int_0^1fv+kv(0)$.

2. 由 Lax--Milgram 定理存在唯一的弱解 $u\in V$, 相应的极小化问题为
\begin{equation*}
\min_{v\in V}\left\{\frac12\int_0^1(v'^2+v^2)-\int_0^1fv-kv(0)\right\}.
\end{equation*}

3. 显然任何弱解 $u$ 属于 $H^2(0,1)$ 且满足在 $(0,1)$ 上几乎处处 $-u''+u=f$,
对一切 $v\in V$ 有 $u'(1)v(1)-u'(0)v(0)=kv(0)$, 即 $u'(1)-u'(0)=k$.

5. $T$ 的特征值由 $\lambda_k=1/\mu_k$ 给出, 其中 $\mu_k$ 对应于在 $(0,1)$
上几乎处处
\begin{equation*}
-u''+u=\mu_ku,\qquad u(1)=u(0),\ u'(1)=u'(0)
\end{equation*}
的非平凡解.因此 $\mu_k\ge1$ 且 $u$ 由
\begin{equation*}
u(x)=A\sin(\sqrt{\mu_k-1}x)+B\cos(\sqrt{\mu_k-1}x)
\end{equation*}
给出, 其中 $\sqrt{\mu_k-1}=2\pi k$, $k=0,1,\dots$.

\solutionhead{8.38}

2. 设 $Tu=\lambda u$, 其中 $u\in H^2(\R)$ 且 $u\not\equiv0$.显然 $\lambda\ne0$
且 $u$ 在 $\R$ 上满足 $-u''+u=\dfrac1\lambda u$.若 $\lambda=1$, 有
$u(x)=Ax+B$.由于 $u\in L^2(\R)$, 推出 $A=B=0$.因此 $1\notin\operatorname{EV}(T)$.
若 $(1/\lambda-1)>0$, 有 $u(x)=A\sin\alpha x+B\cos\alpha x$, 其中
$\alpha=\sqrt{1/\lambda-1}$.条件 $u\in L^2(\R)$ 再次给出 $A=B=0$.类似地若
$(1/\lambda-1)<0$, 除 $u\equiv0$ 外没有解.因此 $\operatorname{EV}(T)=\emptyset$.
$T$ 不能是紧算子.否则由定理  \ref{thm:6.8} 我们会有 $\sigma(T)=\{0\}$, 进而由
推论  \ref{cor:6.10} $T\equiv0$.但显然 $T\not\equiv0$.

3. 若 $\lambda<0$, 例如由 Lax--Milgram 定理以及对一切 $f\in H$ 有 $(Tf,f)\ge0$
的事实, $(T-\lambda I)$ 是由 $H=L^2(\R)$ 到自身的双射.因此 $\lambda\in\rho(T)$.

4. 若 $\lambda>1\ge\norm{T}$, 由命题  \ref{prop:6.7} 我们有 $\lambda\in\rho(T)$.

6. $T$ 不是满射, 因为 $R(T)\subset H^2(\R)$.

7. $T-I$ 不是满射.事实上, 若试图对给定的 $\phi\in L^2(\R)$ 求解 $Tf-f=\phi$,
我们被引向 $-u''+u=f$ (令 $u=Tf$) 与 $u=f+\phi$.因此 $u''=\phi$ 允许解
$u\in H^2$.例如设 $\operatorname{supp}\phi\subset[0,1]$.直接计算给出对一切
$x\le0$ 有 $u(x)=0$, 对一切 $x\ge1$ 有 $u(x)=0$.从而 $u'(0)=u'(1)=0$.由此
推出 $0=u'(1)-u'(0)=\int_0^1\phi$.因此当 $\int_0^1\phi\ne0$ 时方程 $Tf-f=\phi$
没有解 $f\in L^2(\R)$.从而 $T-I$ 不能是满射.

8. $T-\lambda I$ 不是满射.事实上, 若试图求解 $Tf-\lambda f=\phi$, 我们被
引向 $-u''+u=f$ (令 $u=Tf$) 与 $u=\lambda f+\phi$.因此 $-u''+u=\dfrac1\lambda(u-\phi)$.
再次假设 $\operatorname{supp}\phi\subset[0,1]$.在 $[0,1]$ 外会有 $u''=-\mu^2u$,
其中 $\mu=\sqrt{1/\lambda-1}$.因此在 $[0,1]$ 外 $u\equiv0$, 从而
$u(0)=u'(0)=u(1)=u'(1)=0$.方程 $-u''+(1-1/\lambda)u=-\dfrac1\lambda\phi$ 蕴含
对在 $(0,1)$ 上的 $-v''=\mu^2v$ 的任何解 $v$, $\int_0^1\phi v=0$; 例如
$\int_0^1\phi(x)\sin\mu x=0$.因此当 $\int_0^1\phi(x)\sin\mu x\ne0$ 时方程
$Tf-\lambda f=\phi$ 没有解 $f\in L^2(\R)$.从而 $(T-\lambda I)$ 不是满射.

\solutionhead{8.39}

2. 对一切 $v\in\R$, 有 $v^2\le\frac12v^4+\frac12$, 从而
\begin{equation*}
\varphi(v)\ge\frac12\norm{v}_{H^1}^2-\frac14-\norm{f}_{L^2}\norm{v}_{H^1}.
\end{equation*}
因此当 $\norm{v}_{H^1}\to\infty$ 时 $\varphi(v)\to\infty$.

3. 唯一性来自 $\varphi$ 在 $H^1(0,1)$ 上严格凸的事实; 这是函数 $t\mapsto t^4$
在 $\R$ 上严格凸的推论.

4. 我们有
\begin{align*}
\varphi(u+\varepsilon v)&=\frac12\int_0^1(u'^2+2\varepsilon u'v'+\varepsilon^2v'^2)\\
&\quad+\frac14\int_0^1(u^4+4\varepsilon u^3v+6\varepsilon^2u^2v^2+4\varepsilon^3uv^3+\varepsilon^4v^4)-\int_0^1f(u+\varepsilon v).
\end{align*}
写出 $\varphi(u)\le\varphi(u+\varepsilon v)$ 给出
$\int_0^1(u'v'+u^3v-fv)+A_\varepsilon\ge0$, 其中当 $\varepsilon\to0$ 时
$A_\varepsilon\to0$.令 $\varepsilon\to0$ 取极限并选取 $\pm v$ 给出对一切
$v\in H^1(0,1)$,
\begin{equation*}
\int_0^1(u'v'+u^3v-fv)=0.
\end{equation*}

6. 由函数 $t\mapsto t^4$ 的凸性, 对一切 $u,v\in\R$ 有
$\frac14v^4-\frac14u^4\ge u^3(v-u)$.另一方面, 显然对一切 $u,v\in H^1(0,1)$,
在 $(0,1)$ 上几乎处处有 $\frac12v'^2-\frac12u'^2\ge u'(v'-u')$.从而对一切
$u,v\in H^1(0,1)$,
\begin{equation*}
\varphi(v)-\varphi(u)\ge\int_0^1u'(v'-u')+\int_0^1u^3(v-u)-\int_0^1f(v-u).
\end{equation*}
若 $u$ 是 (3) 的解, 则对一切 $v\in H^1(0,1)$ 有
$\int_0^1u'(v'-u')+\int_0^1u^3(v-u)=\int_0^1f(v-u)$, 从而对一切
$v\in H^1(0,1)$, $\varphi(u)\le\varphi(v)$.

9. 我们断言当 $\norm{v}_{H^1}\to\infty$ 时 $\psi(v)\to+\infty$.事实上, 这
归结为证明对每个常数 $C$, 集合 $\{v\in H^1(0,1);\ \psi(v)\le C\}$ 在
$H^1(0,1)$ 中有界.若 $\psi(v)\le C$, 写出
\begin{equation*}
\int_0^1fv=\int_0^1f(v-v(0))+v(0)\le\norm{f}_{L^2}(\norm{v'}_{L^2}+\abs{v(0)}),
\end{equation*}
从而 $\norm{v'}_{L^2}$ 与 $\abs{v(0)}$ 有界 (为什么?).因此
$\norm{v}_{L^2}\le\norm{v'}_{L^2}+\abs{v(0)}$ 也有界, 从而 $\norm{v}_{H^1}$
有界.极小元的唯一性: 验证
$\psi\left(\dfrac{u_1+u_2}{2}\right)\le\frac12(\psi(u_1)+\psi(u_2))$, 且
等号成立当且仅当 $u'_1=u'_2$ 且 $u_1(0)=u_2(0)$, 即 $u_1=u_2$.

我们有
\begin{align*}
\psi(u+\varepsilon v)&=\frac12\int_0^1(u'^2+2\varepsilon u'v'+\varepsilon^2v'^2)\\
&\quad+\frac14(u^4(0)+4\varepsilon u^3(0)v(0)+\dots+\varepsilon^4v^4(0))-\int_0^1f(u+\varepsilon v).
\end{align*}
若 $u$ 是 $\psi$ 的极小元, 写出 $\psi(u)\le\psi(u+\varepsilon v)$, 得到
$\int_0^1(u'v'-fv)+u^3(0)v(0)+B_\varepsilon\ge0$, 其中当 $\varepsilon\to0$ 时
$B_\varepsilon\to0$.令 $\varepsilon\to0$ 取极限并选取 $\pm v$ 给出
\begin{equation}
\int_0^1(u'v'-fv)+u^3(0)v(0)=0\quad\forall v\in H^1(0,1).
\tag{S1}
\label{eq:S8.39a}
\end{equation}
由此推出 $u\in H^2(0,1)$ 满足在 $(0,1)$ 上几乎处处
\begin{equation}
-u''=f.
\tag{S2}
\label{eq:S8.39b}
\end{equation}
回到 \eqref{eq:S8.39a} 并利用 \eqref{eq:S8.39b} 给出对一切 $v\in H^1(0,1)$,
$u'(1)v(1)-u'(0)v(0)+u^3(0)v(0)=0$, 从而
\begin{equation}
u'(1)=0,\qquad u'(0)=u^3(0).
\tag{S3}
\label{eq:S8.39c}
\end{equation}
反之, 满足 \eqref{eq:S8.39b} 与 \eqref{eq:S8.39c} 的任何函数 $u$ 是 $\psi$
的极小元: 论证与第 6 问相同.在这种情形下我们有显式解.\eqref{eq:S8.39b} 的
通解由 $u(x)=-\int_0^x(x-t)f(t)dt+Ax+B$ 给出, 此时 \eqref{eq:S8.39c} 等价于
$A=\int_0^1f(t)dt$ 且 $A=B^3$.

\solutionhead{8.42}

2. 对方程
\begin{equation}
v(x)=p^{1/4}(t)u(t)
\tag{S1}
\label{eq:S8.42a}
\end{equation}
关于 $t$ 求导给出
\begin{equation*}
v'(x)p^{-1/2}(t)=\frac14p^{-3/4}(t)p'(t)u(t)+p^{1/4}(t)u'(t).
\end{equation*}
从而
\begin{equation}
p(t)u'(t)=v'(x)p^{1/4}(t)-\frac14p'(t)u(t)=v'(x)p^{1/4}(t)-\frac14p'(t)p^{-1/4}(t)v(x).
\tag{S2}
\label{eq:S8.42b}
\end{equation}
对 \eqref{eq:S8.42b} 关于 $t$ 求导给出
\begin{equation}
(pu')'=v''(x)p^{-1/4}(t)-\frac14p''(t)p^{-1/4}(t)v(x)+\frac{1}{16}p'(t)^2p^{-5/4}(t)v(x).
\tag{S3}
\label{eq:S8.42c}
\end{equation}
把 \eqref{eq:S8.42c} 与 $(0,1)$ 上的方程 $-(pu')'+qu=\mu u$ 结合得出
\begin{equation}
v''(x)p^{-1/4}(t)-\frac14p''(t)p^{-1/4}(t)v(x)+\frac{1}{16}p'(t)^2p^{-5/4}(t)v(x)=(q(t)-\mu)p^{-1/4}(t)v(x).
\tag{S4}
\label{eq:S8.42d}
\end{equation}
因此 $v$ 满足在 $(0,L)$ 上 $-v''+a(x)v=\mu v$, 其中
\begin{equation*}
a(x)=q(t)+\frac14p''(t)-\frac{1}{16}p'(t)^2p^{-1}(t).
\end{equation*}
